\documentclass[output=paper,colorlinks,citecolor=brown]{langscibook}
\ChapterDOI{10.5281/zenodo.22078589}

\author{Sergei Tatevosov\orcid{0000-0002-1602-0116}\affiliation{Lomonosov Moscow State University}}

\SetupAffiliations{mark style=none}

\title[Event structure and derivational morphology]{Event structure and derivational morphology: A few reflections on the system of Russian}

\abstract{This chapter explores the event structure of simplex imperfective, lexically prefixed perfective, and secondary imperfective verbs and their extended projections. It pursues two related goals. One is to make a case, building on Wolfgang Klein’s work, for a higher subevental complexity of prefixed configurations, be they perfective or secondary imperfective, compared to simplex configurations. The other goal is to investigate the typology of lexical prefixes which are responsible for a result state being part of the event structure. The basic divide is between the verbs where the denotation of a stem undergoes a non-compositional shift after prefixation, and the verbs where this does not happen. The latter, compositional lexically prefixed verbs, are further classified according to the argument structure of a prefix: Monadic prefixes take a single individual argument, the holder of a result state, while dyadic prefixes denote a relation between two entities, the figure and the ground. Among the verbs with monadic prefixes, two important classes are isolated. One contains verbs with argument-sharing prefixes, that is, prefixes whose individual argument is identical to the internal argument of the verb stem. The other contains verbs where the two entities are distinct. Finally, argument-sharing prefixes can be characterized as either restrictive or non-restrictive. The former introduce result states only compatible with a restricted subset of eventualities from the original extension of a verb stem. The latter denote states that can be brought about by every such an eventuality; they are also known under the name “pure perfectivizing” or “empty” prefixes. 

\keywords{event structure, lexical prefixes, result states}
}


\IfFileExists{../localcommands.tex}{
   \addbibresource{../localbibliography.bib}
   \usepackage{tabularx,multicol}
\usepackage{url}
\urlstyle{same}

 
\usepackage{langsci-optional}
\usepackage{langsci-lgr}
\usepackage{langsci-gb4e}

% ADDED BY VOLUME EDITORS

% \usepackage{langsci-osl} % macros specific to Open Slavic Linguistics; PLACED DIRECTLY IN localcommands.tex
\usepackage{siunitx}
\usepackage[linguistics]{forest}
% \usepackage{caption} %borik (?)
\usepackage{pifont} %geist

% 08_muellerreichau

% \usepackage{caption}
\usepackage{subcaption}
\usepackage{cleveref}

\usepackage{drs}
\usepackage{diagbox}

% 09_simik

\usepackage{multirow}
\usepackage{tabto}

% 10_stegovec

\usetikzlibrary{arrows,patterns}
\usepackage{listings}
% \usepackage{multirow}

% 01_gehrke-simik

% \usepackage{comment}

% 13_wagiel

% \usepackage{pifont} % for checkmarks (\ding{51}, \ding{55})


\usepackage{langsci-branding}

   % \newcommand*{\orcid}{}


\makeatletter
\let\thetitle\@title
\let\theauthor\@author
\makeatother

\newcommand{\togglepaper}[1][0]{
   \bibliography{../localbibliography}
   \papernote{\scriptsize\normalfont
     \theauthor.
     \titleTemp.
     To appear in:
     E. Di Tor \& Herr Rausgeberin (ed.).
     Booktitle in localcommands.tex.
     Berlin: Language Science Press. [preliminary page numbering]
   }
   \pagenumbering{roman}
   \setcounter{chapter}{#1}
   \addtocounter{chapter}{-1}
}

\newbool{bookcompile}
\booltrue{bookcompile}
\newcommand{\bookorchapter}[2]{\ifbool{bookcompile}{#1}{#2}}

\AtBeginDocument{\nogltOffset} %removes extra spacing before trans line in examples - SPECIFIC TO OPEN SLAVIC LINGUISTICS, PREVIOUSLY AGREED UPON


%%%%%%%%%%%%%%%%%%%%%%%%%%%%%%%%%%%%%%%%%%%%%%%%%%%%%%%%%%%%%%%%%%%%%
%%      File: langsci-osl.sty
%%    Author: Radek Simik and Language Science Press (http://langsci-press.org)
%%      Date: 2019-02-18
%%   Purpose: This file contains macros for the Open Slavic Linguistics at Language Science Press series and is included 
%%            into langscibook.cls
%%  Language: LaTeX
%%   Licence: 
%%%%%%%%%%%%%%%%%%%%%%%%%%%%%%%%%%%%%%%%%%%%%%%%%%%%%%%%%%%%%%%%%%%%%

% FIXING GRAY

\definecolor{lsDOIGray}{cmyk}{0,0,0,0.45}

% PACKAGES REQUIRED

\RequirePackage{relsize}
\RequirePackage{stmaryrd}
\RequirePackage{array}
\RequirePackage{tabularx}

% % KEYWORDS

% \newcommand{\keywords}[1]{\noindent\textbf{Keywords:} {#1}}

% SEMANTICS MACROS - GENERAL FOR ALL PAPERS

\newcommand{\sib}[1]{$\llbracket${#1}$\rrbracket$} % = semantic interpretation brackets; puts text into double square brackets
\newcommand{\stb}[1]{\langle{#1}\rangle} % = semantic type brackets; puts stuff into semantic type brackets; necessary to use in the form $\st{}$
\newcommand{\cnst}[1]{\textsf{\relsize{-1}\MakeUppercase{#1}}} %creates sans-serif small-caps, used for typesetting logical constants which have no other appropriate symbols (such as Greek letters)

% MINUS HSPACE MACRO

\newcommand{\minsp}[1]{#1\hspace{-2.5pt}} % use for non-glossed elements in glossed examples (typically stars, brackets, slashes, ...)

% CITATION MACROS

\newcommand{\citeposst}[1]{\citeauthor{#1}'s (\citeyear{#1})} % produces Chomsky's (1995)
\newcommand{\citeposstpg}[2]{\citeauthor{#1}'s (\citeyear[#2]{#1})} % produces Chomsky's (1995: page)
\newcommand{\citepossalt}[1]{\citeauthor{#1}'s \citeyear{#1}} % produces Chomsky's 1995
\newcommand{\citepossaltpg}[2]{\citeauthor{#1}'s \citeyear[#2]{#1}} % produces Chomsky's 1995: page

% BARPLOT

\usepackage{pgfplots}
\pgfplotsset{compat=1.18} 
\newcommand{\barplot}[5]{%
  \begin{tikzpicture}
    \begin{axis}[
	xlabel={#1},  
	ylabel={#2}, 
	axis lines*=left, 
        width  = {#3\textwidth},
	height = .3\textheight,
    	nodes near coords, 
	xtick=data,
	x tick label style={},  
	ymin=0,
	symbolic x coords={#4},
	]
	\addplot+[ybar,lsRichGreen!80!black,fill=lsRichGreen] plot coordinates {
	    #5
	}; 
    \end{axis} 
  \end{tikzpicture} 
}



%%%%%%%%%%%%%%%%%%%%%%%%%%%
%%% 04_filip %%%

% \newcommand{\mC}[1]{\multicolumn{1}{c}{#1}} % multicolumn with a single column; ALREADY DEFINED

\usetikzlibrary{arrows, arrows.meta}
\usetikzlibrary{positioning, decorations, 
                decorations.pathreplacing, calligraphy}

%%%%%%%%%%%%%%%%%%%%%%%%%
%%% 08_muellerreichau %%%

\captionsetup[subfigure]{subrefformat=simple,labelformat=simple}
\renewcommand\thesubfigure{(\alph{subfigure})}

%%%%%%%%%%%%%%%%%%%%%%%%%%%
%%% 09_simik

\newcommand{\citepossts}[1]{\citeauthor{#1}' (\citeyear{#1})} % produces Rawlins' (2008)
\newcommand{\citepossalts}[1]{\citeauthor{#1}' \citeyear{#1}} % produces Rawlins' 2008


%%%%%%%%%%%%%%%%%%%%%%%%%%%%%%%
%%% 10_stegovec

\newcommand{\poscite}[1]{\citeauthor{#1}'s (\citeyear{#1})}

%extra commands
\newcommand{\fst}{\textsc{1p}}
\newcommand{\snd}{\textsc{2p}}
\newcommand{\trd}{\textsc{3p}}
\newcommand{\refl}{\textsc{refl}}
\newcommand{\excl}{\textsc{excl}}
\newcommand{\incl}{\textsc{incl}}
\newcommand{\sg}{\textsc{sg}}
\newcommand{\pl}{\textsc{pl}}
\newcommand{\dl}{\textsc{dl}}
\newcommand{\fem}{\textsc{f}}
\newcommand{\masc}{\textsc{m}}
\newcommand{\neut}{\textsc{n}}
\newcommand{\ImpCl}{IMPERATIVE}


\newcommand{\tikzmark}[1]{\tikz[overlay,remember picture] \node (#1) {};}
\forestset{M/.style={
		for tree={s sep=5mm, inner sep=0.5pt, l=5mm,
			calign=fixed edge angles,
			calign primary angle=-65,
			calign secondary angle=65},
	},
	Mid/.style={
		for tree={s sep=0mm, inner sep=0pt, l=0mm,
			calign=fixed edge angles,
			calign primary angle=-62.5,
			calign secondary angle=62.5},
	},
	word/.style={before typesetting nodes={content=\emph{##1}},
		no edge,l=0,for parent={l sep=0}},
	feat/.style={before typesetting nodes={content={{[}##1{]}}},
		no edge,l=0,for parent={l sep=0}},
	feat2/.style={before typesetting nodes={content={##1}},
		no edge,l=0,for parent={l sep=0}},
	llap/.style={
		tikz+={%
			\edef\forest@temp{\noexpand\node[\option{node options},
				anchor=base east,at=(.base east)]}%
			\forest@temp{#1\phantom{\option{environment}}};
		}
	},
	rlap/.style={
		tikz+={%
			\edef\forest@temp{\noexpand\node[\option{node options},
				anchor=base west,at=(.base west)]}%
			\forest@temp{\phantom{\option{environment}}#1};
		}
	}
}

\makeatletter
\newcommand{\coloruline}[2]{%
	\newcommand\temp@reduline{\bgroup\markoverwith
		{\textcolor{#1}{\rule[-0.5ex]{2pt}{0.4pt}}}\ULon}%
	\temp@reduline{#2}%
}

\newcommand{\coloruuline}[2]{%
	\UL@protected\def\temp@uuline{\leavevmode \bgroup
		\UL@setULdepth
		\ifx\UL@on\UL@onin \advance\ULdepth2.8\p@\fi
		\markoverwith{\textcolor{#1}{\lower\ULdepth\hbox
				{\kern-.03em\vbox{\hrule width.2em\kern1\p@\hrule}\kern-.03em}}}%
		\ULon}%
	\temp@uuline{#2}%
}

\newcommand{\coloruwave}[2]{%
	\UL@protected\def\temp@uwave{\leavevmode \bgroup 
		\ifdim \ULdepth=\maxdimen \ULdepth 3.5\p@
		\else \advance\ULdepth2\p@ 
		\fi \markoverwith{\textcolor{#1}{\lower\ULdepth\hbox{\sixly \char58}}}\ULon}
	\font\sixly=lasy6 % does not re-load if already loaded, so no memory drain.
	\temp@uwave{#2}%
}
\makeatother

%%%%%%%%%%%%%%%%%
%%% 13_wagiel %%%

\newcommand\irregularcircle[2]{% radius, irregularity
  \pgfextra {\pgfmathsetmacro\len{(#1)+rand*(#2)}}
  +(0:\len pt)
  \foreach \a in {10,20,...,350}{
    \pgfextra {\pgfmathsetmacro\len{(#1)+rand*(#2)}}
    -- +(\a:\len pt)
  } -- cycle
}


% IF POSSIBLE, CAN WE USE THIS SETTING JUST FOR 13_wagiel?

% \forestset{
% default preamble={for tree={s sep=10mm, inner sep=0, l=0}},
% fairly nice empty nodes/.style={
%         delay={where content={}{shape=coordinate,
%               for siblings={anchor=north}}{}},
% for tree={s sep=4mm}},
% }

   %% hyphenation points for line breaks
%% Normally, automatic hyphenation in LaTeX is very good
%% If a word is mis-hyphenated, add it to this file
%%
%% add information to TeX file before \begin{document} with:
%% %% hyphenation points for line breaks
%% Normally, automatic hyphenation in LaTeX is very good
%% If a word is mis-hyphenated, add it to this file
%%
%% add information to TeX file before \begin{document} with:
%% %% hyphenation points for line breaks
%% Normally, automatic hyphenation in LaTeX is very good
%% If a word is mis-hyphenated, add it to this file
%%
%% add information to TeX file before \begin{document} with:
%% \include{localhyphenation}
\hyphenation{
    par-a-digm
    Cart-wright %filip
    Truswell %docekal
    Shcher-bi-na %simik
    Mün-chen %gehrke
    Mo-ni-ka %gehrke
    spi-sov-né %gehrke
    li-te-ra-tur-no-go %gehrke
    russ-koj %gehrke
    so-ver-šen-no-go %gehrke
    russ-kom %gehrke
    sla-vjan-sko-mu %gehrke
    Zeit-schrift %gehrke
    Na-ta-sha %gehrke
    Zu-stands-pas-siv %gehrke
    Um-gangs-spra-che %gehrke
    Bra-so-vea-nu %geist
}

\hyphenation{
    par-a-digm
    Cart-wright %filip
    Truswell %docekal
    Shcher-bi-na %simik
    Mün-chen %gehrke
    Mo-ni-ka %gehrke
    spi-sov-né %gehrke
    li-te-ra-tur-no-go %gehrke
    russ-koj %gehrke
    so-ver-šen-no-go %gehrke
    russ-kom %gehrke
    sla-vjan-sko-mu %gehrke
    Zeit-schrift %gehrke
    Na-ta-sha %gehrke
    Zu-stands-pas-siv %gehrke
    Um-gangs-spra-che %gehrke
    Bra-so-vea-nu %geist
}

\hyphenation{
    par-a-digm
    Cart-wright %filip
    Truswell %docekal
    Shcher-bi-na %simik
    Mün-chen %gehrke
    Mo-ni-ka %gehrke
    spi-sov-né %gehrke
    li-te-ra-tur-no-go %gehrke
    russ-koj %gehrke
    so-ver-šen-no-go %gehrke
    russ-kom %gehrke
    sla-vjan-sko-mu %gehrke
    Zeit-schrift %gehrke
    Na-ta-sha %gehrke
    Zu-stands-pas-siv %gehrke
    Um-gangs-spra-che %gehrke
    Bra-so-vea-nu %geist
}

   \boolfalse{bookcompile}
   \togglepaper[11]%%chapternumber
}{}

\begin{document}


\maketitle

\section{Event structure and derivational morphology}\label{sec:0}
\hypertarget{RefHeadingToc164647615}{}

The goal of this chapter is two-fold. First, it aims at providing a critical overview of the existing literature on the event structure of various types of verbal predicates in Slavic languages, with special attention to the contribution of verbal derivational morphology to the subevental make-up of an eventuality description. Second, it will focus on two specific topics in the realm of event structure. One is event-structural correlates of (lexical) prefixation and secondary imperfectivization. The other is the contribution of lexical prefixes to the argument structure and interpretation of the entire predicate. 

Throughout the discussion I will consistently keep three theoretical notions apart. \textsc{Argument structure} of a predicate comprises information about its  syntactically projected nominal or clausal arguments and their thematic relations to the event argument. \textsc{Event structure} defines whether an eventuality that falls under an eventuality description is simplex or complex, that is, decomposed into subevental components, what descriptive properties of each component are, and how they are related to each other. To identify the event structure of a predicate, semanticists normally use diagnostics based on the observations that some subevental components can participate in certain processes or undergo certain operations independently from others. 

Both notions should be kept conceptually and representationally distinct from each other and from \textsc{eventuality type}, a partition of predicates into non-over\-lap\-ping classes like states, activities, accomplishments, and achievements, to take \citeposst{Vendler1957} or \citeposst{Dowty1979} systems as a point of departure.\footnote{Membership in such classes is defined by second-order properties of sentences like truth at a point, properties of predicates like the subinterval property (in temporal semantics; \citealt{Bennett.Partee1972,Taylor1977} and subsequent literature) or cumulativity and quantization (in event semantics; \citealt{Krifka1989,Krifka1992,Krifka1998}).} Even though there are clear correlations between these dimensions of the meaning of a predicate, they cannot be fully reduced to each other. The same argument structure can correspond to different eventuality types. For example, transitive verbs with the agent and theme can be either telic (e.g. \textit{open}) or atelic (e.g. \textit{push}). The same eventuality type can be associated with different event structures: both \textit{wipe the floor} and \textit{break the window} are telic, yet the latter, unlike the former, is event-structurally complex (\citealt{Levin.RappaportHovav1995,Rappaport.Hovav.Levin1998} and much further literature).

As \citeauthor{Levin.RappaportHovav1995} and \citeauthor{Rappaport.Hovav.Levin1998} indicate, event-structural complexity of %the 
predicates like \textit{break the window} has several grammatical manifestations. Such predicates disallow indefinite object alternation (cf. \REF{ex:tatevosov:new1a} and \REF{ex:tatevosov:new1b}%*\textit{John} \textit{broke.} and \textit{Terry wiped.}
), are unavailable in configurations with non-subcategorized objects (cf. \REF{ex:tatevosov:new2a} and \REF{ex:tatevosov:new2b}%*\textit{The clumsy child broke the beauty out of the vase.} and \textit{The child rubbed the tiredness out of his eyes.}
) and do not readily accept expressions that specify the result state of an eventuality (cf. \REF{ex:tatevosov:new3a} and \REF{ex:tatevosov:new3b}%*\textit{Kelly broke the dishes off the table.} and \textit{Kelly swept the leaves off the sidewalk.}
) (examples from \citealt{Rappaport.Hovav.Levin1998}: 102--103). 

\ea\label{ex:tatevosov:new1}
\ea[*]{John broke.}\label{ex:tatevosov:new1a}
\ex[]{Terry wiped.}\label{ex:tatevosov:new1b}
\z 
\z 

\ea\label{ex:tatevosov:new2}
\ea[*]{The clumsy child broke the beauty out of the vase.}\label{ex:tatevosov:new2a}
\ex[]{The child rubbed the tiredness out of his eyes.}\label{ex:tatevosov:new2b}
\z 
\z 

\ea\label{ex:tatevosov:new3}
\ea[*]{Kelly broke the dishes off the table.}\label{ex:tatevosov:new3a}
\ex[]{Kelly swept the leaves off the sidewalk.}\label{ex:tatevosov:new3b}
\z 
\z 

\noindent Putting technical details aside, \citeauthor{Levin.RappaportHovav1995} and \citeauthor{Rappaport.Hovav.Levin1998} attribute the difference between result predicates like \textit{break} and manner predicates like \textit{wipe} to the fact that the semantic representation (``lexical conceptual structure'') of the former but not of the latter contains a result state. %They argue that with a few additional assumptions, independently needed, the patterns illustrated above reduce to this difference. For instance, since \textit{off the table} introduces a result state (of being off the  table) it is only licit in those configuration where a result state is not projected by the verb itself. 

The eventuality type of a predicate, on the other hand, does not reduce to event structure. A predicate does not have to have a result state to be telic. \textit{Wipe the table} can be telic without there being a result state, either contributed by the verb or a resultative expression. I will take up issues surrounding event structural complexity in %a 
due order. 

The subsequent discussion is based on %the 
Russian material. Obviously, cross-lin\-guis\-tic variation within the Slavic family is substantial, both in terms of the structure and semantic characteristics of the verbal system. Nevertheless, I hope that the most basic generalizations I will try to establish in what follows for Russian would mostly extend to other Slavic languages, with relevant qualifications and adjustments.\footnote{See \textcitetv{chapters/05_gehrke} for extensive discussion of the cross-Slavic variation in the aspectual domain.}

These minimal clarifications having been made, in the next section I will offer an overview of the verbal derivational system of Russian. Relying on this background, in \sectref{tat:sec:2}--\sectref{tat:sec:3} I will address two central topics of this chapter. \sectref{tat:sec:2} presents the set of data that is best accounted for if one assumes that (verb phrases projected by) lexically prefixed perfective verbs and secondary imperfective verbs contain a result state in their semantic representation, but their simplex imperfective counterparts do not. In \sectref{tat:sec:3}, a more detailed typology of lexical prefixes is set up: the entire set of lexical prefixes is divided into subclasses defined by a number of criteria, including their compositionality, argument structure and type of interaction with a verb stem. 

\section{Derivational morphology}\label{sec:1}
\hypertarget{RefHeadingToc164647616}{}
This section focuses on two topics. First, it aims at surveying basic phenomena in the domain of derivational morphology of Russian. Second, it provides an overview of the existing literature where these phenomena are addressed from various perspectives. I will start with an outline of the Russian verbal system, which can be skipped by the reader familiar with this material. See also a recent survey in \citet{Borik2024} for a detailed and comprehensive overview of the topic. 

\subsection{Verb structure and interpretation}\label{sec:1.1}
\hypertarget{RefHeadingToc164647617}{}
A finite morphological verb in Russian  consists of the three elements in \REF{tat:ex:1}. 

\ea   Verb stem -- (thematic element) -- inflection\label{tat:ex:1}
\z 

A few examples are given in \REF{tat:ex:2}. 

\ea\label{tat:ex:2}
\ea \gll   del-aj-eš\\
    do-\textsc{th-npst.2sg}\\
\glt  `you are doing\textsuperscript{IPFV} / you do\textsuperscript{IPFV}'\label{ex:tatevosov:2a}
\ex \gll   pis-a-l \\
    write-\textsc{th-pst.m}\\
\glt  ‘I (masc.) / you (masc.) / he was writing\textsuperscript{IPFV} / would write\textsuperscript{IPFV}’\label{ex:tatevosov:2b}
\ex \gll   ris-ova-l-i \\
    paint-\textsc{th-pst-pl}\\
\glt  ‘they were painting\textsuperscript{IPFV} / would paint\textsuperscript{IPFV}’\label{ex:tatevosov:2c}
\ex \gll   sox-n-et \\
    dry-\textsc{th-npst.3sg}\\
\glt  ‘it is drying\textsuperscript{IPFV} / dries\textsuperscript{IPFV}’\label{ex:tatevosov:2d}
\ex \gll    ljub-i-l-a\\
    like-\textsc{th-pst-f}\\
\glt  ‘I (fem.) / you (fem.) / she liked\textsuperscript{IPFV}’\label{ex:tatevosov:2e}
\z
\z

\noindent As \REF{ex:tatevosov:2b}, \REF{ex:tatevosov:2c}, and \REF{ex:tatevosov:2e} indicate, in the past tense, which is diachronically a participle, the verb is inflected for number and gender; in the non-past tense, for person and number, as  in \REF{ex:tatevosov:2a} and \REF{ex:tatevosov:2d}. Following \citet{Jakobson1948}, many morphologists (\citealt{Halle1963,Lightner1965} and much subsequent literature) decompose person-number morphemes into a tense morpheme and person-number morpheme proper. Under such an analysis -\textit{eš} of \textit{delaeš} in \REF{ex:tatevosov:2a}, for example, would be treated as involving \textit{-e-} of the non-past/present tense and  \textit{-š} of the 2nd person singular.

Thematic elements define morphological classes of Russian verbs; similar clas\-ses can be found across other Slavic languages as well. Thus, \textit{del-a-} ‘do, make’ in \REF{ex:tatevosov:2a} and \textit{ris-ova-} ‘draw, paint’ in \REF{ex:tatevosov:2c} fall under different classes, since they combine with different thematic elements, \textit{-a-} and \textit{-ova-.}\footnote{In the past and non-past subparadigms, thematic elements normally surface in different forms, cf. \textit{del-}[aj]\textit{-eš} and \textit{del-}[a]\textit{-l}, \textit{ris-}[uj]\textit{-eš} and \textit{ris-}[ova]\textit{-l}. There is a tradition going back to \citet{Jakobson1948}, which assumes a derivational account of this variation, whereby the surface allomorphs are derived from the same underlying representation (e.g. \citealt{Halle1963,Lightner1965} and subsequent work, \citealt{HalleMatushansky2004,Itkin2007}). Other morphologists tend to treat “stem allomorphs” as lexically stored and account for their distribution across subparadigms by invoking an independent mechanism (e.g. \citealt{Brown1998,Rubach.Booij2001}).} Thematic elements are part of the structure of the vast majority of verbs; a very limited number of “athematic” verbs merge inflectional morphology directly on top of the root, cf. \textit{nes-l-a} ‘she carried’. For %the
ease of exposition, in the remainder of this chapter thematic elements will be treated as part of the stem in the example sentences and morpheme-by-morpheme interlinear glosses.

At least some of the thematic elements can be thought of as the realizations of grammatical morphemes that derive a verbal configuration from category-neutral roots in the style of \citet{Marantz1997}, as in \REF{tat:ex:3}, assuming that the verbalizing element is \textit{v} that takes (the projection of) a root as its complement; see \citet{Jablonska2007,Biskup2019,MismasSimonovic2021,Milosavljevic.Arsenijevic2022},  among others, for further discussion of the patterns attested in a number of Slavic languages.\footnote{\citet{Matushansky2009,Matushansky2021,Matushansky2024}, relying mostly on evidence from the morphology of secondary imperfectivization, offers compelling arguments that “thematic elements” are not structurally homogeneous: only a proper subset of them can be analyzed as verbalizers; other “thematic elements” attach higher up in the structure.}

\ea \textsc{Thematic elements as verbalizers}\label{tat:ex:3}

  [\textit{\textsubscript{v}}\textsubscript{P} … \textit{i} … [\textsubscript{${\surd}$P} … \textit{ljub}- … ]] 
\z

\noindent The verbs in \REF{tat:ex:2} contain a thematic element but no further instances of derivational morphology. In what follows I will call such verbs \textsc{derivationally minimal}, or \textsc{simplex}. In most Slavic languages, Russian among them, derivationally minimal verbs are subject to the following descriptive generalization that relates a verb and an aspectual configuration it can occur in: 

\eanoraggedright   \textsc{Simplex verbs}\\\label{ex:tatevosov:4}
Most derivationally minimal verbs can only occur in grammatically imperfective clauses (“simplex imperfective verbs”). A limited number of derivationally minimal verbs can only occur in perfective clauses (``simplex perfective verbs'').\footnote{A number of verbs, labeled “bi-aspectual verbs”, are not subject to \REF{ex:tatevosov:4} and are available in both perfective and imperfective configurations. Such verbs, e.g. \textit{issledova-t'} ‘study, explore’ or \textit{atakova-t’} ‘attack’ are typically thought of as lexically ambiguous; I know of no account that can derive “bi-aspectuality” from other characteristics of a verb. For the purposes of this chapter I put bi-aspectual verbs aside. For a comprehensive overview of existing theories of Russian aspect see \textcitetv{chapters/08_muellerreichau}.}
\z

\noindent \REF{ex:tatevosov:5} offers an example of the simplex imperfective verb \textit{pisa-t’} ‘write’:

\ea\label{ex:tatevosov:5}
\gll Kogda   Nadja   voš-l-a,   Volodja   pisa-l   pis’mo. \\
  when  Nadja  walk.in-\textsc{pst-f}  Volodja  write-\textsc{pst.m}  letter.\textsc{acc}\\

\ea   ‘When Nadja walked in, Volodja was writing a letter.’

\glt  ${\tau}(\textsc{walk in}) \subseteq_T \tau(\textsc{write})$\label{ex:tatevosov:5a}


\ex   Unavailable: ‘When Nadja walked in, Volodja  wrote a letter / had already written a letter.’

\glt  $\tau(\textsc{walk in}) <_T \tau(\textsc{write}); \tau(\textsc{walk in}) >_T \tau(\textsc{write})$\label{ex:tatevosov:5b}
\z
\z

\noindent The sentence in \REF{ex:tatevosov:5} comes out imperfective. It warrants the episodic interpretation in \REF{ex:tatevosov:5a}, where the topic time specified by a temporal clause, the interval in which Nadja’s walking in occurs, is included into the time of writing a letter.\footnote{In the quasi-formal notation in \REF{ex:tatevosov:5}--\REF{tat:ex:6}, $\tau$ is the temporal trace function, $<_T$ and  $\subseteq_T$ are the relations of temporal precedence and temporal inclusion, respectively; $\tau(\textsc{walk in}) <_T \tau(\textsc{solve})$ thus reads as “the time of walking in temporally precedes the time of solving”.}  This is the progressive construal. The perfective interpretation, marked by the non-overlap of walking in and (culminated) writing of a letter in \REF{ex:tatevosov:5b}, is unavailable.

The sentence in \REF{tat:ex:6} with the simplex perfective verb \textit{reši-t’} ‘solve’ shows the opposite range of interpretations: 

\ea\label{tat:ex:6}
\gll Kogda  Nadja  voš-l-a,  Volodja  reši-l  zadač-u.  \\
  when  Nadja  walk.in-\textsc{pst-f}  Volodja  solve-\textsc{pst.m}  quiz-\textsc{acc}\\
\ea   ‘When Nadja walked in, Volodja  solved a quiz  / had already solved a quiz.’

    $\tau(\textsc{walk in}) <_T \tau(\textsc{solve})$; \textsuperscript{?}$\tau(\textsc{walk in}) >_T \tau(\textsc{solve})$; \textsuperscript{?}$\tau(\textsc{walk in}) =_T \tau(\textsc{solve})$\label{ex:tatevosov:6a}

  \ex      Unavailable: ‘When Nadja walked in, Volodja was solving a quiz.’

    $\tau(\textsc{walk in}) \subseteq_T \tau(\textsc{solve})$\label{ex:tatevosov:6b}
\z
\z

\noindent The preferred interpretation of \REF{tat:ex:6} is the one where a complete culminated solving of the quiz follows Nadja’s walking in. An alternative interpretation, less easily available in the null context, is the past-over-past reading where the temporal sequencing is reverse. Yet one more option is a complete temporal overlap of walking in and solving a quiz. These readings are shown in \REF{ex:tatevosov:6a}. Crucially, \REF{tat:ex:6} cannot have the progressive construal in \REF{ex:tatevosov:6b}.

The characterizing property of the aspectual system of most Slavic languages is that all inflected forms of the verbs like \textit{pisa-t’} ‘write’ must occur in imperfective clauses while the forms of \textit{reši-t’} are restricted to perfective clauses. This \textsc{aspectual invariance} has lead aspectologists to the suggestion, standardly accepted since the 19th century, that grammatical aspect is a lexically specified property of the verb in Slavic, \textit{pisa-t’} ‘write’ being lexically imperfective, \textit{reši-t’} lexically perfective. \citet{Tatevosov2011,Tatevosov2015b,Tatevosov2021} argues that this view is fundamentally misguided, and that “aspectual morphology” and “aspectual interpretation” are located at a substantial structural distance, as shown in \REF{tat:ex:7}. 

\ea   \textsc{Verb phrases are aspectless}\label{tat:ex:7}

\ea{} [  …. IPFV …   [\textsubscript{VP} … \textit{pisa}- … ] ]\label{ex:tatevosov:7a}

\ex{} [  …. PFV …   [\textsubscript{VP} … \textit{reši}- … ] ]\label{ex:tatevosov:7b}
\z
\z

\noindent Generalizing over \REF{ex:tatevosov:7a}--\REF{ex:tatevosov:7b}, I will be assuming that at the point of derivation where the verb, either derivationally minimal or not, is fully assembled from the pieces of morphology, no semantic aspect is part of its meaning. This assumption is made explicit in \sectref{tat:sec:2}--\sectref{tat:sec:3}, where all VPs/\textit{v}Ps are treated as aspectless predicates of events.\footnote{For an apporoach along similar lines see \textcitetv{chapters/08_muellerreichau}.}

However, in order to minimize introducing new terminology I will continue to follow the tradition in using \textit{imperfective verbs} and \textit{perfective verbs} as descriptive terms. The reader should keep in mind, however, that \textit{an (im)perfective verb} should be taken as a shorthand for \textit{a verb that can only occur in an (im)per\-fec\-tive configuration if it undergoes no further morphological derivation}. 

Simplex verbs can be combined with various pieces of derivational morphology, characterized below.

\subsection{The ingredients}\label{sec:1.2}
\hypertarget{RefHeadingToc164647618}{}
Russian, as well as other Slavic languages, has an extensive inventory of verbal prefixes, listed in \REF{ex:tatevosov:8}. \REF{ex:tatevosov:8}, according to \citet[][\S 850]{Svedova1980}, forms an exhaustive list:

\ea   \textsc{Verbal prefixes in Russian}\label{ex:tatevosov:8}

  \textit{v(o)-}, \textit{v(o)z(o)-}, \textit{vy-}, \textit{de(z)-}, \textit{dis-}, \textit{do-}, \textit{za-}, \textit{iz(o)-}, \textit{na-, nad(o)-}, \textit{nedo-}, \textit{niz(o)-}, \textit{{о}-}, \textit{{о}b(o)-}, \textit{ot({о})-}, \textit{pere-}, \textit{p{о}-}, \textit{pod(o)-}, \textit{pre-}, \textit{pred(o)-}, \textit{pri}-, \textit{pro}-, \textit{raz(o)-}, \textit{re}-, \textit{s(o)-}, \textit{u-}
\z

\noindent Suffixal derivational morphology comes in two pieces: the so-called secondary imperfective morpheme -\textit{yva}- and the semelfactive morpheme -\textit{nu}-. The latter is illustrated in \REF{tat:ex:9}. 

\ea   \textsc{Semelfactives}\label{tat:ex:9}
\ea \gll tolk-a-t’ \\
push-\textsc{th-inf}\\
\glt `push, give pushes' (\textsc{ipfv})\label{ex:tatevosov:9a}
\ex  \gll tolk-nu-t’ \\
push-\textsc{smfct-inf}\\
\glt  `give a push' (\textsc{pfv})  \label{ex:tatevosov:9b}
\ex tolk-a-nu-t’ \\
push-\textsc{th-smfct-inf}\\
\glt  `give a push' (\textsc{pfv})  \label{ex:tatevosov:9c}
\z
\z 

\noindent Morphologically, the semelfactive allows for two options: -\textit{nu-} either merges directly with the root, as in \REF{ex:tatevosov:9b}, or attaches on top of the thematic element, as in \REF{ex:tatevosov:9c}. In both cases, semantically, a derived predicate denotes atomic pushing eventualities, whereas the non-derived ‘push’ in \REF{ex:tatevosov:9a} has both atoms and their sums in its extension.

In terms of grammatical aspect, the semelfactive is subject to the following generalization:\footnote{The semelfactive morpheme -\textit{nu}- should be distinguished from the thematic element -\textit{nu}- which is part of the derivation of a number of denominal verbs, e.g. \textit{soxnu-t’} ‘dry’. Semantically, the latter are not subject to \REF{tat:ex:10} and fall under the simplex imperfective class. Morphologically, the difference is manifested in the past tense, where the thematic element -\textit{nu}-, unlike the semelfactive -\textit{nu}-, has no surface realization, cf. \textit{sox-(*nu)-l-a} ‘she dried’ and \textit{tolk-*(nu)-l-a} ‘she gave a push’.}  

\ea   \textsc{Aspect of the semelfactive}\label{tat:ex:10}
\ea The semelfactive affix requires its complement to be imperfective.\label{ex:tatevosov:10a}
\ex The semelfactive is perfective.\label{ex:tatevosov:10b}
\z\z

\noindent According to \REF{ex:tatevosov:10a}, the semelfactive must attach to verb stems like \textit{tolka-} ‘push’ in \REF{ex:tatevosov:9a}, which come out imperfective if they stay derivationally minimal. According to \REF{ex:tatevosov:10b}, the \textit{nu}-verbs themselves can only occur in perfective configurations. This is what happens in \REF{ex:tatevosov:9b}--\REF{ex:tatevosov:9c}.

The semelfactive, restrictions on its derivation, differences between semelfatives like \REF{ex:tatevosov:9b} and \REF{ex:tatevosov:9c} remain one of the least explored topics in Slavic aspectology. For more information about the semelfactive the reader can refer to \citet{Markman2008,Janda.Makarova2009,Gribanova2013,TaraldsenMedova.Wiland2019,Matushansky2024}, and literature therein.

In the rest of this section I will mostly focus on prefixation and secondary imperfectivization. Prefixes in Slavic languages are frequently described as “perfecitivizing”. Such a description should not be taken too literally, however. The output of prefixation must be a perfective verb, as captured by the generalization in \REF{tat:ex:11}. 

\eanoraggedright   \textsc{Aspect of prefixed verbs}\label{tat:ex:11}

  Prefixed verbs that undergo no further morphological derivation are perfective.
\z

\noindent However, although the output of prefixation is always perfective, the input to prefixation does not have to be imperfective. \REF{ex:tatevosov:12a}--\REF{ex:tatevosov:12d} are illustrations of a prefix being merged with a perfective verb stem:\footnote{Throughout the chapter \textsc{lp} and \textsc{slp} stand for ``lexical prefixes'' and ``superlexical prefixes'', discussed in more detail in \sectref{sec:1.3} and subsequent sections.}

\ea   \textsc{Prefixes in combination with a perfective verb}\label{tat:ex:12}
\ea \gll \minsp{[} ot- \minsp{[} reši]\textsuperscript{PFV}]\textsuperscript{PFV} -t’  \\
{} \textsc{lp} {} solve \textsc{inf}\\
\glt  ‘dismiss\textsuperscript{PFV}’  \label{ex:tatevosov:12a}
\ex  \gll \minsp{[} pod- \minsp{[} kupi]\textsuperscript{PFV}]\textsuperscript{PFV}{} -t’  \\  
{} \textsc{lp} {} buy \textsc{inf}\\
\glt  ‘bribe\textsuperscript{PFV}’ \label{ex:tatevosov:12b}
\ex \gll \minsp{[} pere- \minsp{[} za- pisa]\textsuperscript{PFV}]\textsuperscript{PFV} -t’ \\
{} \textsc{slp} {} \textsc{lp} write \textsc{inf}   \\
\glt  ‘record\textsuperscript{PFV} again’\label{ex:tatevosov:12c}
\ex \gll   \minsp{[} pod- \minsp{[} za- by]\textsuperscript{PFV}]\textsuperscript{PFV} -t’\\
    {} \textsc{slp} {} \textsc{lp} be \textsc{inf}\\
\glt   ‘forget\textsuperscript{PFV} partially’\label{ex:tatevosov:12d}
\z
\z

\noindent In \REF{ex:tatevosov:12a}--\REF{ex:tatevosov:12b}, the prefixes \textit{ot-} and \textit{pod-} combine with simplex perfective verbs \textit{reši-t’} ‘solve’ and \textit{kupi-t’} ‘buy’. In \REF{ex:tatevosov:12c}--\REF{ex:tatevosov:12d}, \textit{pere}- and \textit{pod-} are affixed to \textit{za-pisa-t’} ‘record’ and \textit{za-by-t’} ‘forget’, which are perfective by virtue of being prefixed and not further suffixed. Therefore, in \REF{ex:tatevosov:12a}--\REF{ex:tatevosov:12d} no “perfectivization” occurs, since the input to prefixation is perfective already.

\REF{ex:tatevosov:12c}--\REF{ex:tatevosov:12d} independently illustrate another prominent property of the Slavic verbal derivational morphology: \textsc{multiple prefixation}. While \REF{ex:tatevosov:12c}--\REF{ex:tatevosov:12d} contain two prefixes each, \REF{ex:tatevosov:13a} is an example of a verb stem with triple prefixation. \REF{ex:tatevosov:13b} shows its derivation in a step-wise manner.

\ea   \textsc{Triple prefixation}\label{tat:ex:13}
\ea
\gll do-pere-za-pisa-t’ \\
    \textsc{slp-slp-lp}-write-\textsc{inf}\\
\glt  ‘complete recording again’\label{ex:tatevosov:13a}
\ex
      \textit{pisa}- ‘write’ \\$\rightarrow$ \textit{za}\textit{-pisa-} ‘record’ \\$\rightarrow$ \textit{pere}\textit{-za-pisa-} ‘record again’ \\$\rightarrow$
\textit{do}\textit{-pere-za-pisa-} ‘complete recording again’\label{ex:tatevosov:13b}
\z
\z

\noindent Constraints on multiple prefixation in Russian have been explored in \citet{Tatevosov2009,Tatevosov2013,Mueller-Reichau2021}, among others.

The secondary imperfective morpheme -\textit{yva}- outputs an imperfective verb:\footnote{The final segment of -\textit{yva}- can be identified with the thematic element \textit{-a(j)-} found in a number of non-derived stems, e.g. \textit{dela-t’} ‘do, make’, see, specifically, \citet{Isacenko1960}: 176--181. For better readability, this information is left implicit in the morpheme-by-morpheme interlinear translation.}  

\eanoraggedright   \textsc{Aspect of secondary imperfective verbs}\label{tat:ex:14}

  \textit{yva}-verbs that undergo no further morphological derivation are imperfective.
\z

\noindent The examples in \REF{ex:tatevosov:15a}--\REF{ex:tatevosov:15d} show the imperfective counterparts of the simplex perfective verb \textit{reši-t'} ‘solve’, prefixed perfective verb \textit{pro-čita-t’} ‘read’, and perfective verbs from \REF{ex:tatevosov:12b} and \REF{ex:tatevosov:12d}. Note that the secondary imperfective morpheme has three surface realizations: -\textit{va}- after vowel-final stems, but -\textit{yva}- or stressed \textit{-á-} after consonant final stems. Throughout the chapter the label “YVA” is used to cover all of them in the morpheme-by-morpheme interlinear translation.\footnote{\citet{Matushansky2009} offers a theory that derives all the three allomorphs from the same underlying representation, /ʊ/, “the third yer”, followed by the thematic element. In verb stems like \textit{(pod)-kup-á-t’} in \REF{ex:tatevosov:15c}, /ʊ/ gets deleted, and its presence in the derivation can only be detected by the obligatory stress on the following thematic element. An extension of this theory is offered in \citealt{Matushansky2021}; an alternative phonological theory has been developed by \citet{Enguehard2017}.}

\ea   \textsc{The -\textit{yva-} morpheme in combination with perfective verbs}\label{tat:ex:15}
\ea \gll  \minsp{[[} reš]\textsuperscript{PFV} -a]\textsuperscript{IPFV} -t’   \\
 {} solve \textsc{yva} \textsc{inf}\\
 \glt ‘solve\textsuperscript{IPFV}’\label{ex:tatevosov:15a}
\ex \gll \minsp{[[} pro- čit]\textsuperscript{PFV} -yva]\textsuperscript{IPFV} -t’  \\ 
{} \textsc{lp} read \textsc{yva} \textsc{inf} \\
\glt ‘read\textsuperscript{IPFV}’\label{ex:tatevosov:15b}
\ex \gll \minsp{[} pod- kup]\textsuperscript{PFV} -a]\textsuperscript{IPFV} -t’\\
{} \textsc{lp} buy \textsc{yva} \textsc{inf} \\
\glt  ‘bribe\textsuperscript{IPFV}’\label{ex:tatevosov:15c}
\ex \gll   \minsp{[[} pod- za- by]\textsuperscript{PFV} -va]\textsuperscript{IPFV} -t’\\
{} \textsc{slp} \textsc{lp} be \textsc{yva} \textsc{inf}\\
\glt   ‘forget\textsuperscript{IPFV} partially’\label{ex:tatevosov:15d}
\z
\z

\noindent One of the issues in the derivational morphology of Slavic that does not seem to have been completely settled comes from the verbs like \REF{ex:tatevosov:16a}--\REF{ex:tatevosov:16c} derived by the affixation of -\textit{yva-} to an imperfective stem:

\ea   \textsc{The -\textit{yva}- morpheme in combination with imperfective verbs}\label{tat:ex:16}
\ea \gll \minsp{[[} siž]\textsuperscript{IPFV} -iva]\textsuperscript{IPFV} -t’ \\
{} sit \textsc{yva} \textsc{inf} \\
\glt  ‘sit occasionally’\label{ex:tatevosov:16a}
\ex \gll \minsp{[[} igr]\textsuperscript{IPFV} -yva]\textsuperscript{IPFV} -t’ \\ 
{} play \textsc{yva} \textsc{inf} \\
\glt  ‘play occasionally’   \label{ex:tatevosov:16b}
\ex \minsp{[[} pis]\textsuperscript{IPFV} -yva]\textsuperscript{IPFV} -t’\\
{} write \textsc{yva} \textsc{inf} \\
\glt  ‘write occasionally’\label{ex:tatevosov:16c}
\z
\z 

\noindent \citet{Filip.Carlson1997} discuss Czech counterparts of \REF{ex:tatevosov:16a}--\REF{ex:tatevosov:16c} and argue for a theory that assumes two distinct \textit{yva’}s: the imperfective one, combining with perfective stems, as in \REF{ex:tatevosov:15a}--\REF{ex:tatevosov:15d},  and the generic one, which only attaches to the imperfective stems, as in \REF{ex:tatevosov:16a}--\REF{ex:tatevosov:16c}.

%One may wonder whether the two \textit{yva’}s can ultimately be reduced to the same underlying semantics, given their complementary distribution across the aspectual stems and the fact that one of the readings of the “imperfective \textit{yva}” is the habitual, or generic one.  

Given the distribution of the two \textit{yva}’s across aspectual stems ([ ... ]\textsuperscript{PFV} $+$ “imperfective \textit{yva}” vs. [ ... ]\textsuperscript{IPFV} $+$ “generic \textit{yva}” ) and the fact that one of the readings of the  “imperfective \textit{yva}” is the habitual / generic one, it looks desirable to reduce them to the same semantics. 
At the moment, however, no compelling semantic theory unifying uses like \REF{tat:ex:15} and \REF{tat:ex:16} has been proposed. 

\subsection{Lexical and superlexical prefixes: an overview}\label{sec:1.3}
\hypertarget{RefHeadingToc164647619}{}
In \REF{ex:tatevosov:8}, 26 phonological strings that appear in the left periphery of a verb stem are listed, some with phonologically conditioned variants like \textit{iz}- and \textit{izo}-. The actual number of prefixal lexical items may be considerably higher, since many elements from \REF{ex:tatevosov:8} come in different varieties, normally identified as “meanings” / “uses” of a prefix. One example is the prefix \textit{pere-} in \REF{tat:ex:17} (see also \citealt{Kagan15}: 118 ff. and literature therein):

\ea\label{tat:ex:17}
\ea     \textsc{Distributive} \textit{pere}-

\gll    pere-lovi-t’  vsex  vrag-ov  \\
    \textsc{slp}-catch-\textsc{inf}  all.\textsc{acc}  enemy-\textsc{acc.pl}\\
\glt    ‘catch all the enemies one after another’\label{ex:tatevosov:17a}
\ex    \textsc{Repetitive} \textit{pere}-

\gll    pere-pisa-t’  stat-ju\\
\textsc{slp}-write-\textsc{inf}  article-\textsc{acc}\\
\glt    ‘re-write the article’\label{ex:tatevosov:17b}
\ex    \textit{Pere}- \textsc{of excess}  

\gll pere-jes-t’\\
\textsc{slp}-eat-\textsc{inf} \\
\glt    ‘overeat’\label{ex:tatevosov:17c}
\ex   \textit{Pere}- \textsc{of comparison}

\gll  pere-ži-t’  dedušk-u\\
\textsc{slp}-live-\textsc{inf}  grandfather-\textsc{acc}\\
\glt    ‘outlive the grandfather’\label{ex:tatevosov:17d}
\z
\z

\noindent \REF{ex:tatevosov:17a}--\REF{ex:tatevosov:17d} show the readings of \textit{pere-} most readily accessible in the null context for a particular combination of \textit{pere-} and a non-derived stem: distributive \textit{pere}-, \REF{ex:tatevosov:17a}, repetitive \textit{pere}-, \REF{ex:tatevosov:17b}, \textit{pere}- of excess, \REF{ex:tatevosov:17c}, and \textit{pere}- of comparison, \REF{ex:tatevosov:17d}.

Different readings are not complementarily distributed over lexical verbs. One can easily find examples where the same combination allows for several interpretations: 

\ea   pere-čita-t’\\\label{tat:ex:18}
  \textsc{slp}-read-\textsc{inf}\\
\ea ‘read one by one’ (distributive)
\ex ‘re-read’ (repetitive)
\ex ‘read too much’ (excess)
\z
\z

\noindent In a similar fashion, the grammar of Russian \citep{Svedova1980} lists 5 meanings for \textit{iz}-, 8 meanings for \textit{ot}-, 9 meanings for \textit{pod}-, 10 meanings for \textit{do-}, and so on. 

The question about examples like \REF{tat:ex:17}--\REF{tat:ex:18} is whether different “meanings” can be reduced to the same underlying semantic representation. There have been many attempts at a unification within different approaches and frameworks, which target some or all of the meanings of different prefixes. One of the most recent studies is \citet{Kagan15}, who offered a degree semantic account for 19 prefixes. 

For obvious reasons, such a unification is highly desirable. If every submeaning of a prefix needs to be listed, the amount of distinct prefixal lexical items starts exceeding several hundred, which makes the acquisition task for a learner of a Slavic language too complicated. However, one can be somewhat skeptical as to whether existing accounts can offer a comprehensive, complete, and fully explicit semantic theory of prefixation where every prefix in \REF{ex:tatevosov:8} is assigned a general underlying meaning that gives rise to the attested surface diversity.\footnote{Remarkably, \citeauthor{Kagan15}'s (\citeyear{Kagan15}:123) account for \textit{pere}- does not cover the repetitive \textit{pere}- in \REF{ex:tatevosov:17b}. There is no obvious way of subsuming its interpretation under the general format she proposes, in which \textit{pere-} indicates that the two relevant interval degrees stay in the inclusion relation.}

Moreover, there are serious empirical reasons to believe that at least some of the meanings of the prefixes in \REF{ex:tatevosov:8} are irreducible to others. Specifically, there is robust evidence that prefixes show radically different morphosyntactic distribution and interpretational properties under different meanings.

Starting at least from \citet{BabkoMalaya1999}, it has been recognized that Slavic prefixes can be divided into two types normally referred to as \textsc{lexical}, or internal, and \textsc{superlexical}, or external prefixes; see e.g. \citet{BabkoMalaya1999,Ramchand2004,Romanova2004,Romanova2007,Svenonius2004,Svenonius2008,DiSciullo.Slabakova2005,Gehrke2008,Tatevosov2008,Tatevosov2009,Tatevosov2013,Zaucer2009,Gribanova2013,Gribanova2015,Tolskaya2015,Biskup2019} on Russian; \citet{Progovac2002, Milicevic2004,Arsenijevic2006,Arsenijevic2007a, Arsenijevic2007b, Arsenijevic2012} on Serbo-Croatian; \citet{Istratkova2004,DiSciullo.Slabakova2005,Slabakova1997,Slabakova2005,Markova2011} on Bulgarian; \citet{Gehrke2008,Biskup2009,Biskup2012,Biskup2015,Biskup2019} on Czech;   and \citet{Zaucer2009,Zaucer2010,Zaucer2012, Zaucer2013} on Slovenian.\footnote{\citeposst{BabkoMalaya1999} observations about Russian parallel in many significant respects \citeposst{DiSciullo1997} work on external and internal prefixation in Romance, specifically, in French and Italian.}

The distribution of prefixal items in \REF{ex:tatevosov:8} across these two classes can be captured by the following generalization:

\eanoraggedright   \textsc{Lexical prefixes (LPs) and superlexical prefixes (SLPs)}\label{tat:ex:19}

   Every phonological string from \REF{ex:tatevosov:8} can be a lexical prefix. Some of the phonological strings in \REF{ex:tatevosov:8} can also be superlexical prefixes.
\z

\noindent According to \REF{tat:ex:19}, many prefixes have both lexical and superlexical manifestations. Different authors (\citealt{BabkoMalaya1999,Svenonius2004,Romanova2007,Gehrke2008,Tolskaya2007,Tolskaya2015,Tatevosov2009,Tatevosov2013,Biskup2019}) identify overlapping but not identical inventories of SLPs. The most inclusive list is shown in \REF{tat:ex:20}.

\ea   \textsc{Superlexical prefixes}\label{tat:ex:20}

  \textit{do-} (completive)\textit{, za-} (inceptive)\textit{, iz(o)-} (totality), \textit{na-} (cumulative) \textit{ot}- (terminative)\textit{, pere-} (distributive), \textit{pere-} (repetitive), \textit{pere-} (excessive)\textit{, p{о}-} (delimitative)\textit{, p{о}-} (distributive)\textit{, pod(o)-} (attenuative)\textit{, pri-} (attenuative)\textit{, pro-} (perdurative)\textit{, raz(o)-} (evolutive)
\z

\noindent As \REF{tat:ex:19}--\REF{tat:ex:20} suggest, a prefix can only be characterized as lexical or superlexical under a particular construal. For instance, \textit{po}- is superlexical as a delimitative or distributive prefix, but lexical elsewhere. Similarly, while the varieties of \textit{pere’}s in \REF{tat:ex:17}--\REF{tat:ex:18} are all superlexical, one can easily find instances of the lexical \textit{pere}-. It is for this reason that completely unifying the prefixal meanings and treating every prefix as a single lexical item may ultimately prove impossible. See \sectref{sec:3.3} for further discussion. 

\subsection{Lexical and superlexical prefixes: a few diagnostics}\label{sec:1.4}
\hypertarget{RefHeadingToc164647620}{}
LPs have a number of characterizing properties that distinguish them from SLPs. Simplifying somewhat, one can group these properties into three clusters: semantics, argument structure, morphosyntax. 

On the morphosyntactic side, the key observation is that in case of multiple prefixation, SLPs occur outside of LPs, never inside. A few illustrations appear in \REF{ex:tatevosov:21}, where bracketing indicates morphological constituency.

\ea   \textsc{SLP outside LP}  \label{ex:tatevosov:21}%    LP outside SLP
\ea \gll \minsp{[} na- \minsp{[[} za- bi] -va]] -t’\\
    {} \textsc{slp} {} \textsc{lp} hit \textsc{yva} \textsc{inf}\\\hfill *za-na-bi-va-t’
\glt  ‘hammer a lot of (e.g., nails)’\label{ex:tatevosov:21a}
\ex \gll    \minsp{[} po- \minsp{[[} o- pis] -yva]] -t’\\
    {} \textsc{slp} {} \textsc{lp} write \textsc{yva} \textsc{inf}\\\hfill *o-po-pis-yva-t’
\glt    ‘describe for a while’\label{ex:tatevosov:21b}
\ex \gll    \minsp{[} za- \minsp{[[} vz- dyx] -a]] -t’  \\
{} \textsc{slp} {} \textsc{lp} breath \textsc{yva} \textsc{inf}\\\hfill  *v(o)z-za-dyx-a-t’ 
\glt    ‘start sighing’\label{ex:tatevosov:21c}
\ex \gll   \minsp{[} po- \minsp{[[} vy -bras] -yva]] -t’   \\
{} \textsc{slp} {} \textsc{lp} throw \textsc{yva} \textsc{inf}\\\hfill *vy-po-bras-yva-t’
\glt ‘throw out one by one’\label{ex:tatevosov:21d}
\ex \gll   \minsp{[} do- \minsp{[} so- bra]] -t’   \\
{} \textsc{slp} {} \textsc{lp} take \textsc{inf}\\\hfill *so-do-bra-t’
\glt ‘finish collecting / assembling’\label{ex:tatevosov:21e}
\ex \gll   \minsp{[} pere- \minsp{[} za- pusti]] -t’   \\
{} \textsc{slp} {} \textsc{lp} let \textsc{inf}\\\hfill *za-pere-pusti-t’
\glt ‘re-start’\label{ex:tatevosov:21f}
\z
\z

\noindent \REF{ex:tatevosov:21a}--\REF{ex:tatevosov:21f} exemplify the cumulative \textit{na}-, delimitative \textit{po}-, inceptive \textit{za}-, distributive \textit{po}-, completive \textit{do}- and repetitive \textit{pere}-. All of them merge with a stem which a lexical prefix is part of. The right column in \REF{ex:tatevosov:21} represents attempts to merge the prefixes in the reverse order, which invariably produce an ungrammatical output.

Semantically, as \citet[][76]{BabkoMalaya1999} points out, SLPs tend to have “systematic and predicable meanings”; LPs, in contrast, are associated with “idiosyncratic meanings”. \citet[][230]{Svenonius2004} indicates that SLPs contribute “temporal or quantizing” meanings, whereas the contribution of LPs is normally spatial or resultative meanings. Similar descriptions can be found in other work on prefixation.\footnote{To do justice to the earlier work on Slavic verbs, it should be pointed out that similar observations had started circulating in the literature long before Babko-Malaya’s work; see e.g. the discussion of prefixal “qualifiers” and “modifiers” in \citet{Isacenko1960}: 222--224.} 

%As an illustration consider \textit{za}-, \textit{po}-, \textit{do}- and two different \textit{pere}-’s in \REF{ex:tatevosov:22}.

%\ea   \textsc{Meaning: SLPs}\label{ex:tatevosov:22}
%\ea \textsc{Inceptive}:\label{ex:tatevosov:22a}

%\textit{za-bole-t’} ‘get sick’, \textit{za-pisa-t’} ‘start writing’     
%\ex \textsc{Delimitative}:\label{ex:tatevosov:22b}

%\textit{po-bole-t’} ‘be sick for a while’, \textit{po-pisa-t’} ‘write for a while’ 
%\ex \textsc{Distributive}:\label{ex:tatevosov:22c}

%\textit{pere-kida-t’} ‘throw one by one’, \textit{pere-kusa-t’} ‘bite one by one’
%\z
%\z 

%\noindent The contribution of each of the prefixes in \REF{tat:ex:20} is easily identifiable, and the combination of a prefix with a stem is compositional. The inceptive verbs in \REF{ex:tatevosov:22a} describe the inception of a process denoted by the non-derived stem. The meaning of the delimitative prefix \textit{po}- in \REF{ex:tatevosov:22b} is taken to be similar to that of measure adverbials (\citealt{Pinon1994,Filip2000}): the prefix indicates that one of the gradable properties of an eventuality (typically, its temporal extent), falls below the standard of comparison or expectation value. The distributive in \REF{ex:tatevosov:22c} describes a plurality of events where each atomic subevent  is associated with a mereological part of (the denotation of) the object.

%In contrast, corresponding lexical prefixes resist simple and straightforward semantic generalizations:   

%\ea    \textsc{Meaning: LPs}\label{ex:tatevosov:23}

%   \ea   \textit{pisa-t’} ‘write’ $\rightarrow$    \textit{za-pisa-t’} ‘record, write down’\label{ex:tatevosov:23a}

%  \textit{ry-t’} ‘dig’  $\rightarrow$   \textit{za-ry-t’} ‘dig in’

%  \textit{stroi-t’} ‘build’  $\rightarrow$   \textit{za-stroi-t’} ‘build up’
%\ex      \textit{ljubi-t’} ‘love’ $\rightarrow$  \textit{po-ljubi-t’} ‘fall in love’\label{ex:tatevosov:23b}

%  \textit{trjasti} ‘shake’ $\rightarrow$ \textit{po-trjasti} ‘strike’

%  (\textit{vesi-t’} ‘weigh’) $\rightarrow$ \textit{po-vesi-t’} ‘hang’
%\ex      \textit{kinu-t’} ‘throw’ $\rightarrow$ \textit{pere-kinu-t’} ‘throw across’\label{ex:tatevosov:23c}
%
%(*\textit{kusi-t’} ‘bite’) $\rightarrow$ \textit{pere-kusi-t’} ‘bite in half’
%\z
%\z

%\noindent In \REF{ex:tatevosov:23a}--\REF{ex:tatevosov:23c} the contribution of a lexical prefix varies across different stems. Intuitively, each of the prefixed verbs in these paradigms differs from a corresponding simplex verb in a different way. Some of the prefixed items fail to entail a  non-prefixed counterpart: in order to strike (\textit{po-trjasti} in \REF{ex:tatevosov:23b}), one does not normally have to shake. Moreover, at least synchronically, many lexically prefixed verbs do not have a semantically related unprefixed counterpart (cf. \textit{vesi-t’} ‘weigh’ and \textit{povesi-t’} ‘hang’ in \REF{ex:tatevosov:23b}). Some do not have such a counterpart at all (e.g. *\textit{kusi-t’} in \REF{ex:tatevosov:23c}).

Another set of observations that tell LPs and SLPs apart concern the interaction of prefixation and argument structure of a non-derived stem. SLPs do not add an argument to the root verb or affect thematic properties of an object (\citealt{Romanova2004,Romanova2007}; \citealt{Ramchand2004}), do not license unselected objects (\citealt{BabkoMalaya1999,Svenonius2004,Romanova2007}; but see \citealt{Zaucer2009,Zaucer2010,Zaucer2012}, and have no consequences for the indefinite object alternation (\citealt{BabkoMalaya1999,Svenonius2004}). LPs can have the opposite properties.\footnote{As \citet[][20--23]{Biskup2019} emphasizes, SLPs can nevertheless impose certain semantic and syntactic restrictions on their arguments. For instance, distributive SLPs need a cumulative (i.e. plural or mass) argument. The cumulative \textit{na}- has been argued to require the internal argument not be a fully projected DP, but a smaller nominal constituent, QP \citep{Pereltsvaig2006}.}

\REF{tat:ex:24} represents a typical case where a lexically prefixed unergative verb acquires the internal argument position, \REF{ex:tatevosov:24b}, but with a superlexical prefix the same stem stays intransitive, \REF{ex:tatevosov:24c}:

\ea \label{tat:ex:24}     
\ea 
\gll rabota-t’  \minsp{(*} den’gi)\\
work-\textsc{inf}  {} money.\textsc{acc}\\\hfill \textsc{simplex stem}
\glt  ‘work’ \label{ex:tatevosov:24a}
\ex 
\gll za-rabota-t’  den’gi\\
LP-work-\textsc{inf}  money.\textsc{acc}\\\hfill \textsc{lp}
\glt  ‘earn money’\label{ex:tatevosov:24b}
\ex 
\gll   za-rabota-t’    \minsp{(*} den’gi)  \\
    SLP-work-\textsc{inf}  {} money.\textsc{acc}\\\hfill \textsc{slp}
\glt  ‘start working’\label{ex:tatevosov:24c}
\z
\z 

\noindent As an illustration of how LPs, unlike SLPs, manipulate thematic properties of arguments and project arguments not subcategorized for by a stem, consider \REF{tat:ex:25}--\REF{tat:ex:27}. 

\ea     \textsc{Non-prefixed}\label{tat:ex:25}

\gll  rubi-t’   \minsp{\{} derevj-a      /   drov-a   /   \minsp{\textsuperscript{??}} plennogo\}\\
  chop-\textsc{inf}  {} trees-\textsc{acc}  {}  firewood-\textsc{acc}  {} {}  captive.\textsc{acc}\\
\glt    (Intended:) ‘chop \{trees / firewood / a captive\}’
\z

\noindent The verb \textit{rubi-t’} ‘chop’ subcategorizes for an affected or effected object (‘trees’ and ‘firewood’ in \REF{tat:ex:25}) and cannot be interpreted as a verb of killing (as in ‘hack a captive’). 

Two effects of lexical prefixation are shown in \REF{ex:tatevosov:26a}--\REF{ex:tatevosov:26b}.

\newpage
\ea   \textsc{LPs}\label{tat:ex:26}
\ea
\gll s-rubi-t’   \minsp{\{} derevj-a   /   \minsp{*} drov-a   /   \minsp{*} plenn-ogo\}\\
\textsc{lp}-chop-\textsc{inf}  {} trees-\textsc{acc}  {} {}  firewood-\textsc{acc}  {} {}    captive-\textsc{acc}\\
\glt     (Intended:) ‘chop down \{the trees / the firewood / the captive\}’\label{ex:tatevosov:26a}
\ex 
\gll   za-rubi-t’   \minsp{\{*} derevj-a   /   \minsp{*} drov-a   /   plenn-ogo\}\\
\textsc{lp}-chop-\textsc{inf}  {}   trees-\textsc{acc} {} {}     firewood-\textsc{acc}  {}  captive-\textsc{acc}\\
\glt      (Intended:) ‘slash \{the trees / the firewood / the captive\}’\label{ex:tatevosov:26b}
\z
\z

\noindent In \REF{ex:tatevosov:26a}, the prefix narrows down the range of options available in \REF{tat:ex:25}: unlike \textit{rubi-t’}, \textit{s-rubi-t’} selects for an affected object ‘trees’, but not for an effected object ‘firewood’. Just like \textit{rubi-t’}, \textit{s-rubi-t’} is not a verb of killing (as in ‘hack a captive’).  In \REF{ex:tatevosov:26b}, prefixation of the lexical \textit{za}- forces the derived verb to take an unselected object ‘the captive’, while both types of arguments the stem \textit{rubi}- subcategorizes for become unavailable. I will address the argument structure of lexically prefixed verbs in more detail in \sectref{tat:sec:3}.

As before, superlexical prefixation has no consequences for the thematic properties of the internal argument projected by the stem. \REF{tat:ex:27} with the cumulative \textit{na}- exhibits the same range of options as \REF{tat:ex:25}.

\ea   \textsc{SLP}\\\label{tat:ex:27}
\gll na-rubi-t’   \minsp{\{} derevj-ev   /    drov   /  \minsp{*} plenn-yx\}\\
  \textsc{slp}-chop-\textsc{inf} {} trees-\textsc{gen.pl}  {}  firewood.\textsc{gen}   {} {} captive-\textsc{gen.pl}\\
\glt (Intended:) ‘chop (down) a quantity of \{trees / firewood / captives\}’
\z

\noindent Finally, \REF{tat:ex:28} shows that lexical prefixation effectively blocks the indefinite object alternation typically available for simplex imperfective verbs.

\ea    \textsc{Indefinite object alternation}\label{tat:ex:28}
\ea
\gll Dima   čita-l   \minsp{(} knig-u).\\
    Dima  read-\textsc{pst.m}  {} book-\textsc{acc}\\
\glt  ‘Dima was reading (a book).’\label{ex:tatevosov:28a}
\ex
\gll    Dima    pro-čita-l   \minsp{*(} knig-u).\\
    Dima  \textsc{lp}-read-\textsc{pst.m}  {} book-\textsc{acc}\\
\glt  ‘Dima read a book.’\label{ex:tatevosov:28b}
\ex
\gll    Dima    po-čita-l     \minsp{(} knig-u).\\
    Dima  \textsc{slp}-read-\textsc{pst.m}  {}  book-\textsc{acc}\\
\glt  ‘Dima spent some time reading (a book).’\label{ex:tatevosov:28c}
\z
\z

\noindent The simplex imperfective verb \textit{čita-t’} ‘read’ in \REF{ex:tatevosov:28a} can occur in a clause with no syntactically projected direct object, which describes the agent’s activity (‘Dima was involved in reading’).\footnote{Verbs like \textit{čita-t’} ‘read’ that allow for indefinite object alternation are all manner verbs in terms of \citet{Rappaport.Hovav.Levin1998} and elsewhere. Overall, Russian seems to conform to Rappaport Hovav \& Levin’s observation that being a manner verb is a necessary, but not a sufficient condition for the alternation.}  This option is preserved in \REF{ex:tatevosov:28c} with the delimitative superlexical \textit{po}-, but not in \REF{ex:tatevosov:28b} with the lexical \textit{pro}- (see, among others, \citealt{BabkoMalaya1999,DimitrovaVulchanova1996}).\footnote{The interaction between superlexical prefixation and indefinite object alternation requires further study. There certainly are SLPs that pattern with the delimitative \textit{po}- in \REF{tat:ex:28} in not affecting the alternation. This is true of the inceptive \textit{za}- (\textit{za-čita-t’} ‘start reading’), perdurative \textit{pro}- (\textit{pro-čita-t’ ves’ den’} ‘spend the entire day reading’), terminative \textit{ot}- (\textit{ot-čita-t’} ‘finish reading’), and, with some reservations, \textit{pere}- of excess (\textit{pere-čita-t’} ‘read too much’), even though not all speakers find these particular prefixed verbs natural. The judgments about the completive \textit{do}- are less clear, since for some completive verbs the alternation seems to be available easier than for others. %Thus, while \textit{do-pe-t’} ‘finish singing’ is perfectly fine without the direct object, many speakers find \textit{do-čita-t’} somewhat degraded in this environment. For the distributive \textit{po}- and \textit{pere}- the object is obligatory, but one can speculate that this can be due to the distributivity itself, since it may require there be individuals to be distributed over events. I am grateful to a reviewer for inviting me to discuss SLPs other than the delimitative \textit{po}- in this context.
} If \REF{ex:tatevosov:28b} without the object NP can be interpreted at all, it either involves object ellipsis (\textit{Žora pro-čita-l knigu, i Dima pro-čita-l \sout{knigu}.}) or a null pronoun in the object position (— `Who read Das Kapital$_i$?' — \textit{Dima pro-čita-l \textit{pro\textsubscript{i}}.}).\footnote{Care should be taken when applying other alleged diagnostics of the lexical or superlexical status of a prefix mentioned in the literature, e.g. their ability to undergo secondary imperfectivization (\citealt{Ramchand2004,Gribanova2013}), form event nominals (\citealt{BabkoMalaya1999,Svenonius2004}) and stative/resultative participles or head infinitival subjects \citep[165]{Gehrke2008}.  Upon closer scrutiny, one can observe that none of these diagnostics works quite the way it is supposed to work; for instance, there is massive evidence that only a proper subset of SLP configurations is incompatible with secondary imperfectivization. See \citet{Tatevosov2013} for an extensive discussion of the relevant constraints on secondary imperfectivization and nominalization.}

Since the early work on the lexical-superlexical distinction, there has been no general agreement as to how this phenomenon should be accounted for. For a theory of prefixation, the entire set of questions about mechanisms of and constraints on word formation needs to be settled, including the central question of whether it happens in the syntax or in a separate module of the grammar. 

Putting aside Slavic academic traditions (e.g. \citealt{Svedova1980, Grzegorczykowa.Laskowski.Wrobel1998, Sticha2018}, a.o.), the majority of existing literature on the LP/SLP divide assumes a syntactic, anti-lexicalist view of prefixation, and it is the view that I will focus on in what follows. One of the relatively rare examples of the opposite approach is \citet{Spencer.Zaretskaya1998} who distinguish between lexical prefixes and “Aktionsart prefixes” (${\approx}$ SLPs in a different terminology) and assume that at least the former occur at the level of lexical conceptual structure.\footnote{\citet{Spencer.Zaretskaya1998} argue that a prefix is a “core predicate” in the sense that it determines the event structure of a verb, whereas the verb stem should be treated as a modifier to this structure.} 

Simplifying somewhat, one can identify two parameters of variation within anti-lexicalist theories. One has to do with projection properties shown in \tabref{tab:1}, which includes minimal references to the work where a particular theoretical position  is assumed or argued for.

\begin{table}
    \centering
    \begin{tabularx}{.85\textwidth}{lll}
    \lsptoprule
    minimal projections & heads & \citet{Svenonius2004} (LPs) \\
    && \citet{Slabakova2005} (LPs)\\
    && \citet{Zaucer2009}\\
    && \citet{Biskup2019}\\\cmidrule(lr){2-3}
    & head-adjuncts & \citet{BabkoMalaya1999}\\\cmidrule(lr){1-3}
    maximal projections & specifiers &\citet{Svenonius2004} (SLPs)\\
    &&\citet{Svenonius2008} (all)\\\cmidrule(lr){2-3}
    & adjuncts & \citet{Slabakova2005} (SLPs)\\\cmidrule(lr){1-3}
    concord morphology & &  \citet{Arsenijevic2012}
    \\\lspbottomrule
    \end{tabularx}
    \caption{The morphosyntax of Slavic prefixes}
    \label{tab:1}
\end{table}

\iffalse 
\begin{table}
    \centering
    \begin{tabularx}{.9\textwidth}{lllll}
    \lsptoprule
    \multicolumn{5}{c}{\textsc{Prefixes are}}\\
    \multicolumn{2}{c}{minimal projections}&\multicolumn{2}{c}{maximal projections}&concord morphology\\\cmidrule(lr){1-5}\\Heads & Head-adjuncts & Specifiers & Adjuncts & \\\citet{Svenonius2004} (LPs)  
    & \citet{BabkoMalaya1999} & \citet{Svenonius2004} (SLPs) & \citet{Slabakova2005} (SLPs) & \citet{Arsenijevic2012}\\
    \citet{Slabakova2005} (LPs) & & \citet{Svenonius2008} (all) & &\\
    \citet{Zaucer2009} & & & &\\
    \citet{Biskup2019} & & & & \\\lspbottomrule
    \end{tabularx}
    \caption{}
    \label{}
\end{table}

\begin{table}
\caption{Projection properties.}
\label{tab:key:1}

\begin{tabularx}{\textwidth}{XXXXX}

\lsptoprule

\multicolumn{2}{c}{Prefixes are minimal projections}  & \multicolumn{2}{c}{Prefixes are maximal projections} & Prefixes are neither minimal nor maximal projections; they are instances of concord morphology\\
Prefixes are heads & Prefixes are head-adjuncts & Prefixes are specifiers & Prefixes are adjuncts & \\
\citealt{Slabakova2005} for LPs (her “internal prefixes”), \citealt{Svenonius2004} for LPs, \citealt{Žaucer2009,Biskup2019} & \citealt{Babko-Malaya1999} & \citealt{Svenonius2004} for SLPs; \citealt{Svenonius2008} for all prefixes & \citealt{Slabakova2005} for SLPs (her “external prefixes”) & \citealt{Aresenijevic2012}\\
\lspbottomrule
\end{tabularx}
\end{table}
\fi 

The dominant view of prefixation is that prefixes are heads that combine with verb stems via head movement or a similar mechanism. One major departure from this view, advanced in the work by Peter Svenonius, is that either some or all prefixes are maximal projections. \citet{Svenonius2004, Svenonius2008} discusses a number of arguments supporting this position, including compelling evidence that prefix $+$ stem combinations exhibit phonological effects characteristic of word boundaries rather than morpheme boundaries. A different alternative is offered by \citet[][and elsewhere]{Arsenijevic2012}, who argues that prefixes are not syntactically represented at all and are to be analyzed as phonological signatures of the agreement relations between V, Asp and a preposition. 

The other, and more significant, major parameter of theoretical variation is established by existing assumptions about the location of prefixes in the hierarchical structure. \tabref{tab:2} represents a number of options attested in the literature.

\begin{table}
    \centering
    \begin{tabularx}{.9\textwidth}{lll}
    \lsptoprule
    LPs & within complement of V & \citet{Svenonius2004}  \\
    &&\citet{Asbury+2007}\\
    &&\citet{Romanova2007}\\
    &&\citet{Gehrke2008}\\
    &&\citet{Zaucer2009}\\\cmidrule(lr){2-3}
    & adjoined to V &\citet{BabkoMalaya1999} \\\cmidrule(lr){2-3}
    & within FP outside VP &\citet{Slabakova2005} \\
    &&\citet{Wiland2012}\\\cmidrule(lr){1-3}
    SLPs& within FP outside highest& \citet{BabkoMalaya1999}\\
    &~~~projection of LP/VP&\citet{Svenonius2004}\\
    &&\citet{Slabakova2005}\\
    &&\citet{Gehrke2008}\\
    &&\citet{Gribanova2013,Gribanova2015}\\\cmidrule(lr){2-3}
    & within complement of V (like LPs) &\citet{Zaucer2009,Zaucer2010,Zaucer2012}
    \\\lspbottomrule
    \end{tabularx}
    \caption{The origin of Slavic prefixes}
    \label{tab:2}
\end{table}


\iffalse
\begin{table}
\caption{Location}
\label{tab:key:2}

\begin{tabularx}{\textwidth}{XXXXX}
\lsptoprule
\multicolumn{3}{c}{ LPs} & \multicolumn{2}{c}{ SLPs}\\
LPs originate within the complement of V & LPs are adjoined to V & LPs originate within a functional projection outside of VP & SLPs originate within functional projections outside of the highest projection of a lexical prefix or VP, whichever is higher & SLPs are resultative, i.e., like LPs originate within the complement of V\\
\citealt{Svenonius2004,Romanova2006}, Asbury \textit{et al.} 2006, \citealt{Gehrke2008,Žaucer2009} & \citealt{Babko-Malaya1999} & \citealt{Slabakova2005,Wiland2012} & \citealt{Babko-Malaya1999,Svenonius2004,Slabakova2005,Gehrke2008,Gribanova2013,Gribanova2015} & \citealt{Žaucer2009,Žaucer2010}, 2012\\
\lspbottomrule
\end{tabularx},
\end{table}
\fi

Most proposals rely on the structural asymmetry between LPs and SLPs, and many on a particular type of asymmetry in \REF{tat:ex:29}, where LPs merge within the complement of V, while SLPs appear outside of VP.

\ea{}   [ … SLP … [\textsubscript{VP} … [ … LP … ] ] ]\label{tat:ex:29}
\z

\noindent \citet{Svenonius2004} argues, specifically, that LPs originate in the R(esult) projection, whereas SLPs adjoin to a functional projection above VP. This explains a bunch of properties that separate the two classes: LPs are correctly predicted to occur inside of SLPs, to have impact on the argument structure, to develop idiomatic meanings easier, and so on. 

However, virtually all logically possible alternatives to \REF{tat:ex:29} can be found in the literature. Some authors assign a different position to LPs; cf. \citeposst{BabkoMalaya1999} head-adjunction and \citeposst{Slabakova2005} analysis in \tabref{tab:2}. Others argue for a different location of SLPs. \citet{Zaucer2009,Zaucer2010,Zaucer2012}, for example, develops a set of carefully designed arguments suggesting that SLPs, just like LPs, originate under V. One of the arguments is based on the observation that at least some superlexicals introduce a result state into the semantic representation of a predicate. If result state descriptions are only generated below V, Žaucer’s conjecture follows. A theory along the lines of \REF{tat:ex:29} needs to assume that SLPs can introduce result states from outside of VP. This obviously disrupts a correspondence between the layers of structure and types of eventuality descriptions that form their denotations, which is taken to be a desideratum in a number of theories, e.g. \citeposst{Ramchand2008} First Phase Syntax, where all result states come out as denotations of R(esult)P/resP. 

A separate question is how SLPs, if they are VP-external, are located with respect to \textit{v} and Asp. Many authors assume, much in line with traditional Slavic linguistics, that “aspectual morphology” contributes directly to the computation of semantic aspect and originate within AspP, i.e. outside of \textit{v}P. \citet{Tatevosov2011, Tatevosov2015b,Tatevosov2021}, however, builds a chain of arguments that pieces of aspectual morphology are \textit{v}P-internal and modify event structure, while true semantic aspectual operators in Asp are phonologically silent. Proposals that treat aspectual morphology as concord morphology (e.g. \citealt{Arsenijevic2012} and subsequent work) assume that prefixes are lower than Asp, too. 

Assessing far-reaching predictions of these proposals would take us too far from the central topic of this chapter, however. Having in mind the alternatives in \tabref{tab:2}, in what follows I will be assuming the architecture of the verbal domain along the lines of \REF{tat:ex:29}. 

In the next section, event-structural correlates of lexical prefixation and secondary imperfectivization are surveyed.

\section{Lexical prefixation, secondary imperfectivization, and event structure}\label{tat:sec:2}
\hypertarget{RefHeadingToc164647621}{}
The goal of this section is to establish a number of generalizations about the event structure of morphologically \textsc{simplex} \textsc{imperfective verbs}, \textsc{verbs derived by lexical prefixation}, and \textsc{secondary imperfective verbs}, shown in \tabref{tab:key:3}.

\begin{table}
\caption{Simplex Imperfective, Perfective, and Secondary Imperfective verbs}
\label{tab:key:3}
\begin{tabularx}{\textwidth}{XXXX}
\lsptoprule
Verb  & Simplex & Perfective & Secondary \\
meaning & Imperfective & & Imperfective\\\midrule
‘write’ & \textit{pisa-t’} \newline write-\textsc{inf} & \textit{na-pisa-t’} \newline \textsc{lp}-write-\textsc{inf}\smallskip & --\\
‘sign’ & -- & \textit{pod-pisa-t’} \newline \textsc{lp}-write-\textsc{inf} & \textit{pod-pis-yva-t’} \newline \textsc{lp}-write-\textsc{yva}-\textsc{inf}\smallskip\\
‘read’ & \textit{čita-t’} \newline read-\textsc{inf} & \textit{pro-čita-t’} \newline \textsc{lp}-read-\textsc{inf} & \textit{pro-čit-yva-t’} \newline \textsc{lp}-read-\textsc{yva}-\textsc{inf}\\
\lspbottomrule
\end{tabularx}
\end{table}


\iffalse
\begin{table}
\caption{Simplex imperfective, perfective, and Secondary imperfective verbs}
\label{tab:key:3}
\begin{tabularx}{\textwidth}{XXXX}
\lsptoprule
Verb  & Simplex & Perfective & Secondary \\
meaning & imperfective & & imperfective\\\hline
‘write’ & \textit{pisa-t’} ‘write\textsuperscript{IPFV}’ \newline write-INF & \textit{na-pisa-t’} ‘write\textsuperscript{PFV}’ \newline LP-write-INF & —\\
‘sign’ & — & \textit{pod-pisa-t’} ‘sign\textsuperscript{PFV}’ \newline LP-write-INF & \textit{pod-pis-yva-t’} ‘sign\textsuperscript{IPFV}’  \newline LP-write\textsc{-}YVA\textsc{-}INF\\
‘read’ & \textit{čita-t’}  ‘read\textsuperscript{IPFV}’ \newline read-INF & \textit{pro-čita-t’} ‘read\textsuperscript{PFV}’  \newline LP-read -INF & \textit{pro-čit-yva-t’} ‘read\textsuperscript{IPFV}’ \newline LP-read\textsc{-YVA-}INF\\
\lspbottomrule
\end{tabularx}
\end{table}
\fi

\tabref{tab:key:3} illustrates three most basic types of relationship between simplex imperfective, lexically prefixed perfective and secondary imperfective verbs. \textit{Na-pisa-t’} ‘write\textsuperscript{PFV}’ as opposed to \textit{pisa-t’} ‘write\textsuperscript{IPFV}’ is frequently referred to in the literature as an instance of “pure perfectivizing”, or “empty” prefixation: the only perceivable contribution of the prefix \textit{na}- in \textit{na-pisa-t’} ‘write\textsuperscript{PFV}’, as \citet[51]{BabkoMalaya1999} puts it, “is to indicate that the process denoted by the verb is completed”. The word “empty” here reflects the intuition that prefixation does not produce any change in the lexical meaning of a verb. \textit{Pisa-t’} ‘write\textsuperscript{IPFV}’ and \textit{na-pisa-t’} ‘write\textsuperscript{PFV}’ are thus considered aspectual partners that form \textsc{aspectual pair}. %The same relationship obtains between \textit{čita-t’} ‘read\textsuperscript{IPFV}’ and \textit{pro-čita-t’} ‘read\textsuperscript{IPFV}’. 

\textit{Pod-pisa-t’} ‘sign\textsuperscript{PFV}’ illustrates “resultative prefixation” where there is more to the meaning of the prefixed verb than “perfectivization”. Intuitively, \textit{pod}- derives a different lexical item from \textit{pisa-t’} ‘write\textsuperscript{IPFV}’, \textit{pod-pisa-t’} ‘sign\textsuperscript{PFV}’. An imperfective partner of \textit{pod-pisa-t’} ‘sign\textsuperscript{PFV}’ is \textit{pod-pis}\textsc{-}\textit{yva}\textsc{-}\textit{t’} ‘sign\textsuperscript{IPFV}’, the product of secondary imperfecitivization.\footnote{A terminological note is due. The notion of “resultative” prefixation (e.g. \textit{pod-pisa-t’}) as opposed to “pure perfectivizing” prefixation (e.g. \textit{na-pisa-t’} ‘write\textsuperscript{PFV}’) goes back to \citet{BabkoMalaya1999}: 51-56. To avoid confusion, one should keep in mind that there is a different tradition (e.g. \citealt{Isacenko1960}) where the term “resultative” has a broader meaning and covers both \textit{na-pisa-t’} ‘write\textsuperscript{PFV}’ and \textit{pod-pisa-t’} ‘sign\textsuperscript{PFV}’. Within this tradition, there is no designated term reserved for verbs like \textit{pod-pisa-t’} ‘sign\textsuperscript{PFV}’.}  

In \sectref{tat:sec:3} I will discuss a more refined typology of lexical prefixes, not restricted to “empty” (\textit{na}- in \textit{na-pisa-t’}; \textit{pro}- in \textit{pro-čita-t’}) and “resultative” ones (\textit{pod}- in \textit{pod}-\textit{pisa}-\textit{t’}). For the purposes of this section these two will suffice, however. 

Whereas most perfective verbs derived by “empty” prefixes, \textit{na-pisa-t’} ‘write\textsuperscript{PFV}’ in \tabref{tab:key:3} among them, do not undergo secondary imperfectivization, some do. This happens to \textit{pro-čita-t’} ‘read\textsuperscript{PFV}’, which has both simplex and secondary imperfective counterparts, \textit{čita-t’} and \textit{pro-čit-yva-t’} ‘read\textsuperscript{IPFV}’. Derivationally related items like \textit{čita-t’} ‘read\textsuperscript{IPFV}’ /  \textit{pro-čita-t’} ‘read\textsuperscript{PFV}’ / \textit{pro-čit-yva-t’} ‘read\textsuperscript{IPFV}’   are frequently referred to in the literature as \textsc{aspectual triples}. In general, it is not predictable if a simplex verb forms a pair like \textit{pisa-t’} ‘write\textsuperscript{IPFV}’ or a triple like \textit{čita-t’} ‘read\textsuperscript{IPFV}’.\footnote{Aspectual triples are not available in every Slavic language. \citet{Klimek-Jankowska.etal2025} indicate that triples are systematically absent in West Slavic languages, as opposed to East Slavic languages, and put forward the hypothesis that this cross-linguistic difference can be reduced to the range of syntactic positions available for the secondary imperfective morpheme.}  

In what follows we will see that event-structurally, lexically prefixed perfective verbs, “pure perfectivizing” and “resultative” verbs alike, differ systematically from simplex, but not from secondary imperfective verbs. 

\subsection{Simplex imperfective and prefixed perfective verbs}\label{sec:2.1}
\hypertarget{RefHeadingToc164647622}{}
The first generalization about the event structure projected by the verbs like those in \tabref{tab:key:3} comes in \REF{tat:ex:30}.

\eanoraggedright    \textsc{Event structure of simplex imperfective and lexically prefixed perfective verbs}\label{tat:ex:30}

  The event structure of simplex imperfective verbs lacks a result state. Lexically prefixed verbs are associated with a more complex event structure containing a result state. 
\z 

\noindent According to \REF{tat:ex:30}, simplex imperfective and lexically prefixed perfective verbs manifest different degrees of subevental complexity, the latter having one subevent more than the former. A natural way of making \REF{tat:ex:30} explicit is to translate it into event-semantic representations of the format in \REF{tat:ex:31}--\REF{tat:ex:32}. 

\eanoraggedright   \textsc{VPs projected by simplex imperfective verbs}: A predicate of events, simplex event structure\label{tat:ex:31}

$\lambda e.\, ...\, e\, ...$
\z

\eanoraggedright   \textsc{VPs projected by prefixed perfective verbs}: A relation between events and states, complex event structure\label{tat:ex:32}

$\lambda s. \lambda e.\, ...\, e\, ... \wedge ...\, s\, ... \wedge R(s)(e)$
\z

\noindent \REF{tat:ex:31} is a predicate of eventualities, whereas \REF{tat:ex:32} is a relation between events and (result) states. The relation R in \REF{tat:ex:32} is typically conceived of as a causal relation. \REF{tat:ex:31}--\REF{tat:ex:32} can be thus thought of as approximating what has been labeled activities and accomplishments in \citeauthor{Dowty1979}'s (\citeyear{Dowty1979} and much further work) line of inquiry, with its decompositional treatment of accomplishments.\footnote{\REF{tat:ex:31} and \REF{tat:ex:32} correspond to \citeposst{Dowty1979} activity and accomplishment templates, $\cnst{do}(\alpha,\, \pi (\alpha,...))$ and $\cnst{do}(\alpha,\, \pi(\alpha,...))\, \cnst{cause}\, (\cnst{become} \rho (\beta,...))$, respectively. The major differences are the following: \citeauthor{Dowty1979} takes the external argument to be part of the representation, employs temporal semantics rather than event semantics, makes use of the \cnst{become} operator, and assigns more structure to the process component.}  Since then, in the literature  configurations like \REF{tat:ex:31} or similar have been sometimes referred to as an activity event structure, while those like \REF{tat:ex:32} as an accomplishment event structure. I will refrain from using this terminology in what follows to avoid confusion with accomplishments and achievements as eventuality types. Rather, \REF{tat:ex:31}--\REF{tat:ex:32} will be merely identified as a simplex and a complex event structure, respectively.

The uninflected tenseless and aspectless VPs from \REF{ex:tatevosov:33a} and \REF{ex:tatevosov:34a} would have the denotations in \REF{ex:tatevosov:33b} and \REF{ex:tatevosov:34b}.\footnote{The representations in \REF{ex:tatevosov:33b} and \REF{ex:tatevosov:34b} assume that external arguments are introduced by a functional head outside of VP and rely on the Davidsonian association of internal arguments with events (\citealt{Kratzer1996,Kratzer2003} and much further work). As a simplification, nominal arguments are represented in \REF{ex:tatevosov:33b} and \REF{ex:tatevosov:34b} as individual constants.}

\eanoraggedright   \textsc{Simplex imperfective verbs}: A predicate of events; simplex event structure\label{ex:tatevosov:33}

 \ea
 \gll Volodja  pisa-l  pis’m-o.\\
     Volodja  write-\textsc{pst.m}  letter-\textsc{acc}\\
\glt ‘Volodja was writing a letter.’\label{ex:tatevosov:33a}
 \ex
   \sib{[\textsubscript{VP} \textit{pisa}- \textit{pis’mo}]} = $\lambda e. \textsc{write}(\textsc{letter})(e)$\label{ex:tatevosov:33b}
 \z
\z

\eanoraggedright   \textsc{Lexically prefixed verbs}: A relation between events and states; complex event structure\label{ex:tatevosov:34}
  \ea
  \gll Volodja  na-pisa-l  pis’m-o.\\
     Volodja  \textsc{lp}-write-\textsc{pst.m}  letter-\textsc{acc}\\
\glt ‘Volodja wrote a letter.’\label{ex:tatevosov:34a}
  \ex     \sib{[\textsubscript{VP} \textit{na-pisa- pis’mo}]} = $\lambda s. \lambda e. \textsc{write}(\textsc{letter})(e) \wedge \cnst{cause}(s)(e) \wedge \textsc{written}(\textsc{letter})(s)$\label{ex:tatevosov:34b}
  \z
  \z

\noindent Essentially, \REF{ex:tatevosov:33b} represents the simplex event structure in \REF{tat:ex:31}, whereas \REF{ex:tatevosov:34b} has the complex structure in \REF{tat:ex:32}.\footnote{The eventive element of \REF{ex:tatevosov:34b} can be further decomposed into the activity and change-of-state, yielding a three-component decomposition, as in \citet{paduceva96},  \citet{Ramchand2008} or \citet{BeaversKoontz-Garboden2020}. As far as I can see, nothing in what follows depends on this possible refinement.}  Much in the spirit of \citet{Klein1995}, \textit{pisa-} ‘write’ and other non-stative simplex imperfectives project VPs that denote predicates of eventualities (of type $\langle v, t \rangle$) (cf. \citeauthor{Klein1995}’s 1-state descriptions). The predicate in \REF{ex:tatevosov:33b} has writing events in its extension, whose theme is a letter.

The representation of \REF{ex:tatevosov:34b} contains, in addition, a target state of being written brought about by a writing event. Prefixed stems like \textit{na-pisa}- thus denote relations between events and states (of type $\langle v, \langle v, t \rangle\rangle$; \citeauthor{Klein1995}’s 2-state descriptions). The eventive part of \REF{ex:tatevosov:34b} corresponds to a “source state” in \citeposst{Klein1995} terminology; the state argument ranges over “target states”. A holder of that state is identical to the theme of writing. The target state component is represented as “\textsc{written}(\textsc{letter})($s$)”, where “\textsc{written}” is of type $\langle e, \langle v, t \rangle\rangle$, a relation between individuals and states. An analysis of the simplex perfective verbs like \textit{da-t’} ‘give\textsuperscript{PFV}, a verb of transfer of possession, would be similar, possibly with minor adjustments.

As we have already seen in the previous section, in the extensive work on prefixation in Slavic languages (\citealt{BabkoMalaya1999,Ramchand2004,Romanova2004,Svenonius2004,Svenonius2008,Arsenijevic2007a,Arsenijevic2007b,Gehrke2008,Tatevosov2009,Tatevosov2013,Zaucer2009,Zaucer2010,Zaucer2013,Biskup2015,Biskup2019}), it has been independently argued that the role of prefixes like \textit{na}- in \REF{ex:tatevosov:34a} is to project a structure interpreted as a target state description. \REF{ex:tatevosov:34b}, which differs from \REF{ex:tatevosov:33b} in the presence of a lexical prefix, makes this fully explicit.

There are a number of empirical arguments for \REF{ex:tatevosov:33b} and \REF{ex:tatevosov:34b} based on converging evidence for higher event-structural complexity of lexically prefixed verbs. The body of literature on predicate decomposition starting from \citet{Dowty1979} (see \citealt{Stechow1996,Rappaport.Hovav.Levin1998,Rapp.vonStechow1999,Kratzer2000,Rothstein2004,Ramchand2008}, among many others) offers a number of diagnostics that allows us to tell the two configurations apart. The list includes, specifically, \textsc{argument realization patterns}, \textsc{scope of adverbials} like `almost' and `again', \textsc{interpretation under negation}.

The first piece of evidence comes from argument realization. As \citet{Rappaport.Hovav.Levin1998} observe, verbs that show the indefinite object alternation (e.g. \textit{sweep}) are all associated with the simplex event structure. If a verb has the structure of accomplishments (e.g., \textit{break}), the alternation is unavailable: 

\ea\label{tat:ex:35}
\ea    John swept the floor.\label{ex:tatevosov:35a}
\ex    John swept.\label{ex:tatevosov:35b}
\z
\z

\ea\label{tat:ex:36}
\ea    John broke the window.\label{ex:tatevosov:36a}
\ex[\#]{John broke.}\label{ex:tatevosov:36b}
\z
\z

\noindent As we have already seen in the previous section, simplex imperfective and lexically prefixed perfective verbs in Slavic differ in exactly the same way. Examples \REF{ex:tatevosov:28a}--\REF{ex:tatevosov:28b} are repeated as \REF{ex:tatevosov:37a}--\REF{ex:tatevosov:37b}.


\ea    \textsc{Indefinite object alternation}\label{tat:ex:37}
\ea
\gll Dima   čita-l   \minsp{(} knig-u).\\
    Dima  read-\textsc{pst.m}  {} book-\textsc{acc}\\
\glt  ‘Dima was reading (a book).’\label{ex:tatevosov:37a}
\ex
\gll    Dima    pro-čita-l   \minsp{*(} knig-u).\\
    Dima  \textsc{lp}-read-\textsc{pst.m}  {} book-\textsc{acc}\\
\glt  ‘Dima read a book.’\label{ex:tatevosov:37b}
\z
\z

\noindent The prefixed verb in \REF{ex:tatevosov:37b} is ungrammatical without the direct object, while its non-prefixed counterpart in \REF{ex:tatevosov:37a} is readily available. This is exactly what one would expect, given \citeposst{Rappaport.Hovav.Levin1998} generalization, if prefixed verbs introduce accomplishment event structures whereas non-prefixed verbs have the structure of activities.

\citet{Rappaport.Hovav.Levin1998} propose that the ungrammaticality of \REF{ex:tatevosov:36b} is due to a principle, operational in the grammar of natural languages, that requires syntactic realization of the structural aspects of the meaning of lexical items. For \textit{break}, associated with a complex event structure, but not for \textit{sweep}, the internal argument is structural. It is construed as the holder of a result state of being broken and as such must be syntactically projected. The internal argument of \textit{sweep} is only licensed as an  argument of the lexical constant of the verb, which makes its syntactic realization optional.

The Slavic pattern in \REF{ex:tatevosov:37a}--\REF{ex:tatevosov:37b}
is addressed in detail in \citet[64--68]{BabkoMalaya1999}; \citet{DimitrovaVulchanova1996,Biskup2019}. Despite being very different as to the specifics of how the structure of prefixed configurations is organized, they seem to fundamentally agree that obligatoriness of the syntactic realization of the object is induced by the prefix, which introduces a result state description.

Another set of diagnostics is based on a further prediction derivable from an analysis along the lines of \REF{tat:ex:31}--\REF{tat:ex:32}. If an event is decomposed into subevents, one should expect to find semantic operations that can target some of the subevents, but not the others. Applying to subeventally more complex structures like \REF{ex:tatevosov:34b}, such operations should yield a different range of interpretations compared to simplex structures like \REF{ex:tatevosov:33b}.

One piece of evidence for \REF{ex:tatevosov:34b} comes from the distribution of adverbials like ‘again’ \citep{Dowty1979,Stechow1996,Beck2005,Lechner.etal2015} and ‘almost’ \citep{Rapp.vonStechow1999}. With the complex event structure, both are known to be able to manifest varying scopal possibilities illustrated informally in \REF{ex:tatevosov:38a}--\REF{ex:tatevosov:38b}. On the low reading in \REF{ex:tatevosov:38a}, the scope of ‘almost’ and ‘again’ only includes the result state. On the “high” reading in \REF{ex:tatevosov:38b}, it extends to the entire eventuality description including the process component.

\ea\label{ex:tatevosov:38}
\ea \textsc{process}  ‘almost’ / ‘again’ [\textsc{result} \textsc{state}]\label{ex:tatevosov:38a}
  \ex   ‘almost’ / ‘again’ [\textsc{process}   \textsc{result} \textsc{state}]\label{ex:tatevosov:38b}
  \z
\z

\noindent Combining with a simplex event structure that lacks a result state altogether, ‘almost’ and ‘again’ are expected not to show this type of ambiguity. In \REF{tat:ex:39}, their position is outside of the only component of the event structure, the process component. The only available reading of ‘almost’ and ‘again’ should thus correspond to the “high” reading in \REF{ex:tatevosov:38b}.

\ea    ‘almost’ / ‘again’ [\textsc{process}]\label{tat:ex:39}
\z

\noindent Consider ‘almost’. Intuitively, ‘almost’ indicates that the scope proposition is false in the evaluation world, but a close-by scalar alternative is true (e.g., \citealt{Penka2006}). Compare the discourses in \REF{ex:tatevosov:40a} and \REF{ex:tatevosov:40b}.

\ea\label{tat:ex:40}
\eanoraggedright  \textit{Context}: — Volodja, how long will it take you to finish this simple letter? You’ve been working on it for an hour!\label{ex:tatevosov:40a}

\gll   —  Počti na-pisa-l!  Ešče  dve  minut-y.\\
    {} almost  \textsc{lp}-write-\textsc{pst.m}  more  two  minute-\textsc{gen.sg}\\
\glt  ‘I almost wrote it. Two more minutes.’
\ex     \textit{Context}: — We need to ask the dean for discretionary funds. — You know,

\gll    ja  počti na-pisa-l  pis’m-o  dekanu,  no pereduma-l.\\
    I  almost   \textsc{lp}-write-\textsc{pst.m}  letter-\textsc{acc}   dean-\textsc{dat}    but change.mind-\textsc{pst.m}\\
\glt  ‘I almost wrote a message to the dean (about that), but then changed my mind.’\label{ex:tatevosov:40b}
\z
\z

\noindent \REF{ex:tatevosov:40a} contextually entails that the the agent’s writing activity has occurred, but the culmination and a subsequent result state have not. \REF{ex:tatevosov:40a} thus instantiates \REF{ex:tatevosov:38a}: the  scope of ‘almost’ only includes the result state. In the null context, this interpretation of \textit{počti na-pisa-t’} ‘almost write’ is preferred.

Another interpretation is observed in \REF{ex:tatevosov:40b}: The agent comes close to performing an entire event, including its eventive component, but changes his mind. No part of a writing-a-message eventuality occurs in the evaluation world. The scope of ‘almost’ in \REF{ex:tatevosov:40b} includes both components of an eventuality, in line with \REF{ex:tatevosov:38b}.

These examples together show, therefore, that exactly as \REF{ex:tatevosov:38} predicts, ‘almost’ can take scope over the entire complex eventuality, \REF{ex:tatevosov:40b}, or over the result state component only, \REF{ex:tatevosov:40a}. This provides us with a piece of evidence for \REF{ex:tatevosov:34b}, a semantic representation containing an activity and a result state as distinct elements of event structure.

When it comes to the event structure of simplex verbs like \textit{pisa-t’} ‘write\textsuperscript{IPFV}’, a few precautions should be taken. As we have already seen, such verbs can only occur in imperfective clauses, which, on the progressive reading, come with the inference that the culmination is not included into the topic time. Consider \REF{tat:ex:41}.

\ea\label{tat:ex:41}
\gll Volodja  počti pisa-l  pism-o.\\
  Volodja  almost  write-\textsc{pst.m}  letter-\textsc{acc}\\
\glt ‘Volodja was almost writing a letter.’
\z

\noindent \REF{tat:ex:41} is unambiguous, indicating that the agent is close to writing a letter in the evaluation world but no writing activity occurs (yet). However, the lack of ambiguity can have more than one explanation. \REF{tat:ex:41} may be unambiguous because simplex imperfective verbs do not have a result state component, as \REF{ex:tatevosov:33b} suggests. But there can be another reason as well: \textit{pisa-t’} ‘write’ in \REF{tat:ex:41} does have a result state in its semantic representation, but this state is not reached in the evaluation world because of the progressive, which renders the low reading of \textit{počti pisa-l} ‘was almost writing’, parallel to that of \textit{počti na-pisa-l} in \REF{ex:tatevosov:40a}, unavailable. Progressive sentences like \REF{tat:ex:41} are therefore uninformative with respect to the event structure of simplex verbs.

Fortunately, there is a way of getting around this difficulty. The imperfective in Russian is available in a few environments where it can describe a culminating eventuality, thus being truth-conditionally equivalent to the perfective. (In the traditional literature on Slavic aspect this is known as \textsc{aspectual competition} between the perfective and the imperfective, see, e.g. \citealt{Maslov1984}; configurations, in which this happens are sometimes called “Maslov’s contexts”). In such an environment, if simplex imperfectives are as subeventally complex as lexically prefixed perfectives, they should exhibit the same range of interpretations as observed in \REF{ex:tatevosov:40a}--\REF{ex:tatevosov:40b}.

One configuration of this type is served by Historical Present (\citealt{Schlenker2004,Eckardt2015}, a.m.o.). Consider \REF{tat:ex:42}.

\eanoraggedright    \textit{Context}: Volodja comes home, sits down at his desk, takes a pen and paper, ...\label{tat:ex:42}

\gll piš-et  pis’m-o  i  lož-it-sja  spa-t’.  \\
    write-\textsc{npst.3sg}  letter-\textsc{acc}  and  lie.down-\textsc{npst.3sg-refl}  sleep-\textsc{inf}\\
\glt  ‘writes a letter, and goes to bed’.
\z

\noindent The narrative in \REF{tat:ex:42} describes a sequence of eventualities each of which, including ‘write a letter’, reaches a culmination; see \citeauthor{chapters/05_gehrke} (\citeyear{chapters/05_gehrke} [this volume]: section 3.3) for an overview of the Historical Present in Czech and Russian. This guarantees that the Historical Present provides a right environment for examining the scopal possibilities of ‘almost’ with a simplex verb. In \REF{tat:ex:43}, the same discourse is modified in order to accommodate ‘almost’. 

\eanoraggedright    \textit{Context}: Volodja comes home, sits down at his desk, takes a pen and paper, ...\label{tat:ex:43}

\gll  i  počti  piš-et  pis’m-o,  kogda razda-ёt-sja   zvonok  v  dver'.\\
  and  almost  write-\textsc{npst.3sg}  letter-\textsc{acc}  when give.out-\textsc{npst.3sg-refl}   ring  in  door.\textsc{acc}\\


\glt ‘and almost writes a letter, when the doorbell rings.’
\z

\noindent The only reading available to \textit{počti} ‘almost’ in \REF{tat:ex:43} is the “high” reading where no activity occurs in the evaluation world. But unlike with the progressive in \REF{tat:ex:41}, this cannot be attributed to the fact that ‘writes’ in \REF{tat:ex:43} cannot describe a culminating eventuality. As \REF{tat:ex:42} shows, it can. Therefore, the only reason for there being no “low” reading of ‘almost’ in \REF{tat:ex:43} is that \textit{pisa-t’} ‘write’ genuinely lacks a result state that \textit{počti} ‘almost’ can take scope over independently of the eventive component. Therefore, \REF{tat:ex:43} argues for \REF{tat:ex:39} and, in this way, for the semantic representation of simplex imperfective verbs along the lines of \REF{ex:tatevosov:33b}.

Another one of “Maslov’s contexts” in which the imperfective is considered aspectually parallel to the perfective is provided by habitual sentences. Consider \REF{tat:ex:44}. 

\eanoraggedright    \textit{Context}: In the morning, Volodja gets up early, sits down at his desk, reads newspapers, ...

\gll   piš-et  pis’m-o   komu-nibud’   iz   druz-ej   i id-et     zavtraka-t’.\\
  write-\textsc{npst.3sg}  letter-\textsc{acc}  who.\textsc{dat-indef}   from  friend-\textsc{gen.pl}   and go-\textsc{npst.3sg}  have.breakfast-\textsc{inf}\\
\glt ‘writes a letter to one of his friends and goes to breakfast.’\label{tat:ex:44}
\z

\noindent \REF{tat:ex:44} describes a regularity verified by multiple occurrences of episodic writing eventualities. Crucially, \REF{tat:ex:44} conveys that each of these eventualities culminates. Therefore, \textit{počti} \textit{pisa-t’} ‘almost write’ within the scope of the habitual provides us with the same environment as Historical Present. If \textit{pisa-t’} ‘write\textsuperscript{IPFV}’ is subeventally complex, \textit{počti pisa-t’} ‘almost write’ is expected to be ambiguous. But it is not: just as \REF{tat:ex:43}, the sentence in \REF{tat:ex:45} can only mean that no writing a letter activity occurs.

\eanoraggedright    \textit{Context}: Arriving in Geneva, Volodja unsuccessfully tries to establish regular correspondence with friends. Every morning after breakfast he sits down at his desk, takes a pen and paper, ...\label{tat:ex:45}

\gll  počti piš-et  pis’m-o   komu-nibud’   iz   druz-ej, no   objazatel’no   zasypa-et. \\
  almost   write-\textsc{npst.3sg}  letter-\textsc{acc}  who.\textsc{dat-indef}   from  friend-\textsc{gen.pl}  but   obligatorily   fall.asleep-\textsc{npst.3sg}\\
\glt ‘almost writes a letter to one of his friends, but always falls asleep.’
\z

\noindent Evidence from the Historical Present and habitual sentences therefore converges: no ambiguity with a decomposition adverb can be detected. Because of the strict parallelism of the two environments, in what follows I will only cite examples with Historical Present.

The same pattern can be brought out by \textit{opjat’} ‘again’. For the sake of space, I will refer the reader to \citet[493--497]{Tatevosov2016} and not reproduce corresponding examples here. 

Before moving on, let me briefly address one more potential complication with the range of interpretations available for decomposition adverbs. As Berit Gehrke (p.c.) points out, there is a class of verbs that are subeventally complex, but nevertheless do not allow for ‘almost’ and ‘again’ to show scope ambiguities. These are ditransitive verbs  discussed in detail in \citet{Bondarenko2018}:

\ea
\gll Maša  opjat’  otda-l-a  Vas-e  knig-u.\\
Maša  again  give-\textsc{pst-f}  Vasja-\textsc{dat}  book-\textsc{acc}\\
\glt `Maša gave Vasja the book, and she had given it to him before.’ \hfill (repetitive) \\
Unavailable: `Maša gave Vasja the book, and he had had it before.’\\\xspace\hfill (restitutive) \\
 \hfill (adapted from \citealt{Bondarenko2018}: 28)
\z 

\noindent \citeauthor{Bondarenko2018} argues extensively that this pattern cannot be derived from any special properties of \textit{opjat’} ‘again’ in Russian. She finds independent arguments that the semantic representation of ditransitives does contain a result state. She concludes, therefore, that the reason for the unavailability of the restitutive reading is syntactic: The result state subevent is not properly represented in the syntax by a phrase that could be scoped over by ‘again’ independently from the causing subevent. If this scenario is empirically real for ditransitives, the question is whether it can be extended to simplex imperfective verbs. That is, can it be the case that \textit{počti} ‘almost’ is unambiguous with \textit{pisa-t’}  ‘write\textsuperscript{IPFV}’ because \textit{pisa-t’} ‘write\textsuperscript{IPFV}’ is like ditransitives: It is subeventally complex, but it does not project a result state denoting phrase that would serve as a landing site for an adverbial? One argument for a negative answer to this question may be based on the interpretation of \textit{-n/t-} passive participles. According to \citet{Paslawska.Stechow2003}, such participles are stativizers: They take a relation between events and states and return a property of states \REF{ex:tatevosov:ii}.

\ea\label{ex:tatevosov:ii}
\sib{\textit{-n/t-}} = $\lambda R_{\langle v, \langle v,t\rangle \rangle}. \lambda s. \exists e [R(s)(e)]$
\z

\noindent With ditransitives, these participles are interpreted just as one would expect given \REF{ex:tatevosov:ii}: \REF{ex:tatevosov:iii} describes a result state that holds of the theme at the reference time.

\ea\label{ex:tatevosov:iii}
\gll V  dannyj  moment  knig-a  otda-n-a  Vase.\\
in  given  moment  book-\textsc{nom}  give-\textsc{n/t-f}  Vasja-\textsc{dat}\\
\glt ‘At the moment the book is in a state of having been given to Vasja.’ 
\z

\noindent \REF{ex:tatevosov:iii} shows that what matters for the interpretation of the \textit{-n/t-} morphology, unlike for ‘almost’, is the presence of the result state in the semantic representation, and not its syntactic manifestation.  With \REF{ex:tatevosov:iii} in mind, consider the minimal pair in \REF{ex:tatevosov:iv}--\REF{ex:tatevosov:v}:

\ea\label{ex:tatevosov:iv}
\gll V  dannyj  moment  statj-a  nakonec  na-pisa-n-a.\\    in  given  moment  article-\textsc{nom}  finally  \textsc{lp}-write-\textsc{n/t-f}\\    
\glt ‘At the moment the paper is finally in a state of having been written.’ 
\z 

\ea[*]
{\gll V  dannyj  moment  statj-a  nakonec  pisa-n-a.\\    in  given  moment  article-\textsc{nom}  finally  write-\textsc{n/t-f}\\   \glt Intended: ‘At the moment the paper is finally a state of having been written.’\label{ex:tatevosov:v}}
\z 

\noindent \REF{ex:tatevosov:v} suggests that the passive participle of \textit{pisa-t’} ‘write\textsuperscript{IPFV}’ is impossible as a description of a result state, unlike its prefixed counterpart in \REF{ex:tatevosov:iv}. (Note that “imperfective” non-prefixed PPs like \textit{pisa-n-a} are available in other semantic environments, \citealt{Borik.Gehrke2018}, \citetv{chapters/05_gehrke}, but there they cannot be analyzed in terms of the stativizer in \REF{ex:tatevosov:ii}.) If simplex verbs had a complex event structure of the $\langle v, \langle v,t \rangle \rangle$ type, the \textit{-n/t-} semantics in \REF{ex:tatevosov:ii} should have been able to apply to it just like to \REF{ex:tatevosov:iii} and \REF{ex:tatevosov:iv}. The fact that it does not happen seems to reinforce the conclusion that such verbs denote a simplex event structure and are not subject to \citeposst{Bondarenko2018} scenario for ditransitives.

Now consider interpretations of the prefixed and non-prefixed verbs under negation. 

\ea    \textsc{Prefixed verb under negation}: ambiguous\label{tat:ex:46}

\gll  Volodja  ne  na-pisa-l  stat’j-u.\\
  Volodja  not  \textsc{lp}-write-\textsc{pst.m}    article-\textsc{acc}\\
\glt ‘Volodja did not write the article.’
\ea No writing activity has been performed.\label{ex:tatevosov:46a}
\ex No writing activity has been completed.\label{ex:tatevosov:46b}
\z
\z

\noindent \REF{tat:ex:46} exhibits at least two readings. On \REF{ex:tatevosov:46a}, the sentence says that there was no writing an article activity and a target state of having been written has never been attained. On \REF{ex:tatevosov:46b}, the existence of the target state is negated, too, but that of the activity is not. On the analysis in \REF{ex:tatevosov:34b}, this ambiguity is expected: since there are two subevental components connected by the conjunction, falsity of any or both of the conjuncts makes the whole negated sentence true, as indicated in the simplified paraphrase in \REF{tat:ex:47}.\footnote{The third logically possible reading (the target state holds but the activity does not) is presumably unavailable for an independent reason: a state of having been written cannot come about without a causal input from an activity.}

\eanoraggedright    It is not the case that [there is a writing an article activity and a target state of having been written].\label{tat:ex:47}
\z

\noindent As with ‘almost’, one should make sure that a simplex imperfective verb is examined in an environment where it is compatible with a culminated construal. \REF{ex:tatevosov:48} is another narrative in the Historical Present.

\eanoraggedright   \textit{Context}: Volodja comes home, sits down at his desk, takes a pen and paper, ...\label{ex:tatevosov:48}

\gll  no  ne  piš-et  pis’m-o,  kak  obyčno,  a  na-bras-yva-et  nebrežnyj  portret  Nad-i.\\
  but  not  \textsc{lp}-write-\textsc{npst.3sg}  letter-\textsc{acc}  as  usual    but  \textsc{lp}-throw-\textsc{yva-npst.3sg}  sketchy  portrait.\textsc{acc}  Nadja-\textsc{gen}\\
\glt ‘but does not write a letter, as usual, but draws a sketch of Nadja’s portrait.’
\z

\noindent \REF{ex:tatevosov:48} is not ambiguous in the way \REF{tat:ex:46} is. It entails that no writing activity occurs. Again, \REF{ex:tatevosov:33b} predicts exactly this. Being based on the predicate of reading events with no target state component, all \REF{ex:tatevosov:48} says is \REF{tat:ex:49}.

\ea    It is not the case that [there was a writing a letter activity].\label{tat:ex:49}
\z

\noindent Crucially, as we have already seen in \REF{tat:ex:42}, telic predicates can describe culminated events in the Historical Present. This guarantees that %the second reading is not unavailable because the target state cannot occur in the evaluation world. I
if a target state was part of the representation of the simplex verb \textit{pisa-t’} ‘write\textsuperscript{IPFV}’, \REF{ex:tatevosov:48} would have had the same range of interpretations as \REF{tat:ex:46}. The fact that this does not happen, once again, provides us with an argument for less subevental complexity of non-prefixed verbs.

This conclusion is reinforced by the behavior of simplex verbs in one more configuration where the imperfective is compatible with the culmination being reached in the evaluation world. Consider \REF{ex:tatevosov:50}.

\ea    \textsc{Non-prefixed stem under negation}: unambiguous\label{ex:tatevosov:50}

\gll  Volodja  ni razu  ne  pisa-l  roman.\\
  Volodja  not once  not  write-\textsc{pst.m}  novel.\textsc{acc}\\
\glt ‘Volodja has never written a novel.’
\ea No writing activity has ever been performed.
\ex Unavailable: No writing activity has ever been completed.
\z
\z

\noindent In \REF{ex:tatevosov:50}, the imperfective clause is taken under the general-factual construal. The general-factual reading has been extensively discussed in the literature (see \citealt{Mehlig1981,Mehlig1995,paduceva96,gronndiss,Groenn2015,Altshuler2012,Arregui.etal2014}, \citetv{chapters/05_gehrke}, among others) and analyzed in a variety of ways. Crucially, under this interpretation, a culminating interpretation is readily available for telic verbs. Hence, the general-factual offers one more environment where event-structural characteristics of \textit{pisa-t’} ‘write\textsuperscript{IPFV}’ and similar verbs can be examined. \REF{ex:tatevosov:50} shows the same pattern as \REF{ex:tatevosov:48}, the one represented in \REF{tat:ex:49}. Because of the lack of subevental complexity, simplex imperfective verbs do not create a structure with a conjunction, hence are not ambiguous under negation.

Before leaving this section, I have to acknowledge one open problem that the view laid out above needs to solve in the future. In a nutshell, this view suggests that simplex imperfective verbs are predicates of events, while their prefixed counterparts appear, in addition, with a result state, the latter thus being a product of prefixation (In \sectref{tat:sec:3}, I will make it a bit more explicit.) 

However, lexical prefixation is possible with simplex perfective verbs, too, as e. g. in \REF{ex:tatevosov:12a}--\REF{ex:tatevosov:12b} above, repeated as \REF{ex:tatevosov:51a}--\REF{ex:tatevosov:51b}. \REF{ex:tatevosov:51c}--\REF{ex:tatevosov:51d} are additional examples of the same type.

\ea   \textsc{Lexical prefixes in combination with a simplex perfective verb}\label{tat:ex:51}
\ea \gll \minsp{[} ot- \minsp{[} reši]\textsuperscript{PFV}]\textsuperscript{PFV} -t’  \\
{} \textsc{lp} {} solve \textsc{inf}\\
\glt  ‘dismiss’  \label{ex:tatevosov:51a}
\ex  \gll \minsp{[} pod- \minsp{[} kupi]\textsuperscript{PFV}]\textsuperscript{PFV}{} -t’  \\  
{} \textsc{lp} {} buy \textsc{inf}\\
\glt  ‘bribe’ \label{ex:tatevosov:51b}
\ex \gll \minsp{[} o- \minsp{[} ses]\textsuperscript{PFV}]\textsuperscript{PFV} -t’\\
{} \textsc{lp} {} sit.down \textsc{inf}\\
\glt ‘settle, settle down’\label{ex:tatevosov:51c}
\ex \gll \minsp{[} za- \minsp{[} sta]\textsuperscript{PFV}]\textsuperscript{PFV} -t’\\
    {} \textsc{lp} {} become \textsc{inf}\\
\glt  ‘catch/find in the middle of sth.’\label{ex:tatevosov:51d}
\z
\z

\noindent Such examples are problematic because simplex perfective verbs, e.g. \textit{reši-t’} ‘make decision, solve\textsuperscript{PFV}’ in \REF{ex:tatevosov:51a}, as a reviewer points out, are arguably as subeventally complex as prefixed perfective verbs discussed in the current section:

\ea   \textsc{Event structure} for \textit{reši-t’} ‘make decision, solve\textsuperscript{PFV}’

$\lambda s. \lambda e.\; ...\; \textsc{decide}(e)\; ... \wedge ...\; \textsc{decided}(s)\; ... \wedge \cnst{cause}(s)(e)$ \label{tat:ex:52}
\z

\noindent \REF{tat:ex:52} for \textit{reši-t’} ‘make decision, solve\textsuperscript{PFV}’ can be supported by the same observations as \REF{ex:tatevosov:34b} for \textit{na-pisa-t’} ‘write\textsuperscript{PFV}’. For instance, \textit{počti reši-t}’ ‘almost solve’ shows the same range of interpretations as \textit{počti na-pisa-t’} ‘almost write’ in \REF{ex:tatevosov:40a}--\REF{ex:tatevosov:40b} above. “Lexical prefixes”, an anonymous reviewer indicates, “are uniformly analyzed as introducing a result state into the event structure of a predicate. In \citet{Rappaport.Hovav.Levin1998}, it is claimed that complex verbs like the English ‘break’ do not appear in resultative constructions precisely because they already contain a result, and adding an additional result is simply not possible \citep[][103]{Rappaport.Hovav.Levin1998}”.

At the moment, I know of no compositional analysis of the verbs like \REF{ex:tatevosov:51a}--\REF{ex:tatevosov:51d}. One possible way of treating them is to assume that Russian is unlike English: a result state can somehow come on top of a complex eventuality description for which its own result state has been specified. If this is the case, \textit{ot-reši-t’} ‘dismiss from office’ would be associated with two result states. One is a state of being decided or solved, lexically assigned to \textit{reši-t’} ‘make decision, solve\textsuperscript{PFV}’, the other, a state of being dismissed from office, would be the contribution of a prefix.

Two event structures that accommodate two result states are shown in \REF{ex:tatevosov:53a}--\REF{ex:tatevosov:53b}, ignoring individual arguments. In \REF{ex:tatevosov:53a} a single eventuality brings about two different states that fall under $\lambda s.\textsc{decided}(s)$ and $\lambda s.\textsc{dismissed}(s)$. In \REF{ex:tatevosov:53b}, an eventuality consisting of eventive and stative parts, \textsc{decide} and \textsc{decided}, brings about an additional state, one from the extension of $\lambda s.\textsc{dismissed}(s)$.

\ea   \textsc{Two hypothetical event structures for \textit{ot-reši-t’} ‘dismiss from office\textsuperscript{PFV}’}\label{ex:tatevosov:53}
\ea $\lambda s.\lambda s'.\lambda e.\; ...\; \textsc{decide}(e)\; ... \wedge ...\; \textsc{decided}(s)\; ... \wedge ...\; \textsc{dismissed}(s')\; ... \newline \wedge \cnst{cause}(s)(e) \wedge \cnst{cause}(s')(e)$\label{ex:tatevosov:53a}
\ex    $\lambda s.\lambda e. \textsc{dismissed}(s) \wedge \cnst{cause}(s)(e) \wedge  \exists e' \exists s' [ e = e' \oplus s' ... \newline \wedge ...\; \textsc{decide}(e')\; ... \wedge ...\; \textsc{decided}(s') \wedge \cnst{cause}(s')(e')]$\label{ex:tatevosov:53b}
\z
\z

\noindent These analytical options, apart from stipulating irreducible cross-linguistic variation between Russian and English, seem to bring up quite a lot of difficult questions that need to be addressed for an analysis to be minimally consistent. 

A general question for \REF{ex:tatevosov:53a} is whether an eventuality can culminate in more than one result state and whether there are other semantic environments where this phenomenon  is observed. For \REF{ex:tatevosov:53b}, which involves sum formation, one may want to find out what triggers this operation and how it is constrained.

Both theories need to determine how the argument structure of a prefixed verb is composed given the event structures in \REF{ex:tatevosov:53a}--\REF{ex:tatevosov:53b} and the fact each of the result states is associated with its own holder. Both \REF{ex:tatevosov:53a} and \REF{ex:tatevosov:53b} face a further empirical complication: They predict that the result state specified for a simplex perfective stem, $\lambda s.\textsc{decided}(s)$, should be a detectable meaning component of a prefixed verb. It is at least dubious if this prediction is borne out for all such verbs. While for \textit{ot-reši-t’} ‘dismiss from office’ it may be argued that a dismissal presupposes making a decision,  it is much less clear if settling in \REF{ex:tatevosov:51c} involves sitting down.

Finally, an analysis has to address the observation that lexically prefixed simplex perfective verbs do not look productive: such verbs form a compact closed class of items, indicating that structures like \REF{ex:tatevosov:53a}--\REF{ex:tatevosov:53b}, if they are real, are not built up freely. I will leave these issues for a further study and will not try to have them settled in what follows.

\subsection{Secondary imperfective verbs}\label{tat:sec:2.2}
\hypertarget{RefHeadingToc164647623}{}
As discussed in \sectref{sec:0}, most secondary imperfective verbs are based on perfective ones, either prefixed or simplex (see \REF{tat:ex:15} in \sectref{sec:1.2}). They pattern with the simplex imperfective verbs as to the range of aspectual interpretations, but event-structurally resemble the lexically prefixed perfective. The latter property is captured by the generalization in \REF{tat:ex:54}. 

 \ea    \textsc{Event structure of secondary imperfective verbs}\label{tat:ex:54}

  The event structure of secondary imperfective verbs, like that of lexically prefixed verbs, contains a result state. 
\z 

\noindent \citet{Tatevosov2016} proposes that the secondary imperfective morpheme takes a relation between events and states generated by a prefixed perfective VP as an argument and, minimally, existentially binds the state variable, creating a derived property of events.\footnote{Nothing in the discussion of event structure of the secondary imperfective seems to hinge on whether eventization shown in \REF{tat:ex:55} exhausts the semantic contribution of the -\textit{yva}- morpheme. In \REF{tat:ex:55}, possible additional components of its meaning are shown as “...”. A similar idea about the contribution of the secondary imperfective, couched withing the DRT framework, is advanced in \textcitetv{chapters/08_muellerreichau}.}  The \textit{-yva-} morpheme is thus  \citeposst{Paslawska.Stechow2003} Eventizer, \REF{tat:ex:55}. 

\ea   \sib{-\textit{yva-}} = $\lambda R.\lambda e.\exists s [... \wedge R(e)(s)]$\label{tat:ex:55}
\z

\noindent In \REF{ex:tatevosov:56b}, the denotation of the aspectless verb phrase representing part of the meaning of \REF{ex:tatevosov:56a} is a predicate of events of reading newspapers that bring about a result state of being read.

\eanoraggedright   \textsc{The secondary imperfective of a lexically prefixed verbs: A derived predicate of events; complex event structure}\label{tat:ex:56}
\ea  \gll Volodja  pro-čit-yva-l  gazet-y.\\
     Volodja  \textsc{lp}-read-\textsc{yva-pst.m}  newspaper-\textsc{acc.pl}\\
\glt ‘Volodja would read newspapers.’\label{ex:tatevosov:56a}\footnote{As with many secondary imperfective members of aspectual triples, the progressive reading of \textit{pro-čit-yva-t’} ‘read\textsuperscript{IPFV}’ in the null context is somewhat degraded; \REF{ex:tatevosov:56a} illustrates its habitual interpretation. Nothing in what follows depends on that.}
\ex    \sib{[ \textit{-yva-} [\textsubscript{VP} \textit{pro-čita-} \textit{gazety} ]]} = $\lambda e.\exists s [\textsc{read}(\textsc{newspapers})(e) \wedge \cnst{cause}(s)(e) \wedge   \textsc{read}(\textsc{newspapers})(s)]$\label{ex:tatevosov:56b}
\z
\z

\noindent \REF{ex:tatevosov:56b}, based on \REF{tat:ex:55}, predicts that the secondary imperfective patterns, event-structurally, with the subeventally complex prefixed perfective in \REF{ex:tatevosov:34b} above, rather than with the simplex imperfective in \REF{ex:tatevosov:33b}.

The most straightforward way of seeing that this prediction is correct is to examine the event-structural behavior of simplex and secondary imperfective elements of aspectual triples like those in \tabref{tab:key:3} above, cf. more examples in \tabref{tab:key:4}. 


\begin{table}
\caption{Aspectual triples}
\label{tab:key:4}
\begin{tabularx}{\textwidth}{XXXX}
\lsptoprule
Meaning & Simplex \newline imperfective & Prefixed \newline perfective & Secondary \newline imperfective\\\midrule
‘drill’ & \textit{sverli-t’} \newline drill-\textsc{inf} & \textit{pro-sverli-t’}\newline \textsc{lp}-drill-\textsc{inf} & \textit{pro-sverl-iva-t’} \newline \textsc{lp}-drill-\textsc{yva-inf}\smallskip\\
‘drink’ & \textit{pi-t’} \newline  drink-\textsc{inf} & \textit{vy-pi-t’}  \newline \textsc{lp}-drink-\textsc{inf} & \textit{vy-pi-va-t’} \newline \textsc{lp}-drink-\textsc{yva-inf}\smallskip\\
‘mow’ & \textit{kosi-t’} \newline mow-\textsc{inf} & \textit{s-kosi-t’} \newline \textsc{lp}-mow-\textsc{inf} & \textit{s-kaš-iva-t’} \newline \textsc{lp}-mow-\textsc{yva-inf}\\
\lspbottomrule
\end{tabularx}
\end{table}

The imperfective members of a triple form minimal pairs contrasting with respect to all the event-structural diagnostics discussed above. 

First, the secondary imperfective does not show the indefinite object alternation: 

\ea   \textsc{Simplex / secondary imperfective and indefinite object alternation}\label{tat:ex:57}
\ea
\gll Po  utr-am  Volodja  čita-l  \minsp{(} svež-ie gazet-y).\\
    on  morning-\textsc{dat.pl}  Volodja  read-\textsc{pst.m}  {} fresh-\textsc{acc.pl} newspaper-\textsc{acc.pl}\\
\glt   ‘In the morning, Volodja would read (fresh newspapers).’\label{ex:tatevosov:57a}
\ex \gll   Po  utr-am  Volodja  pro-čit-yva-l  \minsp{*(} svež-ie gazet-y).\\
    on  morning-\textsc{dat.pl}  Volodja  \textsc{lp}-read-\textsc{yva}-\textsc{pst.m}  {} fresh-\textsc{acc.pl} newspaper-\textsc{acc.pl}\\
\glt  ‘In the morning, Volodja would read fresh newspapers.’\label{ex:tatevosov:57b}
\z
\z

\noindent The simplex imperfective in \REF{ex:tatevosov:57a}, just as in \REF{ex:tatevosov:37a} above, can either project a VP with the direct object or an objectless VP that describes a reading activity. As \REF{ex:tatevosov:57b} shows, if the direct object is not realized in the syntax, the sentence is ungrammatical. In this respect, the secondary imperfective patterns with the prefixed perfective in \REF{tat:ex:58}, see also \REF{ex:tatevosov:37b} above:

\ea\label{tat:ex:58}
\gll Utrom  Volodja  pro-čita-l  \minsp{*(} svež-ie   gazet-y).\\
  in.the.morning  Volodja  \textsc{lp}-read-\textsc{pst.m}  {} fresh-\textsc{acc.pl}   newspaper-\textsc{acc.pl}\\
\glt ‘In the morning, Volodja read fresh newspapers.’
\z

\noindent With respect to the range of interpretations of decomposition adverbs ‘almost’ and ‘again’ and negation, the secondary imperfective, too, patterns with the prefixed perfective rather than with the simplex imperfective. In the Historical Present configuration, parallel to \REF{tat:ex:43}, the secondary imperfective in \REF{ex:tatevosov:59b}, but not the simplex imperfective in \REF{ex:tatevosov:59a}, allows for the narrow scope of \textit{počti} ‘almost’, where it is only associated with the result state:

\eanoraggedright     \textsc{Simplex / secondary imperfective and the interpretation of \textit{počti} ‘almost’}\label{tat:ex:59}

\textit{Context}: Volodja comes home, sits down at his desk, sees a letter from his friend, ...

\ea
\gll i  počti  čita-et  ego,  kogda  razda-ёt-sja zvonok  v  dver' ...\\
    and  almost  read-\textsc{npst.3sg}  it.\textsc{acc}  when  give.out-\textsc{npst.3sg-refl} ring  in  door.\textsc{acc}\\
\glt ‘and almost reads it, when the doorbell rings.’

(i) The reading a letter activity does not start.

(ii) Unavailable: The reading a letter activity is not completed.\label{ex:tatevosov:59a}
\ex
\gll   i  počti  pro-čit-yva-et  ego,  kogda razda-{ёt-sja}  zvonok  v  dver' ...\\
    and  almost  \textsc{lp}-read-\textsc{yva-npst.3sg}  it.\textsc{acc}  when give.out-\textsc{npst.3sg-refl}  ring  in  door.\textsc{acc}\\
\glt ‘and almost reads it, when the doorbell rings.’

(i) The reading a letter activity does not start.

(ii) The reading a letter activity is not completed.\label{ex:tatevosov:59b}
\z
\z

\noindent Availability of the low reading of ‘almost’ separates the secondary imperfective from the simplex imperfective. In this respect the secondary imperfective again patterns together with the prefixed perfective: 

\ea     \textsc{Prefixed perfective and the interpretation of \textit{počti} ‘almost’}\label{ex:tatevosov:60}

\gll  Volodja  počti  pro-čita-l  pis’m-o.\\
  Volodja    almost  \textsc{lp}-read-\textsc{pst.m}  letter-\textsc{acc}\\
\glt ‘Volodja’s reading the letter is close to completion.’ 
\z

\noindent In a similar way, under negation the secondary imperfective is ambiguous in a way the simplex imperfective is not. \REF{tat:ex:61} illustrates that with a general factual sentence parallel to \REF{ex:tatevosov:50}.

\ea    \textsc{Simplex / secondary imperfective under negation}\label{tat:ex:61}
\ea
\gll Volodja  ni  razu  ne  čita-l  \minsp{“} Kapital”.\\
    Volodja  not  once  not  read-\textsc{pst.m}  {} Kapital.\textsc{acc}\\
\glt  ‘Volodja has never read \textit{Das Kapital}.’

(i) No reading activity has ever been performed.

(ii) Unavailable: No reading activity has ever been completed.\label{ex:tatevosov:61a}
\ex
\gll   Volodja   ni   razu  ne  pro-čit-yva-l   \minsp{“} Kapital”.\\
    Volodja  not   once  not  \textsc{lp}-read\textsc{-yva-pst.m} {} Kapital.\textsc{acc}\\
\glt  ‘Volodja has never read \textit{Das Kapital}.’

(i) No reading activity has ever been performed.

(ii) No reading activity has ever been completed.\label{ex:tatevosov:61b}
\z
\z

\noindent \REF{ex:tatevosov:56b} correctly predicts, therefore, that the secondary imperfective sentence is as ambiguous as its perfective prefixed counterpart in \REF{ex:tatevosov:34a}--\REF{ex:tatevosov:34b}. A target state is introduced by the prefix once and for all. Existential binding of the state variable in \REF{ex:tatevosov:56b} does not remove the target state component from the semantic representation. Secondary imperfective verbs are associated with the same complex structure as prefixed perfective verbs, which explains why they pattern together with respect to event-structural diagnostics.\footnote{A reviewer asks what happens with simplex and secondary imperfective verbs outside of aspectual triples discussed in this section%, since “it could very well be that in the case of aspectual triples there is a division of labor between the simple and the secondary imperfective, whereas in the other cases, where there is only one imperfective, it could be that the tests do not yield such clear results”
. I believe that the pattern discussed for \textit{čita-t’} and \textit{pro-čit-yva-t’} ‘write\textsuperscript{IPFV}’ can be replicated with other simplex and secondary imperfective verbs no matter if they are members of an aspectual triple. %This can be illustrated with the simplex \textit{pisa-t’} ‘write\textsuperscript{IPFV}’, which does not have a corresponding secondary imperfective *\textit{na-pis-yva-t’} ‘write\textsuperscript{IPFV}’, and with \textit{vy-pis-yva-t’} ‘write out, copy\textsuperscript{IPFV}’, the secondary imperfective of \textit{vy-pisa-t’} ‘write out, copy\textsuperscript{PFV}’, for which \textit{pisa-t’} ‘write\textsuperscript{IPFV}’ cannot be considered a simplex partner due to the difference in lexical meaning. We have already seen in \REF{tat:ex:43} that for \textit{pisa-t’} ‘write\textsuperscript{IPFV}’ the low reading of \textit{počti}  ‘almost’ is unavailable, just as for the simplex \textit{čita-t’} ‘read\textsuperscript{IPFV}’ in \REF{ex:tatevosov:59a}.
Consider \REF{ex:tatevosov:vi} with \textit{vy-pis-yva-t’} ‘write out, copy\textsuperscript{IPFV}’.

\eanoraggedright \textit{Context}: Volodja needs a quotation for his term paper; he comes home, sits down at his desk, takes a pen and a piece of paper, ...\label{ex:tatevosov:vi}

\gll i  počti vy-pis-yva-et  dlinnuju  citat-u   iz knig-i,  kogda  razda-{ёt-sja}  zvonok  v  dver' ...\\      
and  almost  \textsc{lp}-write-\textsc{yva-npst.3sg}  long  quotation  from      book-\textsc{gen}  when  give.out-\textsc{npst.3sg-refl}   ring  in  door.\textsc{acc}\\
\glt ‘and almost copies a long quotation from a book, when the doorbell rings.’  
\ea The copying activity does not start.  
\ex The copying activity is not completed. 
\z
\z 

\noindent Unlike in \REF{tat:ex:43} and \REF{ex:tatevosov:59a}, but in parallel with \REF{ex:tatevosov:59b}, the secondary imperfective is compatible with the two readings of \textit{počti} ‘almost’, suggesting a complex event structure for \textit{vy-pis-yva-t’} ‘write out, copy\textsuperscript{IPFV}’. There may be secondary imperfectives for which one of the two readings of ‘almost’ is more readily accessible than the other. For instance, for \textit{počti} \textit{pod-pis-yva-et} ‘almost signing’ in the same environment the low reading is more difficult to get, presumably because one has to accommodate a somewhat unnatural context in which it is relevant to convey that the agent has started writing his %her
signature under a document, is maximally close to completion, but is not yet done. Nevertheless, I believe that \REF{ex:tatevosov:vi} represents a regular and systematic pattern.}

In the next section, I focus on the result state component and explore its influence on the interpretation of lexically prefixed verbs in more detail. 

\section{Lexical prefixes, result states, and projection of arguments}\label{tat:sec:3}
\hypertarget{RefHeadingToc164647624}{}\subsection{Introducing the typology}\label{sec:3.1}
\hypertarget{RefHeadingToc164647625}{}
The previous section offers a few pieces of empirical evidence suggesting that lexically prefixed perfective verbs show subevental complexity in contrast with the corresponding simplex verbs. According to \REF{ex:tatevosov:34b}, repeated as \REF{tat:ex:62}, VPs projected by such verbs denote a relation between events and (result) states:

\ea   \textsc{VPs projected by lexically prefixed perfective verbs}\label{tat:ex:62}

  \sib{[\textsubscript{VP} \textit{na-pisa- pis’mo} ]} = $\lambda s.\lambda e. \textsc{write}(\textsc{letter})(e) \wedge \cnst{cause}(s)(e) \wedge \textsc{written}(\textsc{letter})(s)$ 
\z

\noindent The presence of a result state in the event structure of lexically prefixed verbs uniformly tells them apart from simplex imperfective verbs. 

The class of lexically prefixed verbs is, however, not homogeneous. Depending on the setting of a number of semantic parameters, discussed in \citet{BabkoMalaya1999,Romanova2007,Tatevosov2016,Biskup2019}, and many others, lexically prefixed verbs fall into a number of subtypes, summarized in Figure \ref{fig:1}. The terminology used in Figure \ref{fig:1} (\textit{restrictive}, \textit{(non-)argument-sharing} [(n)AS], \textit{mo\-nad\-ic}, etc.) will be introduced and explained shortly.

      \begin{figure}
      \small
      \caption{Typology of lexically prefixed verbs}
	     \begin{forest}
[Lexically prefixed verbs [Compositional [Monadic Ps [AS Ps [Non-restrictive\\(``empty''/\\``pure perfectivizing'') [\textit{na-pisa-t’}\\‘write\textsuperscript{PFV}’\\\textit{pro-čita-t’}\\‘read\textsuperscript{PFV}’ ]] [Restrictive [\textit{za-pisa-t’}\\`write down'\\\textit{za-čita-t’}\\`read out' ]] ] [nAS Ps [\textit{pod-pisa-t’}\\`sign'] ] ]  
[Dyadic Ps [\textit{v-pisa-t’} \\`write into sth.'\\\textit{vy-čita-t’}\\`read from'] ] ]
[Non-compositional [\textit{o-pisa-t’} \\`describe']
]]
         \end{forest}\label{fig:1}
         \end{figure}

%[(“pure perfectivizing”,\\or “empty”)\\\textit{na-pisat’}\\‘write\textsuperscript{PFV}’\\\textit{pro-čitat’}\\‘read\textsuperscript{PFV}’ ]

%[\textit{za-pisa-t’}\\`write down'\\\textit{za-čita-t’}\\`read out' ]

Membership in one of the subtypes in Figure \ref{fig:1} has consequences for the argument structure of a prefixed verb, as well as for the perceived “semantic distance” between a simplex verb and its derived counterpart.  

Throughout this section, I will mostly focus on the prefixed configurations derived from transitive lexical verbs like \textit{pisa-t’} ‘write’ or \textit{čita-t’} ‘read’.\footnote{Throughout this section I will be abstracting away from prefixed verbs whose derivation involves both a prefix and the reflexive clitic -\textit{sja}, e.g. \textit{u-čita-t’-sja} ‘get one’s fill of reading’. See \citet{Zaucer2009} for an analysis of this type of verbs in Slovenian.} The generalizations discussed below are extendable to unaccusative and unergative lexical verbs as well, possibly with minor adjustments; see \citet[70--122]{Romanova2007} for an extensive discussion. For a recent alternative to the typology in \figref{fig:1}, see \citet[ch. 3]{Biskup2019}, who elaborates on and extends the system based on \citet{Bergsma.etal2010}. 

As we have seen above, lexical prefixes contribute a result state into the semantic representation of a verb~/~VP. In what follows, I will be making a straightforward assumption about the descriptive content of predicates of such states: 

\eanoraggedright   \textsc{Content of state descriptions}\label{tat:ex:64}

  Descriptive properties of result states associated with lexical prefixes in a given lexical configuration V are specific enough to differentiate states introduced by different prefixes in combination with V. 
\z

\noindent \textit{Na-} in \textit{na-pisa-t’} ‘write’ in \REF{tat:ex:62} is understood as the description of states of having been written, $\lambda s.\textsc{written}(x)(s)$, for a fixed $x$. \textit{Pod}- in \textit{pod-pisa-t’} ‘sign’, on the other hand, comes with the description of different type of states, states of having been signed, $\lambda s.\textsc{signed}(x)(s)$, for a fixed $x$. Crucially, since the non-derived verb lacks a state argument altogether, all such information can only be the contribution of a prefix, which makes some or other version of \REF{tat:ex:64} inevitable. 

A separate question is how this information is computed and by what mechanism. This constitutes one of the most difficult topics in the semantic literature on Slavic verbs, which I will address in \sectref{sec:3.3}. Anticipating this discussion, I can mention two specific ways of thinking about such a mechanism. On the one hand, prefixes can be conceived of as collections of correspondences between verb stems and state descriptions. On this view, part of the lexical specification of \textit{pod}- in \textit{pod-pisa-t’} ‘sign\textsuperscript{PFV}’ is a function that maps the stem \textit{pisa}- ‘write’ to a state of being signed, “write $\mapsto$ signed”. On the other hand, one can try to assign a general  meaning to every prefix (or a restricted set of meanings, if one assumes that prefixes are polysemous). Since prefixes are mostly drawn from the prepositional inventory, such a meaning can naturally be thought of in terms of spatial relations. For \textit{pod}- that would yield ‘under’, the meaning of its prepositional counterpart, cf. \textit{pod stolom} ‘under the table’. The state of being signed, then, would be derived as an outcome of a semantic interplay between the stem, the prefix and event structural information. 

For the purposes of the current section, however, I will be abstracting away from this dichotomy. In order to keep the presentation simple, I will assume that prefixes are merely relations between individuals and states, and that descriptive properties of those states are always specified in every lexical context, in accordance with \REF{tat:ex:64}: 

\ea
\ea  \sib{\textit{na-}} = $\lambda x.\lambda s.\textsc{written}(x)(s)$, in the context of \textit{pisa}- ‘write’
\ex   \sib{\textit{pod-}} = $\lambda x.\lambda s. \textsc{signed}(x)(s)$, in the context of \textit{pisa}- ‘write’
\z
\z

\noindent As was already mentioned above, the discussion below is based on the assumption that prefixes and the secondary imperfective do not denote directly any aspectual  verb stem meanings. Verb phrases, with all the relevant pieces of morphology merged with the stem, are aspectless predicates of events, relations between events and states or the like. 

\subsection{Compositionality}\label{tat:sec:3.2}
\hypertarget{RefHeadingToc164647626}{}
The first parameter that divides lexically prefixed verbs into two major groups is whether \textit{only} events from the original extension of a non-prefixed verb can be found in the extension of its prefixed counterpart.\footnote{A technical clarification is due at this point. Above, it was assumed that transitive verbs denote relations between individuals and events, hence their extensions are not (characteristic functions of) sets of eventualities, but sets of ordered pairs of individuals and eventualities. An expression like “an eventuality is in the extension of a non-derived verb” is therefore to be understood as a shorthand for “an eventuality is an eventive member of an ordered pair in the extension of a non-derived verb”.}  This parameter will be called \textsc{compositionality} and defined as in \REF{ex:tatevosov:66}.

\eanoraggedright   \textsc{Compositionality}\label{ex:tatevosov:66}

Both inside and outside of a lexically prefixed configuration, the verb stem denotes events with the same descriptive properties. 
\z

\noindent \REF{ex:tatevosov:66} tells apart prefixed verbs like \REF{ex:tatevosov:67} and \REF{ex:tatevosov:68}, on the one hand, and \REF{tat:ex:69}, on the other.

\ea\label{ex:tatevosov:67}
\gll Volodja na-pisa-l pis’m-o.\\
  Volodja  \textsc{lp}-write-\textsc{pst.m} letter-\textsc{acc}\\
\glt ‘Volodja wrote a letter.’
\z

\ea\label{ex:tatevosov:68}
\gll Volodja pod-pisa-l dogovor.  \\
  Volodja  \textsc{lp}-write-\textsc{pst.m}  agreement.\textsc{acc}\\
\glt ‘Volodja signed an agreement.’
\z

\ea\label{tat:ex:69}
\gll Volodja o-pisa-l scen-u  prestupleni-ja.\\
  Volodja  \textsc{lp}-write-\textsc{pst.m}  scene-\textsc{acc}  crime-\textsc{gen}\\
\glt ‘Volodja described the crime scene.’
\z

\noindent Events that \textit{na-pisa-t’} ‘write’ in \REF{ex:tatevosov:67} and \textit{pod-pisa-t’} ‘sign’ in \REF{ex:tatevosov:68} describe are writing events, those that fall under the extension of \textit{pisa-t’} ‘write’, the prefixless counterpart of \REF{ex:tatevosov:67}--\REF{tat:ex:69}. The agent who writes a letter or signs an agreement has to write. Both \REF{ex:tatevosov:67} and \REF{ex:tatevosov:68} thus entail \REF{ex:tatevosov:70} (abstracting away from the aspectual difference between them).

\ea\label{ex:tatevosov:70}
\gll Volodja  \minsp{(} čto-to)  pisa-l.\\
  Volodja  {} something.\textsc{acc}  write-\textsc{pst.m}\\
\glt ‘Volodja was writing / has written (something).’
\z

\noindent Whatever validates \REF{ex:tatevosov:70} as the description of a writing eventuality should also be part of the description of the eventive components of \REF{ex:tatevosov:67}--\REF{ex:tatevosov:68}. To put it differently, events that fall under \textit{na-pisa-t’} ‘write’ and \textit{pod-pisa-t’} ‘sign’ form (not necessarily proper) subsets of the extension of \textit{pisa-t’} ‘write’. This is shown in \REF{ex:tatevosov:71}.\footnote{Note that the object of a prefixed verb derived from the simplex imperfective \textit{pisa-t’} can but does not have to be the theme of writing. In \REF{tat:ex:64}, for example, the agreement is an object of signing without being  an object of writing. For this reason, the right-hand part of the formulas in \REF{ex:tatevosov:71}--\REF{ex:tatevosov:72} contain the most inclusive set of writing events: the set of wiring events in which something is written. This topic will be taken up shortly, in \sectref{sec:3.3}.}

\ea\label{ex:tatevosov:71}
\ea    $\{e: {\exists} s[$\sib{\textit{na-pisa- pis’mo}}$(s)(e)]\} \subseteq \{e: \exists x[$\sib{\textit{pisa}-}$(x)(e)]\}$\label{ex:tatevosov:71a}
  \ex    $\{e: \exists s[$\sib{\textit{pod-pisa- dogovor}}$(s)(e)]\} \subseteq \{e: \exists x[$\sib{\textit{pisa}-}$(x)(e)]\}$\label{ex:tatevosov:71b}
\z
\z

\noindent In contrast, \REF{tat:ex:69} fails to entail \REF{ex:tatevosov:70}. \REF{tat:ex:69} can but does not have to involve a writing activity: to describe an object one does not have to write; for instance, a verbal description would suffice. Therefore%Since not only writing events fall under \textit{o-pisa-t’} ‘describe’
, the set of describing events is not a subset of writing events:%included into it: 

\ea   $\{e: \exists s[$\sib{\textit{o-pisa- scenu prestuplenija}}$(s)(e)]\} \nsubseteq \{e: \exists x[$\sib{\textit{pisa}-}$(x)(e)]\}$\label{ex:tatevosov:72}
\z

\noindent If the prefix only contributes a property of (result) states into a complex structure, events that bring those states about must be introduced by the verb stem. One has to conclude, therefore, that the extensions of \textit{pisa-} in \textit{o-pisa-t’} ‘describe’ and of the unprefixed \textit{pisa}- are different. After the prefixation of \textit{o-}, the extension of \textit{pisa-} widens: it acquires exactly those eventualities where something is described with no writing being involved.\footnote{\label{fn:30}In addition, the extension loses those writing events that do not take an argument into a state of being described. This is, however, what \textit{o-pisa-t’} ‘describe’ has in common with \textit{pod-pisa-t’} ‘sign’. Not all writing events culminate in a state of having been signed. Likewise, there are writing events irrelevant for bringing about a state of being described. I will address such “narrowing-down” effects of lexical prefixation shortly, in \sectref{sec:3.3.2}--\sectref{sec:3.3.3}. I am grateful to an anonymous reviewer who encouraged to introduce this qualification here.}  There is no obvious way of deriving this widening in a compositional way. \REF{ex:tatevosov:66} captures exactly this property of lexically prefixed verbs like \textit{o-pisa-t’} ‘describe’. On \textit{o-pisa-t’}, \REF{ex:tatevosov:66} fails.

While the eventive component of \textit{o-pisa-t’} ‘describe’ is not included into that of \textit{pisa-t’} ‘write’, \REF{ex:tatevosov:72}, these components do have a non-empty overlap. This overlap obtains if a description is done in writing. But one can easily find prefixed verbs with no such an overlap. For example, \textit{ot-kry-t’} ‘open’ is derived from \textit{kry-t’} ‘cover’ by the lexical prefix \textit{ot-} (cf. the preposition \textit{ot} ‘from’; diachronically, \textit{ot-kry-t’} ‘open’ is ‘from-cover’). In an opening eventuality nothing is getting covered, so for \textit{kry-t’} ‘cover’ and \textit{ot-kry-t’} ‘open’ \REF{ex:tatevosov:73} holds.

\ea    $\{e: \exists s[$\sib{\textit{otkry- dver’}}$(s)(e)]\} \cap \{e: \exists x[$\sib{\textit{kry-}}$(x)(e)]\} = \varnothing$\label{ex:tatevosov:73}
\z

\noindent Therefore, prefixes like \textit{ot}- in \textit{ot-kry-t’} ‘open’ or \textit{o-} in \textit{o-pisa-t’} ‘describe’ are \textsc{non-compositional.} They induce a non-compositional shift in the denotation of a verb stem whereby after prefixation it starts describing events that had not been there originally. \textsc{Compositional} are prefixes with which this does not happen. 

In what follows, I will focus on compositional prefixes like \textit{na}- and \textit{pod}- in \REF{ex:tatevosov:67} and \REF{ex:tatevosov:68}, for which \REF{ex:tatevosov:66} does hold.

\subsection{Argument structure of a prefix}
\label{sec:3.3}
\hypertarget{RefHeadingToc164647627}{}
(Result) states that compositional lexical prefixes introduce into the event structure can be further divided into several classes. The most basic division is induced by whether a state involves one or two participants. Accordingly, the argument structure of a prefix can be \textsc{monadic} or \textsc{dyadic}, to use \citeposst{Hale.Keyser2002} term.

\eanoraggedright   \textsc{Participants of a result state}\label{ex:tatevosov:74}

\textsc{Monadic prefixes} are associated with result states with a single participant, its holder. \textsc{Dyadic prefixes} introduce states corresponding to a relation between two entities, figure and ground, and take two arguments.  
\z

\noindent An important background assumption underlying \REF{ex:tatevosov:74} is that a prefix is only considered dyadic if both of its arguments can be associated with overt material in the syntax. Otherwise, a prefix is treated as monadic. This strategy, which follows \citet{Romanova2007}, is not the only admissible one, of course. In other systems, \citet{BabkoMalaya1999} and \citet{Biskup2019} among them, (near) all prefixes are dyadic, but in certain cases one of the arguments is obligatorily left implicit for some or other reason.

These two approaches seem to manifest different theoretical desiderata. \REF{ex:tatevosov:74} prioritizes minimizing implicit arguments in the argument structure of a prefix whose empirical reality is difficult or impossible to assess. The alternative aims at keeping the semantics and syntax of prefixes as uniform as possible. Whether the two approaches make different empirical predictions should ultimately depend on the representational and interpretational status of implicit arguments (see e.g. \citealt{Landau2010,Bhatt.Pancheva2017}) assumed within the latter type of theory. This topic, however, requires more examination than one can find in the current literature.

Monadic prefixes fall into three subtypes characterized below in \sectref{sec:3.3.1}--\sectref{sec:3.3.4}. In \sectref{sec:3.3.5}, dyadic prefixes are briefly discussed.

\subsubsection{Monadic prefixes and argument identity}\label{sec:3.3.1}
\hypertarget{RefHeadingToc164647628}{}
\REF{ex:tatevosov:67}--\REF{ex:tatevosov:68} with \textit{na-pisa-t’} ‘write\textsuperscript{PFV}’ and \textit{pod-pisa-t’} ‘sign\textsuperscript{PFV}’ above are both examples of event structures in which the prefix is monadic. A characteristic property of verbs with monadic prefixes is that the argument corresponding to the holder of result state is typically realized in the direct object position, rather than appearing in the subject or indirect object slots.\footnote{I am grateful for an anonymous reviewer for encouraging me to make this clarification.}

\textit{Na-pisa-t’} and \textit{pod-pisa-t’} instantiate two major subtypes of event structures with monadic prefixes. \textit{Na-pisa-t’} projects VPs where the object is \textsc{subcategorized} for by the verb stem. This is not the case with \textit{pod-pisa-t’}. 

In \REF{ex:tatevosov:67}, \textit{pis’mo} ‘letter’ is construed  both as the theme of writing and the holder of a result state of having been written. In contrast, \textit{pis’mo} ‘letter’ in \REF{ex:tatevosov:68} is the holder of the result state of having been signed without being the theme of writing.

One more illustration appears in \REF{ex:tatevosov:75}--\REF{ex:tatevosov:76}.

\ea\label{ex:tatevosov:75}
\gll Volodja  e-l  \minsp{\{} jablok-o  /  \minsp{\#} puz-o\}.\\
  Volodja     eat-\textsc{pst.m}    {} apple-\textsc{acc} {} {} belly-\textsc{acc}\\
\glt   (Intended:) ‘Volodja was eating/ate \{an apple / a belly\}.’
\z

\ea\label{ex:tatevosov:76}
 \gll Volodja na-e-l  \minsp{\{} puz-o  /    \minsp{*} jablok-o\}.\\
  Volodja \textsc{lp}-eat-\textsc{pst.m} {} belly-\textsc{acc} {} {} apple-\textsc{acc}          \\
\glt  (Intended:) ‘Volodja acquired \{a belly / an apple\} by eating.’
\z

\noindent The verb \textit{es-t’} ‘eat’, when it appears in the unprefixed configuration in \REF{ex:tatevosov:75}, is a verb of consumption that takes a DP denoting edible substances as its object. In \REF{ex:tatevosov:76}, the object of \textit{na-est’} ‘acquire by eating’, derived from \textit{es}-\textit{t’} ‘eat’ by the lexical prefix \textit{na-}, is \textit{puzo} ‘belly’, which is not consumed but effected in the course of eating. In \REF{ex:tatevosov:75}, it can only be understood as the object of eating. Similar effects have been observed in \REF{ex:tatevosov:26a}--\REF{ex:tatevosov:26b} in \sectref{sec:1}, where lexical prefixes were said, descriptively, to be able to introduce non-subcategorized arguments or change the thematic properties of the object.

These observations lead to establishing two classes of monadic prefixes in \REF{ex:tatevosov:77}, based on identity of the holder of a result state and the internal argument of a verb stem. I will call them \textsc{argument-sharing} (AS) \textsc{prefixes} and \textsc{non-argument-sharing} (nAS) \textsc{prefixes}:

\eanoraggedright   \textsc{Argument identity}\label{ex:tatevosov:77}
\eanoraggedright \textsc{Argument-sharing prefixes}: The argument of a prefix, the holder of a result state, is identical to the  internal argument of the non-derived verb stem.
  \ex   \textsc{Non-argument-sharing prefixes}: The holder of a result state is different from the internal argument of the non-derived verb stem.
\z
\z

\largerpage
\noindent A simple heuristic that separates the two types of prefixes is based on whether a sentence with a lexically prefixed verb entails a corresponding prefixless sentence with the same object (again, abstracting away from aspectual differences). Whereas the entailment from \textit{na-pisa-t’ pis’mo} to \textit{pisa-t’ pis’mo} goes through, that from \textit{pod-pisa-t’ pis’mo} to \textit{pisa-t’ pis’mo} does not, \REF{ex:tatevosov:78a}--\REF{ex:tatevosov:78b}, or, in set-theoretic terms, \REF{ex:tatevosov:79a}--\REF{ex:tatevosov:79b}.

\ea\label{ex:tatevosov:78}
\ea\label{ex:tatevosov:78a}
\gll    Nadja na-pisa-l-a pis’m-o. $\rightarrow$ Nadja pisa-l-a pis’m-o. \\
      Nadja \textsc{lp}-write-\textsc{pst-f}  letter-\textsc{acc}   {}   Nadja write-\textsc{pst-f}  letter-\textsc{acc}\\
\glt  ‘Nadja wrote the letter.’   entails   ‘Nadja was writing the letter.’
\ex\label{ex:tatevosov:78b}
\gll    Nadja   pod-pisa-l-a   pis’m-o   $\nrightarrow $    Nadja   pisa-l-a   pis’m-o \\
Nadja \textsc{lp}-write-\textsc{pst-f}  letter-\textsc{acc}  {}    Nadja write-\textsc{pst-f}  letter-\textsc{acc}\\
\glt     ‘Nadja signed the letter.’  does not entail    ‘Nadja was writing the letter.’
\z
\z

\ea\label{ex:tatevosov:79}
\ea   $\{e: \exists s$ \sib{\textit{na-pisa- pis’mo}}$(s)(e)\} \subseteq \{e:$ \sib{\textit{pisa- pis’mo}}$(e) \}$\label{ex:tatevosov:79a}
\ex    $\{e: \exists s$ \sib{\textit{podpisa- pis’mo}}$(s)(e)\} \nsubseteq \{e:$ \sib{\textit{pisa- pis’mo}}$(e) \}$\label{ex:tatevosov:79b}
\z
\z

\noindent \textit{Na-pisa-t'} ‘write’ thus contains an argument-sharing prefix, whereas \textit{pod-pisa-t’} ‘sign’ and \textit{na-est’} ‘acquire by eating’ in \REF{ex:tatevosov:76} do not.

\subsubsection{nAS prefixes: putting events and states together}\label{sec:3.3.2}
\hypertarget{RefHeadingToc164647629}{}
The phenomena discussed in the previous section pose a number of questions about how a result state description projected by a prefix is integrated into a larger configuration. What is the relationship between the verb stem and the projection of a prefix? How is the argument structure of their combination computed? How can the difference between AS and nAS prefixes be accounted for?

These questions are not specific to prefixal configurations in Slavic languages. Exactly the same type of issues are brought out by other types of resultative constructions cross-linguistically, e.g., by the familiar resultatives in English and German. 

Consider the resultatives in \REF{ex:tatevosov:80}, which are nAS, since the objects \textit{the pub, seine Familie} ‘his family’ and \textit{the baby} are not subcategorized for by the verb.

\ea\label{ex:tatevosov:80}
\ea They drank the pub dry. \hfill \citep{Rappaport.Hovav.Levin2001}\label{ex:tatevosov:80a}
\ex   
\gll Er hat seine Familie magenkrank gekocht.\\
he \textsc{aux} self's family stomach.sick cooked\\\hfill \citep{Kratzer2005}\label{ex:tatevosov:80b}
\glt   ‘He cooked his family stomach sick.’
\ex   John sang the baby asleep. \hfill \citep{Rothstein2004}\label{ex:tatevosov:80c}
\z
\z

\noindent The composition of resultative XPs like those in \REF{ex:tatevosov:80a}--\REF{ex:tatevosov:80c} has been addressed in a variety of studies (e.g.  \citealt{Rappaport.Hovav.Levin2001,Rothstein2004,Kratzer2005}, to mention just a few). Even though details of specific proposals vary, there seems to be a fundamental agreement that arguments not subcategorized for by the verb stem are licensed as arguments of a result state expression like \textit{dry}, \textit{magenkrank} ‘stomach sick’, and \textit{asleep}.

Along similar lines, \citet{BabkoMalaya1999,Romanova2007,Gehrke2008,Tatevosov2015a,Biskup2019} all argue that in a prefixal configuration like \REF{ex:tatevosov:68} or \REF{ex:tatevosov:76} the object DP originates as an argument of a prefix. \citeauthor{Gehrke2008}'s (\citeyear{Gehrke2008}: 164ff.) proposal, for example, assumes that lexical (= internal, in her terminology) prefixes head a functional projection that V takes as its complement, the object DP being merged in its Spec, see \figref{fig:2}.

\begin{figure}
\caption{Proposal according to \citet{Gehrke2008}}
\begin{forest}
[VP [\textit{Spec} ] [V' [V ] [PredP [DP\textsubscript{int} ] [Pred' [Pred [\textsc{internal prefix} ] ] [(XP) ] ] ] ] ]
\end{forest}\label{fig:2}
\end{figure}

Less uniformity across existing theories is attested as to how a result state description is integrated into a larger configuration and what happens to the internal argument of the verb stem, if present. 

I will introduce a toy fragment that offers a possible treatment of nAS-prefixed verbs like \textit{pod-pisa-t’} ‘sign’ in \REF{ex:tatevosov:68}, repeated as \REF{ex:tatevosov:82}, and discuss  several proposals found in the literature against this background.

\ea\label{ex:tatevosov:82}
\gll Volodja  pod-pisa-l  dogovor.  \\
  Volodja  \textsc{lp}-write-\textsc{pst.m}  agreement.\textsc{acc}\\
\glt ‘Volodja signed an agreement.’
\z

\noindent This fragment will be used to assess theories that vary along the two parameters in \REF{ex:tatevosov:83} and \REF{ex:tatevosov:84}.

\ea   \textsc{Result states}\label{ex:tatevosov:83}
\ea Result states are integrated by a special rule of construal.\label{ex:tatevosov:83a}
\ex Result states are introduced by a special morphosyntactic device.\label{ex:tatevosov:83b}
\z
\z

\ea   \textsc{Internal arguments}\label{ex:tatevosov:84}
\ea Internal arguments get existentially bound.\label{ex:tatevosov:84a}
\ex  Internal arguments are not there to begin with.\label{ex:tatevosov:84b}
\z
\z

\noindent First, let us assume, as before, that the semantic representation of a non-prefixed transitive verb looks like \REF{ex:tatevosov:85}, taking \textit{pisa}- ‘write’ as an example:

\ea   \sib{\textit{pisa-}} = $\lambda x.\lambda e. \textsc{write}(x)(e)$\label{ex:tatevosov:85}
\z

\noindent Following \REF{tat:ex:64}, two irreducible ingredients of the semantic representation of a monadic prefix should be a predicate of states, specifying its descriptive properties, and an individual argument construed as the holder of that state. For the prefix \textit{pod}- in \REF{ex:tatevosov:82} this yields \REF{ex:tatevosov:86}.

\ea   \sib{\textit{pod-}} = $\lambda x.\lambda s. \textsc{signed}(x)(s)$, in the lexical context of \textit{pisa}- ‘write’\label{ex:tatevosov:86}
\z

\noindent As was mentioned above for the moment I am putting aside the question of how the descriptive properties of a specific state in a specific lexical context are computed. This will be  the topic of \sectref{sec:3.4}.

With these assumptions, a possible outline of the system that combines the denotations of the two expressions may look like \REF{ex:tatevosov:87a}--\REF{ex:tatevosov:87h}.

\ea\label{ex:tatevosov:87}
\ea    \sib{\textit{pod-}} $= \lambda x.\lambda s. \textsc{signed}(x)(s)$        \hfill (Assumption in \REF{tat:ex:64})\label{ex:tatevosov:87a}
\ex    \sib{\textit{pisa}-} $= \lambda x.\lambda e. \textsc{write}(x)(e)$         \hfill (Lexicon)\label{ex:tatevosov:87b}
\ex    \sib{\textit{pis’mo pod-}} $= \lambda s. \textsc{signed}(\textsc{letter})(s)$      \hfill (FA)\label{ex:tatevosov:87c}
\ex  \sib{CS} $= \lambda P.\lambda s.\lambda e. P(s) \wedge \cnst{cause}(s)(e)$      \hfill (Causative Shift)\label{ex:tatevosov:87d}
\ex    \sib{CS [ \textit{pis’mo pod-} ]} $= \lambda s.\lambda e. \textsc{signed}(\textsc{letter})(s) \wedge \cnst{cause}(s)(e)$   \hfill (FA)\label{ex:tatevosov:87e}
\ex    \sib{$\exists_x$} $= \lambda R. \lambda e. \exists x [R(x)(e)]$         \hfill (Existential Closure)\label{ex:tatevosov:87f}
\ex    \sib{$\exists_x$ \textit{pisa}-} $= \lambda e.\exists x [\textsc{write}(x)(e)]$ \hfill (FA)\label{ex:tatevosov:87g}
\ex \sib{[$\exists_x$ \textit{pisa}-] [CS [ \textit{pis’mo pod-} ]]} \hfill (Event Identification)\\$= \lambda s.\lambda e.\exists x [\textsc{write}(x)(e) \wedge \textsc{signed}(\textsc{letter})(s) \wedge \cnst{cause}(s)(e)]$\label{ex:tatevosov:87h}
\z
\z

\noindent In \REF{ex:tatevosov:87c}, the denotation of the prefix, \REF{ex:tatevosov:87a}, merges with its argument, the holder of a state that falls under $\lambda s.\textsc{signed}(x)(s)$. (For simplicity, I keep on representing DPs as individual constants. Also, I ignore the subsequent movement of the DP for case reasons.)

The resulting predicate of states undergoes Causative Shift, \REF{ex:tatevosov:87d}--\REF{ex:tatevosov:87e}, which makes a relation between events and states, causally connected, out of a state description. As a consequence, states from the extension of \REF{ex:tatevosov:87c} are construed as result states, those brought about by a causing eventuality. A rationale behind Causative Shift can be thought of in terms of naturalness of a state. It captures the intuition that states like being signed described by \textit{pod-} in examples like \REF{ex:tatevosov:82} are not found in non-derived environments across languages. While languages systematically have lexical items like \textit{sick} or \textit{open} whose denotations are states of being sick and open, no language that we know possesses a lexical item for states like being signed. Such states only come out as denotations of derived state descriptions, e.g. “passive participles” in Germanic or Slavic languages.\footnote{\citet[and subsequent work]{KoontzGarboden.Beavers2017} make a similar case for non-derived change of state verbs whose roots introduce a state description. They argue such roots fall into two types: some come with entailments of change, being similar in that respect to a state of being signed in \REF{ex:tatevosov:82}, while others (“property concept states”) do not. I am grateful to an anonymous reviewer for turning my attention to this parallelism.}

In \REF{ex:tatevosov:87f}--\REF{ex:tatevosov:87g}, the internal argument of \textit{pisa-t’} ‘write’ from \REF{ex:tatevosov:87b} undergoes existential closure, yielding a predicate of events. This captures the pattern in \REF{ex:tatevosov:71b} whereby sentences with \textit{pod-pisa-t’} ‘sign’ entail that something has been written in the course of signing a letter.

Finally, the output of Causative Shift in \REF{ex:tatevosov:87e} and the verb stem denotation in \REF{ex:tatevosov:87g} combine by a version of \citeposst{Kratzer1996} Event Identification. In \REF{ex:tatevosov:87h}, Event Identification restricts the eventive members of the relation in \REF{ex:tatevosov:87e} to those that fall under \REF{ex:tatevosov:87g}. \REF{ex:tatevosov:87h} is thus the relation between events and states such that an event in which something has been written causes the letter to enter a state of being signed. This is exactly the meaning of \REF{ex:tatevosov:82}.

A similar derivational scenario can be assumed for unergative lexically prefixed verbs like \textit{za-rabota-t’} ‘earn’ in \REF{ex:tatevosov:24a}--\REF{ex:tatevosov:24b} above or \textit{ob-rabota-t'} ‘process’ in \REF{ex:tatevosov:88a}, derived from \textit{rabota-t'} ‘work’ in \REF{ex:tatevosov:88b}, where the lexical prefixation leads to transitivization.

\ea\label{ex:tatevosov:88}
\ea[]{\gll   Komitet  ob-rabota-l  sto  zajavok.  \\
committee  \textsc{lp}-work-\textsc{pst.3sg}  hundred.\textsc{acc}  application.\textsc{gen.pl}  \\
\glt  ‘The committee processed one hundred applications.’\label{ex:tatevosov:88a}}
\ex[*]{\gll   Komitet  ra-bota-l  sto  zajavok.  \\
committee  work-\textsc{pst.3sg}  hundred.\textsc{acc}  application.\textsc{gen.pl}\\
\glt  Intended: ‘The committee was working one hundred applications.’\label{ex:tatevosov:88b}}
\z
\z

\noindent Arguably, the only difference between \REF{ex:tatevosov:88a} and \REF{ex:tatevosov:82} is that unergatives lack the internal argument altogether and thus denote predicates of eventualities to begin with: $\lambda e.\textsc{work}(e)$. As a result, Existential Closure in \REF{ex:tatevosov:87f} does not happen and Event Identification applies to this denotation directly. This creates a relation between working events and result states of the applications having been processed:

\ea   \sib{[ \textit{rabota}- ] [CS [ \textit{sto anket ob-} ]]} $= \lambda s.\lambda e. \textsc{work}(e) \wedge \textsc{processed}(\textsc{one hundred applications})(s) \wedge \cnst{cause}(s)(e)$\label{ex:tatevosov:89}
\z

\noindent The derivation in \REF{ex:tatevosov:87}, apart from lexical information, functional application, and event identification independently argued for in the literature, makes use of two more specific mechanisms for deriving the denotation of a prefixed VP: Causative Shift and Existential Closure. These are exactly the parts of the derivation where existing theories of resultative predication vary along the two dimensions in \REF{ex:tatevosov:83}--\REF{ex:tatevosov:84}.

\newpage
The first dimension in \REF{ex:tatevosov:83} has to do with specific assumptions about the way the eventive and stative components of a complex description are combined. In \REF{ex:tatevosov:87}, this is accomplished by Causative Shift. Establishing a relation between events and result states is an unavoidable ingredient of any theory of resultative predication. However, the relation can in be conceived of as an output of a rule of construal, as in \REF{ex:tatevosov:83a}, or as the denotation of a phonologically null morpheme, as in \REF{ex:tatevosov:83b}.

Some theories assume a special rule of construal. For English resultatives, for example, \citet{Rappaport.Hovav.Levin2001} propose Template Augmentation, while \citet{Rothstein2004} makes use of the \cnst{tpconnect} relation. \citeposst{Beck2005} Principle R falls under this type of approach, too. \citet{Kratzer2005} considers both options for the derivation of Germanic resultatives, and having discussed a number of pros and contras, opts for the morphological solution. She poses a special derivational morpheme represented in the syntax which is interpreted as introducing the relation of immediate causation. 

For Slavic languages, both types of semantic theories of lexical prefixation have been proposed. Some, in line with \citet{Kratzer2005}, assume a syntactically represented (and typically silent) grammatical element. For instance, \citet[53]{BabkoMalaya1999} suggests a morpheme that adjoins to a prefix, takes its denotation as an argument and yields a relation to which the verb stem denotation applies at the next stage of derivation. \citet[ch. 3]{Biskup2019} introduces a similar element as the denotation of a functional head located in between the verb stem and the projection of a prefix. On the other hand, \citet{Gehrke2008} argues that Rothstein’s architecture is better suited than Kratzer’s to account for attested constraints on the derivation of resultative configurations cross-linguistically; see \citet[61--63]{Gehrke2008} for specifics of her critical assessment of \citet{Kratzer2005} and \citet{Rothstein2004}.  

It is not immediately obvious, however, whether the choice between a special rule of construal or a silent functional element necessarily makes different predictions with respect to the event structure or argument structure of nAS lexically prefixed verbs. While \REF{ex:tatevosov:87d} assumes the latter, I believe that nothing in what follows should be incompatible with the alternative.

The other parameter of theoretical variation in \REF{ex:tatevosov:84a}--\REF{ex:tatevosov:84b} is how a theory treats an argument that appears as an object in an unprefixed configuration and is perceived as a subcategorized argument of a verb stem. For the transitive verb stem \textit{pisa}-, for example, this argument is an object effected in the course of an event (‘write a letter’ and the like).

The sketch in \REF{ex:tatevosov:87a}--\REF{ex:tatevosov:87h} relies on the three assumptions in \REF{ex:tatevosov:90}.

\eanoraggedright  \textsc{Argument projection in nAS configurations}\label{ex:tatevosov:90}
\eanoraggedright Lexically, transitive verbs denote relations between events and (internal)     arguments.\label{ex:tatevosov:90a}
\ex The argument of a prefix is syntactically realized.\label{ex:tatevosov:90b}
\ex The argument of a verb stem gets existentially bound.\label{ex:tatevosov:90c}
\z
\z

\noindent The same or a similar set of assumptions can be found, for example in \citet{Rothstein2004} for English resultatives and in \citet[53]{BabkoMalaya1999}  for Russian resultative lexical prefixes. A related proposal appears in \citet{Bruening2021}, where a verb can merge with a head that prevents it from projecting its internal argument and existentially binds it. (The other such head makes the implicit object definite.)

An alternative advanced by \citet{Kratzer2005} for Germanic and \citet{Biskup2019} for Slavic languages does not assume \REF{ex:tatevosov:90a}, thus eliminating \REF{ex:tatevosov:90c}. Specifically, \citet[48-50]{Biskup2019} suggests that prefixless transitive verbs denote predicates of eventualities, while their internal arguments are introduced by the functional structure within the complement of V. \citeauthor{Biskup2019}’s proposal amounts to treating the event /  argument structure of prefixed and unprefixed configurations as being products of different derivations, unlike in \REF{ex:tatevosov:87a}--\REF{ex:tatevosov:87h}, where the former is derived from the latter.

Assessing these two types of theories, one can observe that neither is exempt from certain vulnerabilities. Theories assuming that internal arguments are not true arguments of a verb stem need additional stipulations to explain selectional restrictions that stems massively impose on internal arguments, but not on external arguments (as well as other asymmetries between external and internal arguments discussed, e.g., in \citealt{Kratzer2003}). Theories that assume existential closure of the internal argument (or a similar mechanism) bear the burden of finding an independent motivation for this operation. For example, \citet{Tatevosov2015a} argues that at least for some nAS lexically prefixed verbs existential closure of the internal argument results in coercion, which prevents the derivation from outputting a predicate with an empty extension. However, it is not clear whether such considerations extend to all prefixal configurations.

Properly motivating something like the existential closure may also be substantial for preventing the system from generating ungrammatical structures in which both the internal argument of a stem  and the argument of a prefix are projected. For \textit{pod-pisa-t’} ‘sign\textsuperscript{PFV}’, as a reviewer points out, that would yield something like \REF{ex:tatevosov:91} with the meaning ‘one’s name has been written, and the agreement is consequently in a signed state’.

\ea[*]
{\gll Volodja  pod-pisa-l  imja     dogovor.  \\
  Volodja  \textsc{lp}-write-\textsc{pst.m}  name.\textsc{acc}  agreement.\textsc{acc}\\
\glt Intended: ‘Volodja signed the name an agreement.’\label{ex:tatevosov:91}}
\z

\noindent One strategy of ruling \REF{ex:tatevosov:91} out may be to explore possible ways in which it  could have come about and find a reason why such derivations are problematic. To produce \REF{ex:tatevosov:91} out of the available ingredients one may need, assuming the syntax in \REF{ex:tatevosov:92}, a derivation along the lines of \REF{ex:tatevosov:93}.

\ea{}   [\textsubscript{VP}  \textit{imja}  [\textsubscript{V'} \textit{pisa-} [  CS \textit{pismo pod}- ] ]  ]\label{ex:tatevosov:92}
\z

\ea\label{ex:tatevosov:93}
\ea \sib{\textit{pisa}-} $= \lambda x.\lambda e. \textsc{write}(x)(e)$           \hfill($=$ \REF{ex:tatevosov:87b})\label{ex:tatevosov:93a}
\ex    \sib{CS [ \textit{pis’mo pod-} ]} $= \lambda s.\lambda e. \textsc{signed}(\textsc{letter})(s) \wedge \cnst{cause}(s)(e)$  \hfill (= \REF{ex:tatevosov:87e})\label{ex:tatevosov:93b}
\ex    \sib{\textit{pisa}- [CS [ \textit{pis’mo pod-} ]]} \hfill (Generalized Event                       Identification)\label{ex:tatevosov:93c}\\$= \lambda x. \lambda s.\lambda e.\textsc{write}(x)(e) \wedge \textsc{signed}(\textsc{letter})(s) \wedge \cnst{cause}(s)(e)$
\ex    \sib{\textit{imja} \textit{pisa}- [CS [ \textit{pis’mo pod-} ]]} $= \lambda s.\lambda e.\textsc{write}(\textsc{name})(e) \wedge \textsc{signed}(\textsc{letter})(s) \wedge \cnst{cause}(s)(e)$ \hfill (FA)\label{ex:tatevosov:93d}
\z
\z

\noindent Unlike in \REF{ex:tatevosov:87}, the internal argument position does not undergo existential closure in \REF{ex:tatevosov:93}, but stays open. Then, to combine \REF{ex:tatevosov:93a} with \REF{ex:tatevosov:93b} one would need a rule of construal that can be thought of as a generalized version of event identification. \citeposst{Kratzer1996} rule only identifies events that fall under a particular description (of type $\langle v, t \rangle$) with events that stand in a particular relation to individuals (of type $\langle e, \langle v, t \rangle\rangle$). Generalized event identification (GEI) would identify events that are elements of any relations to any other entities:

\ea   \textsc{Generalized event identification} (GEI)\label{ex:tatevosov:94}

If $\alpha$ is a branching node and \{$\beta$, $\gamma$\} is the set of its daughters,

\sib{$\beta$} is in $D_{\langle\sigma_{\alpha_{1}}, \langle..., \langle \sigma_{\alpha_{n}}, \langle v,t\rangle\rangle\rangle\rangle}$, and \sib{$\gamma$} is in $D_{\langle\sigma_{\beta_{1}}, \langle ..., \langle \sigma_{\beta_{k}}, \langle v,t\rangle\rangle\rangle\rangle}$, \\
then \sib{$\alpha$} $= \lambda \omega_{\alpha_{1}}.\, ...\, . \lambda \omega_{\alpha_{n}}. \lambda \omega_{\beta_{1}}.\, ...\, . \lambda \omega_{\beta_{k}}.\lambda e$. \sib{$\alpha$}$(\omega_{\alpha_{1}})(...)(\omega_{\alpha_{n}})(e) =$ \\\sib{$\beta$}$(\omega_{\beta_{1}})(...)(\omega_{\beta_{k}})(e) = 1$, where $\sigma_{\alpha_{i}}, \sigma_{\beta_{j}}$ and the like are types, and $\omega_{\alpha_{i}}, \omega_{\beta_{j}}$ are corresponding variables.  
\z

\noindent Now, with GEI in \REF{ex:tatevosov:94}, in \REF{ex:tatevosov:93c}--\REF{ex:tatevosov:93d} we have identified events of writing the name with events that bring about a state of an agreement having been signed.

Given that \REF{ex:tatevosov:91} is drastically ungrammatical, something must be wrong with this scenario, and the question is what exactly. One can come up with a number of possible answers. GEI may not be a rule of construal that is available in natural languages in the first place. This is not implausible, given that cross-linguistically, we do not find descriptions, either morphosyntactically simplex or derived, of complex eventualities \textit{e} such that \textit{e} is, for example, an eventuality of Olga eating a sausage and Cathy swimming in the pool, or the like. Another reason may be that while there is nothing wrong with GEI in itself, there are output constraints that descriptions like \REF{ex:tatevosov:93d} fail to meet. For instance, one can speculate that larger syntactic configurations in which VPs like \REF{ex:tatevosov:92} occur fall short of the sources of structural case for both nominals. Also, it is not unimaginable that \REF{ex:tatevosov:93d} does not meet some or other semantic well-formedness requirement based, e.g., on fine-grainedness considerations.

As one can see, at this juncture we may take several ways to go. However, discussing this in further detail should be a topic for a separate study, which I am not attempting in what follows. 

One more potentially sensitive issue with the existential closure in \REF{ex:tatevosov:87f} may have to do with verbs of creation, as a reviewer points out. ‘Write’ is such a verb, and its internal argument, the one that undergoes $\exists_x$ in \REF{ex:tatevosov:87g}, is an effected object. Effected objects only exist at the end of an event of creation. For this reason, the reviewer indicates, existentially bound objects are only warranted in perfective sentences, which entail the culmination, but not in imperfective sentences that come with no such entailment. However, there is one aspect of \REF{ex:tatevosov:87} that may make it less problematic than initially appears. Consider first a predicate like \REF{ex:tatevosov:ifn33}.

\ea\label{ex:tatevosov:ifn33}
$\lambda e.\exists x.\textsc{write}(x)(e) \wedge \textsc{letter}(x)$ 
\z

\noindent In \REF{ex:tatevosov:ifn33}, an object that falls under \textsc{letter} only exists at the end of writing, and \REF{ex:tatevosov:ifn33} falls short of capturing that. \REF{ex:tatevosov:87g}, repeated as \REF{ex:tatevosov:iifn33}, is crucially different in that an effected individual falls under no particular description; it is only identified as an entity staying in the write relation to an event.

\ea\label{ex:tatevosov:iifn33}
$\lambda e.\exists x.\textsc{write}(x)(e)$
\z

\noindent \textsc{Write}, as any verb of creation, is incremental (e.g., \citealt{Pinon2008}). This means, due to mapping to subobjects, that in any writing subevent there is, trivially, something written. This written entity may not qualify as an entity that can be described by \textsc{letter} or, in fact, by any other (nominal) predicate. But \REF{ex:tatevosov:iifn33}, unlike \REF{ex:tatevosov:ifn33}, does not require that. For there to be something written it is enough that some writing has occurred. I believe that the combination of \textsc{write} being incremental and the lack of any descriptive content associated with the effected object circumvents the problem that would have emerged otherwise and does in fact emerge with any predicate like \REF{ex:tatevosov:ifn33}.

\subsubsection{AS prefixes}\label{sec:3.3.3}
\hypertarget{RefHeadingToc164647630}{}
Having in mind the spectrum of possibilities for nAS lexical prefixes discussed above, one can address the derivation of AS prefixal configurations. One of those are “pure perfectivizing”, or “empty” prefixes like \textit{na}- in \textit{na-pisa-t’} ‘write\textsuperscript{PFV}’, or \textit{pro}- in \textit{pro-čita-t’} ‘read\textsuperscript{PFV}’, whereby the prefixation does not induce any perceivable change of the lexical meaning. With such prefixes, an entity undergoing writing or reading is identical to one attaining a state of having been written or read. In \sectref{tat:sec:2}, for \textit{na-pisa-t’} \textit{pis’mo} ‘write a letter’ the representation repeated in \REF{ex:tatevosov:95} is assumed.

\ea   \sib{[\textsubscript{VP} \textit{na-pisa- pis’mo}]} $= \lambda s.\lambda e. \textsc{write}(\textsc{letter})(e) \wedge \cnst{cause}(s)(e) \wedge  \textsc{written}(\textsc{letter})(s)$\label{ex:tatevosov:95}
\z

\noindent In \REF{ex:tatevosov:95}, the fact that the letter is understood as the object of writing and the holder of a state of having been written is fully explicit: the object DP figures in \REF{ex:tatevosov:95} as an argument of both the verb and the prefix \textit{na}-. To achieve this result, one may need to make certain assumptions about a corresponding syntactic structure and syntax-semantic interface. If, for example, a DP is allowed to merge in more than one thematic position, and if such positions are interpreted as variables bound by the topmost copy of the DP, \REF{ex:tatevosov:95} is not difficult to derive. As an example of a theory of this type one can think of \citeauthor{Ramchand2008}'s (\citeyear{Ramchand2008} and elsewhere) First Phase Syntax and its later developments (for a recent discussion see \citealt{Svenonius2016}).

As before, \REF{ex:tatevosov:95} is not the only admissible option. In the literature, one can find two logically possible alternatives to \REF{ex:tatevosov:95}. First, the object is construed as being an argument of an AS verb without being an argument of a prefix. Second, the other way around, the object is an argument of a prefix, but not of a verb.

The first approach is what \citet[55-56]{BabkoMalaya1999} proposes for pure perfectivizing prefixes, which she keeps apart from the rest of the lexical prefixes. Rather than introducing a result state, empty prefixes take the denotation of a verb stem as an argument and modify it by contributing the entailment of “completion”. However, this comes at a cost: VPs containing empty prefixes are predicted to manifest less subevental complexity than other VPs projected by lexically prefixed verbs. As discussed in \sectref{tat:sec:2}, this is not the actual pattern attested in Russian and other Slavic languages. It is not clear if the analysis can be amended in order to incorporate observations from \sectref{tat:sec:2} while maintaining that the object is an argument of the stem only. 

Under the second approach, the object of AS-prefixed verbs can be treated on a par with the object of nAS-prefixed verbs, that is, as an argument of a prefix only. This is the path \citet{Biskup2019} takes, assuming no derivational asymmetry between AS and nAS configurations. \citeposst{Romanova2007} approach falls under this type, too. This implies that the internal argument of the verb stem is treated identically in both AS and nAS configurations. As was mentioned above, for \citet{Biskup2019}, this argument is absent altogether. 

If, on the other hand, one assumes a derivational scenario for nAS-prefixes along the lines of \REF{ex:tatevosov:87}, where the argument of the stem gets existentially bound, the derivation of a parallel AS-configuration would output the relation in \REF{ex:tatevosov:96} instead of \REF{ex:tatevosov:95}.

\ea   \sib{[\textsubscript{VP} \textit{na-pisa- pis’mo}]} $= \lambda s.\lambda e.\exists x [\textsc{write}(x)(e) \wedge \cnst{cause}(s)(e) \wedge \textsc{written}(\textsc{letter})(s)]$\label{ex:tatevosov:96}
\z

\noindent At first sight, \REF{ex:tatevosov:96} looks very wrong: the relation is predicted to have events in it where something other than a letter undergoes writing. One can speculate, however, that in our world as well as in other worlds where causal regularities resemble ours, the extensions of \REF{ex:tatevosov:95} and \REF{ex:tatevosov:96} may come out indistinguishable. The predicate of states $\lambda s.\textsc{written}(\textsc{letter})(s)$ comes, according to \REF{tat:ex:64}, with very specific descriptive properties. States that fall under this predicate %in the lexical context of \textit{pisa-} ‘write’
are states of having been written, not other kinds of states that may result from writing events. World knowledge, part of which is knowledge of causal dependencies, tells us that for any $x$ to attain such a state% of being written
, $x$ has to undergo writing, otherwise an event and a state would fail to fall under \cnst{cause}. This restricts the events that can enter the relation in \REF{ex:tatevosov:96} to those where the theme of writing is the same as the holder of the state, hence the relation in \REF{ex:tatevosov:96} would reduce to the relation in \REF{ex:tatevosov:95}.

\subsubsection{Restrictive and non-restrictive AS prefixes}\label{sec:3.3.4}
\hypertarget{RefHeadingToc164647631}{}
AS lexically prefixed verbs do not form an entirely homogeneous class. As the above discussion suggests, this class contains all the verbs with “empty” prefixes. It does not only contain such verbs, however. One can identify many prefixes that are AS without being empty. 

Consider \REF{ex:tatevosov:97}--\REF{ex:tatevosov:98}.

\ea\label{ex:tatevosov:97}
\gll Volodja  pro-čita-l   doklad.\\
  Volodja  \textsc{lp}-read-\textsc{pst.m}  report.\textsc{acc}\\
\glt ‘Volodja read the report.’
\z

\ea\label{ex:tatevosov:98}
\gll Volodja   za-čita-l   doklad.\\
  Volodja  \textsc{lp}-read-\textsc{pst.m}  report.\textsc{acc}\\
\glt ‘Volodja read the report out loud.’
\z

\noindent Both \REF{ex:tatevosov:97} and \REF{ex:tatevosov:98} pass the diagnostic for argument-sharedness in \REF{ex:tatevosov:78}--\REF{ex:tatevosov:79}, since both entail, modulo the aspectual difference, the prefixless sentence in \REF{ex:tatevosov:99} with the same object.

\ea\label{ex:tatevosov:99}
\gll Volodja   čita-l   doklad.\\
  Volodja  read-\textsc{pst.m}  report.\textsc{acc}\\
\glt ‘Volodja was reading the report.’
\z

\noindent While \textit{pro-čita-t’} ‘read\textsuperscript{PFV}’ in \REF{ex:tatevosov:97} is a paradigmatic example of “pure perfectivization”, \textit{za-čita-t’}  in \REF{ex:tatevosov:98} has a more specific meaning: it comes with the inference that a reading of the report is aloud and, possibly, in front of an audience.

What \REF{ex:tatevosov:97}--\REF{ex:tatevosov:98} and lots of similar examples tell us can be captured in the following way. Every event of reading a report can take the report into a state associated with the prefix \textit{pro}-, a state of having been read. No matter what specific properties a reading event has, whether reading is silent or aloud, slow or fast, thoughtful or superficial, it can culminate in a state that falls under $\lambda s.\textsc{read}(\textsc{report})(s)$. But only a subset of reading events can culminate in a state of having been read out aloud in front of the audience, associated with \textit{za}-. Events of silent reading do not enter a causal chain where the maximal element is such a state. For \textit{za-čita-t’} ‘read out aloud’ and  \textit{pro-čita-t’} ‘read\textsuperscript{PFV}’, therefore, the following holds:

\ea   For any $x$,\label{ex:tatevosov:100}
\ea
$\{e: \exists s [\cnst{cause}(s)(e) \wedge \textsc{read}(x)(s) ] \} = \{e: \textsc{read}(x)(e) \}$

\glt    ‘Events that cause a state of having been read are reading events.’\label{ex:tatevosov:100a}
\ex     $\{e: \exists s [ \cnst{cause}(s)(e) \wedge \textsc{read out aloud}(x)(s)]\} \subset \{e:  \textsc{read}(x)(e)\}$

\glt    ‘Events that cause a state of having been read out aloud are a proper subset   of reading events.’\label{ex:tatevosov:100b}
\z
\z

\noindent Prefixes like \textit{pro}- and \textit{za}- can thus be called \textsc{non-restrictive} and \textsc{restrictive}, respectively. According to \REF{ex:tatevosov:100a}, the set of events that bring about a \textit{pro}-state of having been read is identical to the entire set of culminating reading events that fall under the unprefixed verb. \textit{Pro-} is thus non-restrictive on the extension of \textit{čita}- ‘read’. \textit{Za}- is restrictive in the sense that events that cause a \textit{za}-state of having been read out aloud form a restricted proper subset of reading events, \REF{ex:tatevosov:100b}. For \REF{ex:tatevosov:100a}--\REF{ex:tatevosov:100b}, this is represented diagrammatically in \figref{fig:tatevosov:101}.

\begin{figure}
\begin{tikzpicture}
\tikzstyle{TextNode} = [text width=5cm, anchor=north west]
% Upper descriptive row
\node[TextNode] at (0,0) (desc-1) {$\{e: \exists s [\cnst{cause}(s)(e) \wedge{}$\newline $\textsc{read}(\textsc{report})(s)]\}$\\ $= \{e: \textsc{read}(\textsc{report})(e)\}$};
\node[TextNode] at (6,0) (desc-3) {$\{s: \textsc{read}(\textsc{report})(s)\}$\\‘states of the report having been read’};

\node (p1) [circle,fill=black,inner sep=2pt] at (1,-2.5) {};
\node (p2) [circle,fill=black,inner sep=2pt] at (2,-3.5) {};
\node (p2fit) [draw, fit={(p2)}, rectangle, rounded corners=4pt, inner xsep=1cm, inner ysep=.25cm] {};
\node (p1p2fit) [draw, fit={(p1) (p2fit)}, rectangle, rounded corners=4pt, inner sep=.25cm] {};

\node (p3) [circle,fill=black,inner sep=2pt] at (8,-2.5) {};
\node (p4) [circle,fill=black,inner sep=2pt] at (8,-3.5) {};

\node (p4fit) [draw, fit={(p4)}, rectangle, rotate=45, inner sep=.25cm] {};
\node (p3p4fit) [draw, fit={(p3) (p4fit)}, rectangle, rotate=-45, minimum size=1.75cm] {};

% Lower descriptive row
\node[TextNode] at (0,-5) (desc-2) {$\{e: \exists s [\cnst{cause}(s)(e) \wedge{}$\newline $\textsc{read.aloud}(\textsc{report})(s)]\}$\\
‘events that cause the report to attain a state of having been read aloud in front of the audience’};
\node[TextNode] at (6,-5) (desc-4) {$\{s: \textsc{read.aloud}(\textsc{report})(s)\}$\\
‘states of the report having been read aloud in front of the audience'};

% Paths
\draw[dashed,-{Triangle[]}] (p1p2fit) -> (p3p4fit);
\draw[dashed,-{Triangle[]}] (p2fit) -> (p4fit);
\foreach \i in {1,2,3,4}{\draw (desc-\i) -- (p\i);}
\end{tikzpicture}
    \caption{Process and result state components of \textit{pro-čita-t'} `read\textsuperscript{PFV}' and \textit{za-čita-t'} `read\textsuperscript{PFV} out aloud'}
    \label{fig:tatevosov:101}
\end{figure}

It is for this reason that the restrictive AS prefixes produce intuitions that the lexical meaning of a non-derived verb has been affected by prefixation. The non-derived predicate has the entire set of reading eventualities in its extension. But when it comes to bringing about a state of having been read aloud in front of the audience, associated with the prefix \textit{za}- in \REF{ex:tatevosov:98}, only a proper subset of them is involved, a different eventuality description. Crucially, restrictive AS-prefixation is always a strengthening of the input lexical meaning, never a weakening, and the proper subset relation in \REF{ex:tatevosov:100b} captures that in a principled way. This case is thus crucially different from examples like \REF{tat:ex:69} with \textit{o-pisa-t'} ‘describe’ in \sectref{tat:sec:3.2}. Whereas both \textit{za-čita-t’} ‘read out aloud’ and \textit{o-pisa-t’} ‘describe’ only make use of a proper part of events from the original extension of the unprefixed verb, \textit{o-pisa-t’} ‘describe’, but not \textit{za-čita-t’} ‘read out loud’ acquires additional events that have not been there before prefixation; see also footnote \ref{fn:30}.

These considerations allow us to give the notion of “empty prefixation” the following content:

\eanoraggedright   \textsc{Empty prefixes as non-restrictive AS prefixes}\label{ex:tatevosov:102}
\eanoraggedright Empty prefixes denote result states compatible with every eventuality in the extension of a non-derived verb.\label{ex:tatevosov:102a}
\ex A lexical prefix P is empty, or non-restrictive with respect to the verb stem V iff $\forall x\forall e [$\sib{V}$(x)(e) \rightarrow \Diamond \exists s[$\sib{P}$(x)(s) \wedge \cnst{cause}(s)(e)]]$\label{ex:tatevosov:102b}
\z
\z

\noindent \REF{ex:tatevosov:102b} generalizes what is said in \REF{ex:tatevosov:100a} about the non-restrictive prefix \textit{pro}- in \textit{pro-čita} ‘read\textsuperscript{PFV}’. According to \REF{ex:tatevosov:102b}, P is empty with respect to V just in case the following holds: if V’s denotation is true of an event and an individual, this event can take this individual into a state described by P. This is not so for restrictive prefixes: for them, there are events that stay in the V-relation to an individual but cannot culminate in a P-state that would hold of the same individual. As we have just seen, not every reading event can bring about a state of having been read aloud in front of the audience described by the prefix \textit{za}-.  The $\Diamond$ operator in \REF{ex:tatevosov:102b} takes care of the fact that the extension of a non-derived verb in a world can contain events that do not culminate in that world; see \citet{Martin.Demirdache2020} for a recent overview of predicates of non-culminating eventualities.

Note that \REF{ex:tatevosov:102b} only defines empty prefixation for a class of verbs that are relations between individuals and events, i.e. verbs that take an internal argument. This seems to make a welcome prediction: empty prefixation is characteristic of transitive and unaccusative, but not unergative verbs. The reader is referred to \citet[ch. 4]{Romanova2007} for further discussion.

According to \REF{ex:tatevosov:102a}--\REF{ex:tatevosov:102b}, empty prefixes are not empty after all. Just like other lexical prefixes, they introduce a result state into the event structure of a prefixed configuration. “Emptiness”, that is, the lack of a perceivable lexical meaning change resulting from prefixation, is an epiphenomenon of there being no narrowing down effect of the type represented in \REF{ex:tatevosov:100b} and \figref{fig:tatevosov:101}. According to \REF{ex:tatevosov:102b}, for any event $e$ and individual $x$, once $e$ and $x$ are in the relation established by a non-derived verb, x can attain a state brought about by $e$ that falls under the denotation of an empty prefix. As we have just seen, this is what happens in case of \textit{pro-čita-t’} ‘read\textsuperscript{PFV}’ as opposed to \textit{za}-\textit{čita-t’} ‘read out aloud\textsuperscript{PFV}’.

This view of “emptiness” bears certain resemblance to the idea circulating in the literature since \citet{Vey1952} and \citet{Schooneveld1958}, see, e.g., \citet[][193--194]{Filip1999}. As \citet[177-178]{Gehrke2008} puts it, “it could be argued that ‘empty’ prefixes still express particular meanings in isolation, but that they contribute a redundant meaning when they combine with the verb, since they supply information that is already presupposed in the simple verb”.\footnote{Effects of this type may not be confined to Slavic prefixation. I am grateful to an anonymous reviewer for turning my attention to the parallelism between “empty” prefixation discussed here and the role of the clitic \textit{se} attached to verbs of consumption in Spanish (\citealt{Campanini.Schafer2011}). In \textit{comerse una manzana} `eat.\textsc{refl} an apple', the reviewer indicates, “\textit{se} expresses, according to \citet{Campanini.Schafer2011}, a relation of ingestion of the theme by the external argument, but this relation is already expressed by \textit{comer una manzana}” `eat an apple'.} 

Not infrequently, this redundancy is understood in terms of a literal overlap of the lexical meaning of a prefix and a verb. For instance, \citet{Janda.etal2013} suggest that the prefix \textit{u-} indicates that the figure is moving away from the ground, and the same meaning component can be found in the description of \textit{krast’} ‘steal’: stealing involves movement of a stolen entity away from its possessor. For \citeauthor{Janda.etal2013}, \textit{u-krast’} ‘steal\textsuperscript{PFV}’ is therefore expected to be an instance of “empty prefixation”, since the former meaning is contained in the latter. 

With \REF{ex:tatevosov:102}, what “is already presupposed in a simple verb” can receive a slightly different characterization. It can be understood as the culmination potential of a simplex verb, the ability of every eventuality from its extension to bring an object into a certain result state. A non-restrictive prefix happens to denote exactly such a state, so its attachment has no perceivable effects on lexical semantics. It is in this sense that the non-restrictive AS prefixes supply presupposed information (even if, technically, they are not redundant, since they introduce a result state into the event structure that had not been there before).

Typically, simplex imperfective verbs have one perfective counterpart derived by “empty” prefixation. Nothing in \REF{ex:tatevosov:102} suggests, however, that this always has to be the case. In fact, one can easily find cases of AS lexical prefixation where several prefixes can be characterized as pure perfectivizing with respect to the same stem. In such a case speakers judge different prefixal counterparts of the non-derived verb synonymous. An example is \textit{žari-t’} ‘fry, roast’ and its prefixed partners in \REF{ex:tatevosov:103} \citep{Polivanova1975}.

\ea\label{ex:tatevosov:103}
\ea
\gll   Volodja   žari-l   kartošk-u.  \\
    Volodja   fry-\textsc{pst.m}  potato-\textsc{acc}\\
\glt  ‘Volodja was frying potatoes.’\label{ex:tatevosov:103a}
\ex
\gll   Volodja   \minsp{\{} po-žari-l /  pod-žari-l /  za-žari-l /  iz-žari-l\} kartošk-u.\\
Volodja  {} \textsc{lp}-fry-\textsc{pst.m} {}   \textsc{lp}-fry-\textsc{pst.m} {}   \textsc{lp}-fry-\textsc{pst.m} {}   \textsc{lp}-fry-\textsc{pst.m} potato-\textsc{acc}\\
\glt  ‘Volodja fried potatoes.’\label{ex:tatevosov:103b}
\z
\z

\noindent The prefixed verbs \textit{po-žari-t’}, \textit{pod-žari-t’}, \textit{za-žari-t’}, and \textit{iz-žari-t'} in \REF{ex:tatevosov:103b} are all associated with result states with more or less the same descriptive properties which can be brought about by any frying eventuality described by \textit{žari-t’} ‘fry, roast’ in \REF{ex:tatevosov:103a}. As a result, truth-conditional differences between the four prefixed variants are not detectable, so none of them can be a candidate for the status of the only “pure perfective” aspectual partner of the non-derived verb.

On the other hand, a non-derived verb can fail to have an aspectual partner derived by empty prefixation. One typical instance of such a pattern occurs if the extension of a verb is too wide and heterogeneous to satisfy \REF{ex:tatevosov:102b}. \REF{ex:tatevosov:104a}--\REF{ex:tatevosov:104b} illustrate the verb \textit{kopa-t’} ‘dig’.

\ea\label{ex:tatevosov:104}
\ea 
\gll    Volodja   kopa-l   jam-u. \\
    Volodja   dig-\textsc{pst-m}   hole-\textsc{acc} \\
\glt  ‘Volodja was digging a hole.’\label{ex:tatevosov:104a}
\ex
\gll    Volodja   kopa-l   ogorod. \\
    Volodja   dig-\textsc{pst.m} vegetable.garden.\textsc{acc}\\
\glt  ‘Volodja was digging up a vegetable garden.’\label{ex:tatevosov:104b}
\z
\z

\noindent As \REF{ex:tatevosov:104a}--\REF{ex:tatevosov:104b} suggest, the relation between individuals and events denoted by \textit{kopa-t’} ‘dig’ contains at least two types of ordered pairs: those where digging events are coupled with effected objects like \textit{jama} ‘a hole’ in \REF{ex:tatevosov:104a}, and those where events are related to affected objects like \textit{ogorod} ‘vegetable garden’ in \REF{ex:tatevosov:104b}. However, no prefix can denote result states attained by both types of entities in the course of digging. Rather, those states are distributed over different prefixes, as \REF{ex:tatevosov:105}--\REF{ex:tatevosov:106} show.

\ea\label{ex:tatevosov:105}
\ea
\gll   Volodja  vy-kopa-l  jam-u.  \\
    Volodja  \textsc{lp}-dig-\textsc{pst.m}  hole-\textsc{acc}  \\
\glt  ‘Volodja dug a hole.’\label{ex:tatevosov:105a}
\ex[*]
{\gll   Volodja  vy-kopa-l  ogorod.  \\
    Volodja  \textsc{lp}-dig-\textsc{pst.m}  vegetable.garden.\textsc{acc}\\
\glt  Intended: ‘Volodja dug up a vegetable garden.’\label{ex:tatevosov:105b}}
\z
\z 

\ea\label{ex:tatevosov:106}
\ea[*]
{\gll   Volodja  vs-kopa-l  jam-u.  \\
    Volodja  \textsc{lp}-dig-\textsc{pst.m}  hole-\textsc{acc}  \\
\glt  Intended: ‘Volodja dug a hole.’\label{ex:tatevosov:106a}}
\ex
\gll   Volodja  vs-kopa-l  ogorod.  \\
    Volodja  \textsc{lp}-dig-\textsc{pst.m}  vegetable.garden.\textsc{acc}\\
\glt  ‘Volodja dug up a vegetable garden.’\label{ex:tatevosov:106b}
\z
\z

\noindent In \REF{ex:tatevosov:105a}--\REF{ex:tatevosov:105b}, the prefix \textit{vy-} expresses a state of having been effected as the result of digging. In \REF{ex:tatevosov:106a}--\REF{ex:tatevosov:106b}, the prefix \textit{vs}- introduces a state of having been (totally) affected by digging. If one recognizes two distinct relations \textit{kopa-t’\textsubscript{EFF}} `dig with an effected object’ and \textit{kopa-t’\textsubscript{AFF}} `dig with an affected object’, both being subsets of the denotation of \textit{kopa-t’}, one can argue that \textit{vy}- is pure perfectivizing with respect to \textit{kopa-t’\textsubscript{EFF}},  while \textit{vs}- is so with respect to \textit{kopa-t’\textsubscript{AFF}}. Given the highly degraded status of \REF{ex:tatevosov:105b} and \REF{ex:tatevosov:106a}, neither prefix can be characterized as pure perfectivizing with respect to \textit{kopa-t’} as a whole, i.e., with respect to \textit{kopa-t’\textsubscript{EFF}} $\,\cup$ \textit{kopa-t’\textsubscript{AFF}}.

The above observation only holds, of course, if one believes that \textit{kopa-t’\textsubscript{EFF}}{ } and \textit{kopa-t’\textsubscript{AFF}} fall under the denotation of the same unprefixed verb, \textit{kopa-t’} ‘dig\textsuperscript{IPFV}’. If this is not the case, and \textit{kopa-t’\textsubscript{EFF}} and \textit{kopa-t’\textsubscript{AFF}} form the denotations of two distinct, but homophonous lexical items, each can be said to have its own pure perfectivizing partner. As a reviewer points out, this analysis may be warranted by “the verbs’ distinct meanings, s-selectional properties, and behavior under prefixation”. On the other hand, there is a reason to assume that both \textit{vy-kopa-t’} ‘dig\textsuperscript{PFV}’ and  \textit{vs-kopa-t’} in \REF{ex:tatevosov:105}--\REF{ex:tatevosov:106} ‘dig\textsuperscript{PFV}’ are derivatives of the same unprefixed verb: distinct result states are brought about by an activity with exactly the same descriptive properties: ‘to break up and move soil using a tool, a machine, or your hands’, as Cambridge Dictionary puts it. If the lexical meaning of \textit{kopa-t’} only entails these  characteristics of the activity and is underspecified with respect to the type of object, the fact that we find both  \textit{kopa-t’\textsubscript{EFF}} and \textit{kopa-t’\textsubscript{AFF}} relations as part of its denotation comes as no surprise.

Another class of non-derived verbs that fail to have a single aspectual partner derived by empty prefixation are manner of motion verbs like \textit{lete-t’} ‘fly’ that do not come with any specific lexical entailments about the source and goal of motion. The AS lexically prefixed counterparts of \textit{lete-t’} ‘fly’, \textit{u-lete-t’} ‘fly away’ and \textit{pri-lete-t’}  ‘arrive flying’ introduce the state ‘be outside of the deictic center as the result of flying’ or ‘be at the deictic center as the result of flying’. These states are too specific to be compatible with any flying event denoted by \textit{lete-t’} ‘fly’, hence neither \textit{u-lete-t’} ‘fly away’ nor \textit{pri-lete-t’} ‘arrive flying’ come across as involving empty prefixation; both should be regarded as instances of restrictive prefixation. 

\subsubsection{Dyadic prefixes}\label{sec:3.3.5}
\hypertarget{RefHeadingToc164647632}{}
Dyadic prefixes are singled out by a prominent property of theirs: verbs with such prefixes have an extra argument which is not there in the corresponding unprefixed configuration. Consider \REF{ex:tatevosov:107}.

\ea\label{ex:tatevosov:107}
\gll Volodja  v-pisa-l  dopolnitel’n-oe  uslovi-e v  dogovor.\\
Volodja  \textsc{lp}-write-\textsc{pst.m}  additional-\textsc{acc}  condition-\textsc{acc} in  agreement.\textsc{acc}\\
\glt ‘Volodja wrote an additional condition into the agreement.’
\z

\noindent \REF{ex:tatevosov:107} entails a result state ‘an additional condition is in the agreement’, with the goal PP \textit{v dogovor} ‘in the agreement’ specifying the location of the condition after an eventuality culminates. A non-prefixed counterpart of \REF{ex:tatevosov:107} in \REF{ex:tatevosov:108} is highly degraded.

\ea[\textsuperscript{??}]
{\gll Volodja  pisa-l  dopolnitel’n-oe   uslovi-e  v  dogovor.  \\
     Volodja  write-\textsc{pst.m}  additional-\textsc{acc}   condition-\textsc{acc}  in  agreement.\textsc{acc}\\
\glt Intended: ‘Volodja was writing an additional condition into the agreement.’\label{ex:tatevosov:108}}
\z

\noindent Verbs like \textit{pisa-t’} ‘write’ under normal circumstances do not describe a movement in the physical space and for that reason do not take goal (and source) adjuncts, as in \REF{ex:tatevosov:108}. Given \REF{ex:tatevosov:108}, one can conclude that the goal PP in \REF{ex:tatevosov:107} can only appear as an argument of the prefix. In \REF{ex:tatevosov:107} the prefix \textit{v-}, therefore, introduces a state of the condition being in the agreement; such a state is attained as the result of writing. The prefix can thus be treated as dyadic: one specifying a relation between its external argument (‘an additional condition’; Figure)  and its internal argument (‘the agreement’; Ground). For various analyses of the syntactic structure projected by such prefixes see \citet{Ramchand2004,Svenonius2004,Svenonius2008,Romanova2007,Biskup2019}. Figure \ref{fig:3} illustrates the structure of \REF{ex:tatevosov:107} along the lines of \citet{Svenonius2004}.

\begin{figure}
\begin{forest}
[VP [V [R [\textit{v-} ] ] [V [\textit{pisa-} ] ] ] [RP [DP [\textsc{figure}\\\textit{uslovie}, roof ] ] [R' [\textit{t}\textsubscript{R} ] [PP [\textsc{ground}\\\textit{v dogovor}, roof ] ] ] ] ]
\end{forest}
\caption{Proposal according to \citet{Svenonius2004}}
\label{fig:3}
\end{figure}

Another example of a dyadic lexical prefix and its unprefixed counterpart is shown in \REF{ex:tatevosov:110}--\REF{ex:tatevosov:111}.

\ea\label{ex:tatevosov:110}
\gll Volodja  is-pisa-l  tetrad’  stix-ami.  \\
  Volodja  \textsc{lp}-write-\textsc{pst.m}  notebook.\textsc{acc}  poem-\textsc{instr.pl}\\
\glt ‘Volodja filled the notebook with poems.’
\z

\ea[*]
{\gll Volodja  pisa-l  tetrad’  stix-ami.  \\
   Volodja  write-\textsc{pst.m}  notebook.\textsc{acc}  poem-\textsc{instr.pl}\\
\glt Intended: ‘Volodja was writing the notebook with poems.’\label{ex:tatevosov:111}}
\z

\noindent In \REF{ex:tatevosov:110}, the direct object \textit{tetrad’} ‘notebook’ enters a result state of being filled with poems. As in \REF{ex:tatevosov:108}, the unprefixed counterpart of \REF{ex:tatevosov:110} shown in \REF{ex:tatevosov:111} is severely ungrammatical. However, \REF{ex:tatevosov:110} with \textit{is-pisa-t’} ‘fill sth. with sth. written’ differs from \REF{ex:tatevosov:107} with \textit{v-pisa-t’} ‘write into’ in one significant respect.

In \REF{ex:tatevosov:107}, the theme of writing, \textit{uslovie} ‘condition’, is realized as the direct object of the prefixed verb. In \REF{ex:tatevosov:110}, the theme of writing are poems, which however, appears as an instrumental DP \textit{stixami} ‘with poems’, not as the direct object. The direct object is \textit{tetrad’} ‘notebook’, which denotes the location of the poems that come about in the course of writing. The DP \textit{tetrad’} ‘notebook’ is thus Ground, the internal argument of the prefix, whereas \textit{stixami} is the Figure, its external argument.

\citet[110ff.]{Romanova2007} argues that this configuration should be analyzed in parallel with the verbal passive. The passive feature on the prefix does the regular job of the passive: descriptively, it induces “demotion” of the external argument, which makes the internal argument eligible for “promotion”. As a result, the final position of \textit{tetrad’} in \REF{ex:tatevosov:110} is the one where the external argument of prefixes is normally found, as, e.g., in \REF{ex:tatevosov:107}. The parallelism between configurations like \REF{ex:tatevosov:110} and verbal passives is reinforced by the fact that in both cases the demoted external argument is realized as an adjunct bearing the instrumental case.\footnote{I am grateful to Berit Gehrke (p.c.) who noted the parallelism between examples like \REF{ex:tatevosov:110} and the \textit{spray}/\textit{load}-alternation (as in \textrm{\textit{John loaded the hay onto the wagon}} \textrm{and} \textrm{\textit{John loaded the wagon with the hay}}). Remarkably, \citeposst{Larson1990} analysis of the alternation involves the demotion of the locatum (\textit{the hay}) to the adjunct position and A-movement of the location (\textit{the wagon}) to the specifier position; see \citet{Beavers2017} for further discussion. Essentially, this is what \citet{Romanova2007} proposes for examples like \REF{ex:tatevosov:110}.}

This completes the survey of the typology of prefixes set up in Figure \ref{fig:1}. Before closing this chapter, however, a few reflections on the descriptive properties of results states in light of the generalization in \REF{tat:ex:64} above are due.

\subsection{Descriptive properties of result states}\label{sec:3.4}
\hypertarget{RefHeadingToc164647633}{}
In this section I will address a few related open questions that all have to do with how the descriptive properties of a state predicate in the denotation of a lexical prefix are computed. The toy fragment in section \sectref{tat:sec:2.2} has been abstracted away from this issue.

To begin with, one should recognize two conflicting types of empirical evidence concerning the meaning of prefixes. On the one hand, descriptive properties of states they introduce must be specific enough for a theory to be able to capture semantic judgments about the differences between combinations of the same verb stem with different prefixes. In \textit{na-pisa-t’} ‘write\textsuperscript{PFV}’ writing events lead the object to a state of having been written. In \textit{pod-pisa-t’}, the object enters a state of having been signed. In both cases, writing events fall under \textit{pisa-t’} ‘write’. The source of the difference between the two verbs must therefore be descriptive properties of states associated with the prefixes \textit{na-} and \textit{pod-}. This implies that the meaning of prefixes is rather fine-grained and specific, as is pointed out in \REF{tat:ex:64} above. 

On the other hand, the descriptive properties of a state associated with a lexical prefix cannot be too specific, since every prefix should be able to combine with an open class of verb stems. A predicate of states of having been written cannot be lexically associated with the prefix \textit{na-}. Otherwise if the same prefix combines with \textit{risova-t’} ‘paint’ to produce \textit{na-risova-t’} ‘paint\textsuperscript{PFV}’, one would get a set of painting events that bring about a state of having been written, which obviously makes no sense. 

At the moment, I know of no fully articulated theory that can successfully reconcile the challenges of having a rather specific state description for any combination of a prefix and a verb while keeping the meaning of a prefix general enough for it not to be subject to severe lexical restrictions. However, the space of possibilities shaping the content of such a theory seems to be more or less clear. With a certain simplification, it can be reduced to two clear-cut strategies (and various mixed combinations of their ingredients), which can be called \textsc{functional} strategy and \textsc{intrinsic meaning} strategy.

The functional strategy would assume that in order to build up a state description in a specific lexical context a prefix needs access to the information about the event description denoted by the verb stem. A straightforward way of implementing this would be to treat (the projection of) a prefix as a (possibly very partial) function that takes the denotation of a verb as an argument and picks out an appropriate predicate of states. 

Above, in \REF{ex:tatevosov:87}, the (shifted) denotation of the projection of a prefix and that of the verb stem, repeated as \REF{ex:tatevosov:112a}--\REF{ex:tatevosov:112b}, are combined by Event Identification. The output is a relation between events and states in \REF{ex:tatevosov:112c}.

\largerpage
\ea\label{ex:tatevosov:112}
\eanoraggedright      The denotation of the projection of a prefix (after Causative Shift)\label{ex:tatevosov:112a}

$\lambda s.\lambda e. \textsc{signed}(\textsc{letter})(s) \wedge \cnst{cause}(s)(e)$

\ex      The denotation of a verb stem (after existential closure of the internal   argument)\label{ex:tatevosov:112b}

$\lambda e. \exists x [\textsc{write}(x)(e)]$

\ex      The denotation of the VP projected by a prefixed verb\label{ex:tatevosov:112c}

$\lambda s.\lambda e.\exists x [\textsc{write}(x)(e) \wedge \textsc{signed}(\textsc{letter})(s)\wedge \cnst{cause}(s)(e) ]$
\z
\z

\noindent The simplified fragment in \REF{ex:tatevosov:87}/\REF{ex:tatevosov:112} was abstracted away from how the descriptive properties of a state in \REF{ex:tatevosov:112a} were computed. But given the above considerations, the denotation in \REF{ex:tatevosov:112a} may need to be amended as in \REF{ex:tatevosov:113a}. It would apply to \REF{ex:tatevosov:112b} yielding \REF{ex:tatevosov:113b} as the denotation of \textit{pod-pisa-t’} ‘sign’ instead of \REF{ex:tatevosov:112c}.

\ea\label{ex:tatevosov:113}
\ea      The denotation of the projection of a prefix\label{ex:tatevosov:113a}

$\lambda P.\lambda s. \lambda e. \textsc{pod}(P)(\textsc{letter})(s) \wedge \cnst{cause}(s)(e) \wedge P(e)$

\ex      The denotation of the VP projected by a prefixed verb\label{ex:tatevosov:113b}

$\lambda s.\lambda e.\exists x [ \textsc{pod}(\lambda e'.\exists x[\textsc{write}(x)(e')])(\textsc{letter})(s) \wedge \cnst{cause}(s)(e) \wedge \textsc{write}(x)(e) ]$
\z
\z

\noindent In \REF{ex:tatevosov:113a}, \textit{pod-} is of type $\langle\langle v,t \rangle, \langle v, \langle v,t\rangle\rangle\rangle$. It does not name a relation between events and states directly, as in \REF{ex:tatevosov:112a}. Rather, this is done by \textsc{pod}($\lambda e'.\exists x [\textsc{write}(x)(e')])$ in \REF{ex:tatevosov:113b}. Every state description created by a prefix is thus relativized to the event predicate denoted by  a verb stem it combines with. Combining with the stem \textit{pisa}- ‘write’, \textit{pod-} yields a relation between individuals and states of being signed:

\ea   $\textsc{pod}(\lambda e'.\exists x [\textsc{write}(x)(e')]) = \lambda x.\lambda s.\textsc{signed}(x)(s)$\label{ex:tatevosov:114}
\z

\noindent The denotation of the prefix itself can be represented as the function in \REF{ex:tatevosov:115}.

\ea   \sib{\textit{pod-}} $= \lambda x.\lambda P.\lambda s.\lambda e. \textsc{pod}(P)(x)(s) \wedge \cnst{cause}(s)(e) \wedge P(e)$\label{ex:tatevosov:115}
\z

\largerpage[2]
\noindent Technically, this should work. However, this type of theory reveals an obvious vulnerability. The meaning of a prefix is now conceived of as a mapping from event predicates to relations between individuals and states, as in \REF{ex:tatevosov:113a}. For the \textit{pod}-verbs in \REF{ex:tatevosov:116}, this mapping would look like \REF{ex:tatevosov:117}.

\ea\label{ex:tatevosov:116}
\ea    \textit{pod-pisa-} ‘sign’   $\leftarrow$   \textit{pod} $+$ \textit{pisa-} ‘write’
\ex     \textit{pod-dela-} ‘fake’   $\leftarrow$   \textit{pod} $+$ \textit{dela-} ‘do, make’
\ex     \textit{pod-stroi-} ‘set up’   $\leftarrow$   \textit{pod} $+$ \textit{stroi-} ‘build’
\ex      \textit{pod-kova-} ‘shoe (a horse)’   $\leftarrow$   \textit{pod} $+$ \textit{kova-} ‘forge’
\ex ...
\z
\z

\ea   \textsc{The extension of \textit{pod}-, quasi-formally}\label{ex:tatevosov:117}

\{ \textit{pisa-} ‘write’ $\mapsto $ $\lambda x.\lambda s.\textsc{signed}(x)(s)$, \textit{dela-} ‘do, make’ $\mapsto $ $\lambda x. \lambda s.\textsc{faked}(x)(s)$, \textit{stroi-} ‘build’ $\mapsto$ $\lambda x.\lambda s.\textsc{set up}(x)(s)$, \textit{kova-} ‘forge’ $\mapsto$ $\lambda x.\lambda s.\textsc{shoed}(x)(s)$, ... \}
\z

\noindent If nothing else is said, each combination of a prefix and a stem is unique, and the extension of every prefix starts looking like a list of such combinations. Strictly speaking, this gives us as many submeanings of a prefix as there are stems it can combine with. This may make the acquisition task for a learner of Russian and other Slavic languages rather costly. For every prefix, one needs to learn a list like \REF{ex:tatevosov:117}: a collection of state descriptions which it maps verbs from its domain to. Whether one can convince oneself that this is how the system works may ultimately depend on experimental evidence from L1 and L2 acquisition and, as Berit Gehrke (p.c.) points out, on evidence from “languages with less transparent morphology, in which there are different lexical items for each of the meanings”. At the moment, however, this seems to be more a prospect for a future inquiry than a working research agenda.

It should be noted, finally, that making the denotation of a verb stem an argument of a prefix, as in \REF{ex:tatevosov:115}, makes it inevitable that the semantics incorporates the effects of causative shift. This is so because the verb denotation is not only used to calculate a state description, but also, as before, contributes eventualities that cause such a state. As a result, the semantics of a prefix now duplicates the event structure of the entire complex event description: all its subevental components are already present in \REF{ex:tatevosov:115}. Whether this is a welcome consequence is a separate question that I am not trying to address here.

An alternative to the view just outlined is to assume that lexical prefixes have  intrinsic meanings in and by themselves and to specify how they interact with the meaning of a verb stem. This interaction should not be conceived of as verb-specific and must follow from independent principles and constraints.

A prominent line of inquiry within this \textsc{intrinsic meaning} strategy is to explicate the meaning of a prefix in terms of properties of states defined by a spatial configuration between the figure and the ground. Of those two, one receives syntactic realization, while the other can (but does not have to) be implicit. Typically, given that the vast majority of prefixes have prepositional counterparts, the meaning of a prefix is thought of in terms of a relation established by the corresponding preposition, e.g., \textit{na} ‘on, at the surface of’, \textit{pod} ‘under’,  \textit{u} ‘at’, and so on. 

For instance, the prefix \textit{pod}- can be regarded as introducing a state of the figure being under the ground. In \textit{pod-pisa-t’} ‘sign’ \citep[67]{Biskup2019}, the figure is a signature, which has no syntactic realization. The ground is an entity under which a signature appears, and it projects as the direct object. \citeauthor{Biskup2019}’s semantics for \textit{pod-pisa-t' pis'mo} 'sign the letter', adjusted to the format used throughout this chapter, is shown in \REF{ex:tatevosov:118}.\footnote{\REF{ex:tatevosov:118} differs from \citeposst{Biskup2019} original proposal in two respects: Existentially bound implicit arguments are represented in \citeauthor{Biskup2019} as free variables; verb stems come with no individual arguments.}

\ea   $\lambda s.\lambda e.\exists x [\textsc{under}(\textsc{letter})(x)(s) \wedge \cnst{cause}(s)(e) \wedge \textsc{write}(x)(e)]$\label{ex:tatevosov:118}
\z

\noindent One can take a step further and argue that certain figure/ground configurations may not be confined to the spatial domain and extend to other domains as well, most importantly, to the temporal domain. An early exemplar of this approach can be found in \citet{Flier1975}, who points out that “verbal prefixes [...] ultimately make reference to abstract delimitation, dimension, and direction, eliciting met\-a\-phor\-i\-cal interpretations of these notions in nonspatial universes” \citep[219]{Flier1975}. In \citet{Flier1975}, one can find an analysis of the four “spanned” Russian prefixes \textit{pere}-, \textit{pro}-, \textit{ob}-, and \textit{po}- along these lines. With a certain idealization, \citeposst{Kagan15} degree semantics for a number of Russian prefixes can be regarded as a further development of this idea, one that takes degree structures to be an analytical tool that can capture semantic regularities across domains. 

Another possible implementation of the intrinsic meaning strategy would be to set up a polysemy model of the semantics of lexical prefixes, whereby every prefix represents a family of meanings structured in a particular way. A prominent example of this line of thought has been developed in a series of publications by Laura Janda and colleagues of hers. In \citet{Janda.etal2012}, for example, every prefix is associated with a “radial category”, where submeanings are organized into a network, centered around the prototype, as in \figref{fig:tatevosov:119}, which represents the semantics of the prefix \textit{u-}. 

\begin{figure}
\begin{forest}
for tree={font=\scshape}
[keep/save
[place/fit,name=place
[move\\away,font=\scshape\bfseries,name=away
[move\\downwards,name=down
[reduce,name=reduce]
]
]
]
]
\begin{scope}[nodes={align=center,font=\scshape}]
\draw (reduce.west) -- ++(-2cm,0pt) node [left] (control) {control} -- (down.west);
\path let \p1=(control), \p2=(place) in node at (\x1,\y2) (cover) {cover\\completely};
\draw (place) -- (cover.east) -- (away.west);
\path let \p1=(control), \p2=(away) in node at (\x1,\y2) (depart) {depart\\from norm};
\draw (away.west) -- (depart);
\draw (reduce.east) -- ++(2cm,0pt) node [right] (harm) {harm} -- (down.east);
\path let \p1=(harm), \p2=(away) in node at (\x1,\y2) (perceive) {perceive};
\draw (perceive) -- (away);
\end{scope}
\end{forest}
\caption{Radial category for the prefix \textit{u}- \citep{Janda.etal2012}}
\label{fig:tatevosov:119}
\end{figure}

In \figref{fig:tatevosov:119}, the prototype is \textsc{move} \textsc{away}, which manifests itself in the meaning of verbs like \textit{u-beža-t’} ‘run away’. Other submeanings that form the network in \figref{fig:tatevosov:119} can be defined in terms of properties they share with the parent node(s). For instance, the link between moving away and moving downwards is established “because when an object moves away, it sinks below the horizon” \citep[249]{Janda.etal2012}.

The intrinsic meaning strategy does not run into the acquisition problem inherent to the previous one. If a theory is equipped with an elaborated machinery of making spatial configurations, movements and paths entirely explicit, it should ultimately be able to develop a model in which the general meaning of a prefix interacts with the meaning of a stem in a regular and predictable way, deriving the range of event structural and argument structural characteristics of prefixed verbs discussed above. %Note as well that since the prefix does not take the verb stem as its argument, the causative component does not have to be made part of its meaning unlike in \REF{ex:tatevosov:115}.

At the moment, however, building up such a model looks more like a theoretical desideratum than like an articulated theory. One obvious difficulty in realizing this desideratum is that if a prefix is assigned a single general meaning, it can be too devoid of specific content for a theory to derive testable predictions about the descriptive properties of a result state in every lexical configuration. For instance, assume ‘under’ as the meaning of the prefix \textit{pod}-. Looking at \textit{pod-pisa-t’} ‘sign $=$ under-write’, \textit{pod-rubi-t}’ ‘under-cut’ or \textit{pod-žari-t’} ‘under-fry’ one can try to convince oneself that all of these verbs involve something happening at the bottom of  or under the entity denoted by the internal argument. However, this is much less obvious in case of \textit{pod-stroi-t’} ‘set up, frame, trick $=$ under-build’ or \textit{pod-stegnu-t’} ‘hurry up, spur $=$ under-lash’. If a car accident has been set up, can one say that the car accident ends up in the under-relation to some other entity? Failure of a theory to address questions like this makes its empirical prospects dubious. 

Polysemy models typically come with a more fine-grained semantic structure that can potentially be more successful in assigning the right descriptive properties to result states. However, they face additional challenges. One is to identify and motivate the set of individual submeanings of a prefix and relations between them. For example, for \figref{fig:tatevosov:119} to be fully convincing, one has to present elaborated arguments showing that the prefix \textit{u-} has exactly 10 submeanings, that these meanings are exactly \{\textsc{move} \textsc{away,} \textsc{move} \textsc{downwards,} \textsc{control,} \textsc{reduce,} \textsc{harm,} ...\}, as in \figref{fig:tatevosov:119}, and that they are connected to each other as indicated. If structures like \figref{fig:tatevosov:119} are taken for granted rather than thoroughly argued for on empirical grounds, predictions of a theory would be difficult if not impossible to assess. 

\section{Summary}\label{tat:sec:4} %4. /
\hypertarget{RefHeadingToc164647634}{}
After an overview of the landscape of the derivational morphology of Russian, taken to represent several typical characteristics of the entire Slavic family, this chapter addressed two main topics. One is the event structure of simplex imperfective, lexically prefixed perfective, and secondary imperfective verbs and their extended projections. The other is lexical prefixation, with special attention to the parameters of variation across lexical prefixes.

In \sectref{tat:sec:2}, relying on the idea circulating in the literature since \citet{Klein1995}, I presented a few pieces of data that argue for a higher subevental complexity of prefixed configurations, be they perfective or secondary imperfective, compared to simplex configurations. This is manifested in a number of phenomena that are typically taken to diagnose event-structural complexity: scope of decomposition adverbs, interpretation under negation, argument realization patterns. In event-semantic terms, therefore, VPs projected by prefixed perfective verbs are proposed to be analyzed as denoting relations between events and result states. Secondary imperfectives can be argued to share this event-structural make up with prefixed perfectives. Their prefixless imperfective counterparts, in contrast, lack a result state.  

The results of \sectref{tat:sec:2}, as well as numerous considerations found in the literature, strongly suggest that the presence of a result state in the event structure of VP can be attributed to lexical prefixes. However, lexical prefixes do not form a homogeneous class and fall into a number of subclasses, identified and motivated in \sectref{tat:sec:3}. The basic divide within the entire set of lexically prefixed verbs is that between the verbs where the denotation of a stem undergoes a non-compositional shift after prefixation, and the verbs where this does not happen. The latter, compositional lexically prefixed verbs, are further classified according to the argument structure of a prefix. I proposed to identify monadic prefixes that take a single individual argument, the holder of a result state, and dyadic prefixes which denote a relation between two entities, the figure and the ground. Among the verbs with monadic prefixes, two important classes are isolated. One contains verbs with argument-sharing prefixes, that is, prefixes whose individual argument is identical to the internal argument of the verb stem. The other contains verbs where the two entities are distinct. Finally, argument sharing prefixes can be characterized as either restrictive or non-restrictive. The former introduce result states only compatible with a restricted subset of eventualities from the original extension of a verb stem. The latter denote states that can be brought about by every such an eventuality; they are also known under the name “pure perfectivizing” or “empty” prefixes.  

\section*{Abbreviations}

\begin{tabularx}{.5\textwidth}{@{}lQ}
\textsc{2}&second person\\
\textsc{3}&third person\\
\textsc{acc}&accusative\\
\textsc{dat}&dative\\
\textsc{f}&feminine\\
\textsc{gen}&genitive\\
\textsc{indef}&indefinite\\
\textsc{inf}&infinitive\\
\textsc{instr}&instrumental\\
\textsc{ipfv}&imperfective\\
\textsc{lp}&lexical prefix\\
\textsc{m}&masculine\\
\textsc{n}&neuter\\
\end{tabularx}%
\begin{tabularx}{.5\textwidth}{lQ@{}}
\textsc{n/t}&past passive participle morphology\\
\textsc{npst}&non-past\\
\textsc{pfv}&perfective\\
\textsc{pl}&plural\\
\textsc{pst}&past\\
\textsc{refl}&reflexive\\
\textsc{sg}&singular\\
\textsc{slp}&superlexical prefix\\
\textsc{smfct}&semelfactive\\
\textsc{th}&thematic element\\
\textsc{yva}&the \textit{-yva-} morpheme, ``secondary imperfective''\\
%\textsc{}&\\
%&\\ % this dummy row achieves correct vertical alignment of both tables
\end{tabularx}

\section*{Acknowledgments}

I would like to express my immense gratitude to Berit Gehrke and anonymous reviewers for their benevolent attention to this work and helpful comments. This chapter owes a lot to the conversations with Daniel Altshuler, Tanya Bondarenko, Sabine Iatridou, Hans-Robert Mehlig, Elena Viktorovna Paducheva, Rou\-my\-ana Pancheva, Anna Pazelskaya, David Pesetsky, Mitya Privoznov, Olav Muel\-ler-Reichau about syntax and semantics of aspect and many other beautiful things. 

\sloppy\printbibliography[heading=subbibliography,notkeyword=this]

\end{document}
